%%%%%%%%%%%%%%%%%%%%%%%%%%%%%%%%%%%%%%%%%%%%%%%%%%%%%%%%%%%%%%%%%%%%%%
% Lớp 11 | Giới hạn. Hàm số liên tục | Giới hạn của dãy số
% Dạng 1D3?1-2: Phương pháp đặt thừa số chung (lim hữu hạn)
% Số câu: 268  (sinh tự động bởi tach_cau.py; từ nay thầy sửa trực tiếp file này)
%%%%%%%%%%%%%%%%%%%%%%%%%%%%%%%%%%%%%%%%%%%%%%%%%%%%%%%%%%%%%%%%%%%%%%

%%%%%%%%%%%%%%%% TRẮC NGHIỆM 4 PHƯƠNG ÁN %%%%%%%%%%%%%%%%

% nguồn: Ok_Toan11_TracNghiem_Lan2
\begin{ex}%[1D3N1-2]%[ma:70b2dfd0]
	 $\lim\dfrac{1-n^2}{2n^2+1}$ bằng
	 \choice
		 {$0$}
		{$\dfrac{1}{2}$} 
		 {$\dfrac{1}{3}$}
		 {\True $-\dfrac{1}{2}$}
	 \loigiai{
		Ta có $\lim\dfrac{1-n^2}{2n^2+1}=\lim\dfrac{\dfrac{1}{n^2}-1}{2+\dfrac{1}{n^2}}=\dfrac{0-1}{2+0}=-\dfrac{1}{2}$. 
	} 
\end{ex}

% nguồn: Ok_Toan11_TracNghiem_Lan2
\begin{ex}%[1D3N1-2]%[ma:fa000cc8]
	Tính $\lim\dfrac{3}{2n+3}$ được kết  quả bằng
	\choice 
		 {$\dfrac{3}{2}$}
		 {\True $0$}
		{$1$} 
		 {$+\infty$}
	\loigiai{ 
		Ta có $\lim\dfrac{3}{2n+3}=\lim\dfrac{\dfrac{3}{n}}{2+\dfrac{3}{n}}=\dfrac0{2+0}=0$.\\ 
	} 
\end{ex}

% nguồn: Ok_Toan11_TracNghiem_Lan2
\begin{ex}%[1D3N1-2]%[ma:cbf47b8c]
	 Tính giới hạn $I=\lim\dfrac{2n+2\,017}{3n+2\,018}$.
	 \choice
		{\True  $I=\dfrac{2}{3}$}
		{$I=\dfrac{3}{2}$} 
		{$I=\dfrac{2\,017}{2\,018}$} 
		{$I=1$} 
	\loigiai{ 
		Ta  có $I=\lim\dfrac{2n+2\,017}{3n+2\,018}=\lim\dfrac{2+\dfrac{2\,017}{n}}{3+\dfrac{2\,018}{n}}=\dfrac{2+0}{3+0}=\dfrac{2}{3}$.
	} 
\end{ex}

% nguồn: Ok_Toan11_TracNghiem_Lan2
\begin{ex}%[1D3N1-2]%[ma:95564896]
	Cho $I=\lim\dfrac{\sqrt{4n^2+5}+n}{4n-\sqrt{n^2+1}}$.  Khi đó giá trị của $I$ là
	 \choice
		 {\True $I=1$}
		 {$I=\dfrac{5}{3}$}
		 {$I=-1$}
		{$I=\dfrac{3}{4}$} 
	\loigiai{ 
		Ta có $I=\lim\dfrac{\sqrt{4n^2+5}+n}{4n-\sqrt{n^2+1}}=\lim\dfrac{\sqrt{4+\dfrac{5}{n^2}}+1}{4-\sqrt{1+\dfrac{1}{n^2}}}=\dfrac{\sqrt{4+0}+1}{4-\sqrt{1+0}}=1$. 
	} 
\end{ex}

% nguồn: Ok_Toan11_TracNghiem_Lan2
\begin{ex}%[1D3N1-2]%[ma:f8439a0b]
	Giá trị $ A=\lim\dfrac{n^2+n}{12n^2+1}$  bằng
	\choice 
		{\True $\dfrac{1}{12}$} 
		 {$0$}
		 {$\dfrac{1}{6}$}
		 {$\dfrac{1}{24}$}
	\loigiai{ 
		Ta có $  A=\lim\dfrac{n^2+n}{12n^2+1}=\lim\dfrac{1+\dfrac{1}{n}}{12+\dfrac{1}{n^2}}=\dfrac{1+0}{12+0}=\dfrac{1}{12}$.\\
	Vậy  $ A=\dfrac{1}{12}$.}
\end{ex}

% nguồn: Ok_Toan11_TracNghiem_Lan2
\begin{ex}%[1D3N1-2]%[ma:695a4235]
	Giới hạn $\lim\dfrac{1-2n}{3n+1}$ bằng  bao nhiêu?
	\choice 
		{$\dfrac{2}{3}$} 
		 {$\dfrac{1}{3}$}
		{$1$} 
		 {\True $-\dfrac{2}{3}$}
	\loigiai{ 
		 Ta có $\lim\dfrac{1-2n}{3n+1}=\lim\dfrac{\dfrac{1}{n}-2}{3+\dfrac{1}{n}}=\dfrac{0-2}{3+0}=-\dfrac{2}{3}$.
	} 
\end{ex}

% nguồn: Ok_Toan11_TracNghiem_Lan2
\begin{ex}%[1D3N1-2]%[ma:12eedef4]
	$\lim\dfrac{3n^2+1}{n^2-2}$ bằng 
	 \choice
		{\True  $3$}
		 {$0$}
		 {$\dfrac{1}{2}$}
		{$-\dfrac{1}{2}$} 
	 \loigiai{
		Ta có $\lim\dfrac{3n^2+1}{n^2-2}=\lim\dfrac{3+\dfrac{1}{n^2}}{1-\dfrac{2}{n^2}}=\dfrac{3+0}{1-0}=3$. 
	 }
\end{ex}

% nguồn: Ok_Toan11_TracNghiem_Lan2
\begin{ex}%[1D3N1-2]%[ma:4e9ad7d7]
	Tính $\lim\dfrac{-3n+4}{3+n}$ được kết  quả bằng
	\choice 
		 {$-1$}
		 {$-\infty$}
		{\True $-3$} 
		{$4$} 
	 \loigiai{
		Ta  có $\lim\dfrac{-3n+4}{3+n}=\lim\dfrac{-3+\dfrac{4}{n}}{\dfrac{3}{n}+1}=\dfrac{-3+0}{0+1}=-3$.\\
	 }
\end{ex}

% nguồn: Ok_Toan11_TracNghiem_Lan2
\begin{ex}%[1D3N1-2]%[ma:43958057]
	Cho dãy  số $(u_n)$ với $u_n=\dfrac{1+2+3+\ldots+n}{n^2+1}$. Mệnh đề nào sau đây đúng?
	\choice 
		 {$\lim u_n=0$}
		{\True  $\lim u_n=\dfrac{1}{2}$}
		{Dãy số $(u_n)$ không có giới hạn khi $n\to+\infty$} 
		{$\lim  u_n=1$}
	\loigiai{ 
		Ta có  $\lim u_n=\lim\dfrac{1+2+3+\ldots+n}{n^2+1}=\lim\dfrac{n\left(n+1\right)}{2\left(n^2+1\right)}=\dfrac{1}{2}$.
	 }
\end{ex}

% nguồn: Ok_Toan11_TracNghiem_Lan2
\begin{ex}%[1D3N1-2]%[ma:4d84b724]
	$\lim\dfrac{n^3+4n-5}{3n^3+n^2+7}$  bằng
	\choice 
		 {$1$}
		{\True $\dfrac{1}{3}$} 
		 {$\dfrac{1}{4}$}
		{$\dfrac{1}{2}$} 
	\loigiai{ 
	Ta  có $\lim\dfrac{n^3+4n-5}{3n^3+n^2+7}=\lim\dfrac{1+\dfrac{4}{n^2}-\dfrac{5}{n^3}}{3+\dfrac{1}{n}+\dfrac{7}{n^3}}=\dfrac{1+0-0}{3+0+0}=\dfrac{1}{3}$.}
\end{ex}

% nguồn: Ok_Toan11_TracNghiem_Lan2
\begin{ex}%[1D3N1-2]%[ma:33686ce2]
	Tính $\lim\dfrac{5n+3}{2n+1}$. 
	\choice 
		{$1$} 
		{$+\infty  $}
		 {$2$}
		{\True  $\dfrac{5}{2}$}
	 \loigiai{
	Ta có $\lim\dfrac{5n+3}{2n+1}=\lim\dfrac{5+\dfrac{3}{n}}{2+\dfrac{1}{n}}=\dfrac{5+0}{2+0}=\dfrac{5}{2}$.} 
\end{ex}

% nguồn: Ok_Toan11_TracNghiem_Lan2
\begin{ex}%[1D3N1-2]%[ma:5eaa2a97]
	Tính $\lim\dfrac{8n^2+3n-1}{4+5n+2n^2}$. 
	\choice 
		{$2$} 
		{$-\dfrac{1}{2}$} 
		 {\True $4$}
		 {$-\dfrac{1}{4}$}
	 \loigiai{
	 Ta có $\lim\dfrac{8n^2+3n-1}{4+5n+2n^2}=\lim\dfrac{8+\dfrac{3}{n}-\dfrac{1}{n^2}}{\dfrac{4}{n^2}+\dfrac{5}{n}+2}=\dfrac{8+0-0}{0+0+2}=4$.}
\end{ex}

% nguồn: Ok_Toan11_TracNghiem_Lan2
\begin{ex}%[1D3N1-2]%[ma:e03b7d8f]
	 Tính $L=\lim\dfrac{n^3+n^2+1}{2\,018-3n^3}$.
	\choice 
		 {$\dfrac{1}{2\,018}$}
		 {$-3$}
		 {$+\infty$}
		 {\True $-\dfrac{1}{3}$}
	 \loigiai{
		Ta có $L=\lim\dfrac{n^3+n^2+1}{2\,018-3n^3}=\lim\dfrac{1+\dfrac{1}{n}+\dfrac{1}{n^3}}{\dfrac{2\,018}{n^3}-3}=\dfrac{1+0+0}{0-3}=-\dfrac{1}{3}$. 
	} 
\end{ex}

% nguồn: Ok_Toan11_TracNghiem_Lan2
\begin{ex}%[1D3N1-2]%[ma:41f31fb9]
	 $\lim\dfrac{2n+1}{n+1}$ bằng
	\choice 
		{$1$} 
		{\True  $2$}
		{$-2$} 
		 {$+\infty $}
	\loigiai{ 
	Ta  có $\lim\dfrac{2n+1}{n+1}=\lim\dfrac{2+\dfrac{1}{n}}{1+\dfrac{1}{n}}=\dfrac{2+0}{1+0}=2$.}
\end{ex}

% nguồn: Ok_Toan11_TracNghiem_Lan2
\begin{ex}%[1D3N1-2]%[ma:08884192]
	Tính giới hạn $I=\lim\dfrac{10n+3}{3n-15}$ ta được kết quả 
	 \choice
		{$I=-\dfrac{10}{3}$} 
		 {\True $I=\dfrac{10}{3}$}
		{$I=\dfrac{3}{10}$} 
		{$I=-\dfrac{2}{5}$} 
	\loigiai{ 
	 Ta có $I=\lim\dfrac{10n+3}{3n-15}=\lim\dfrac{10+\dfrac{3}{n}}{3-\dfrac{15}{n}}=\dfrac{10+0}{3-0}=\dfrac{10}{3}$.}
\end{ex}

% nguồn: Ok_Toan11_TracNghiem_Lan2
\begin{ex}%[1D3N1-2]%[ma:ac9c00d6]
	Tìm giới hạn  $I=\lim\dfrac{3n-2}{n+3}$.
	 \choice
		 {$I=-\dfrac{2}{3}$}
		{$I=1$} 
		 {\True $I=3$}
		{$k\in\mathbb{Z}$} 
	\loigiai{ 
		 Ta có $I=\lim\dfrac{3n-2}{n+3}=\lim\dfrac{3-\dfrac{2}{n}}{1+\dfrac{3}{n}}=\dfrac{3-0}{1+0}=3$.
	} 
\end{ex}

% nguồn: Ok_Toan11_TracNghiem_Lan2
\begin{ex}%[1D3N1-2]%[ma:9bce9395]
	$\lim\left(\sqrt{n^2-3n+1}-n\right)$  bằng
	 \choice
		{$-3$} 
		 {$+\infty$}
		{$0$} 
		{\True  $-\dfrac{3}{2}$}
	\loigiai{ 
		 Ta có
		\begin{align*}
			\lim\left(\sqrt{n^2-3n+1}-n\right) 
			&=\lim\dfrac{\left(\sqrt{n^2-3n+1}-n\right)\left(\sqrt{n^2-3n+1}+n\right)}{\sqrt{n^2-3n+1}+n}\\ 
			&=\lim\dfrac{-3n+1}{\sqrt{n^2-3n+1}+n}\\ 
			&=\lim\dfrac{-3+\dfrac{1}{n}}{\sqrt{1-\dfrac{3}{n}+\dfrac{1}{n^2}}+1}\\ 
			 &=-\dfrac{3}{2}
		\end{align*}
		Vậy $\lim\left(\sqrt{n^2-3n+1}-n\right)=-\dfrac{3}{2}$. 
	} 
\end{ex}

% nguồn: Ok_Toan11_TracNghiem_Lan2
\begin{ex}%[1D3N1-2]%[ma:ccba13d7]
	Kết quả của $\lim\dfrac{n-2}{3n+1}$ bằng 
	\choice 
		{\True  $\dfrac{1}{3}$}
		 {$-\dfrac{1}{3}$}
		 {$-2$}
		 {$1$}
	 \loigiai{
	Ta  có $\lim\dfrac{n-2}{3n+1}=\lim\dfrac{n\left(1-\dfrac{2}{n}\right)}{n\left(3+\dfrac{1}{n}\right)}=\lim\dfrac{1-\dfrac{2}{n}}{3+\dfrac{1}{n}}=\dfrac{1-0}{3+0}=\dfrac{1}{3}$.}
\end{ex}

% nguồn: Ok_Toan11_TracNghiem_Lan2
\begin{ex}%[1D3N1-2]%[ma:6e8e2b65]
	Tính  giới hạn: $\lim\dfrac{3n-\sqrt{4n^2+n+1}}{n+\sqrt{n^2-2n-2}}$.
	 \choice
		 {$-2$}
		{$-\dfrac{3}{2}$} 
		 {$+\infty$}
		{\True $\dfrac{1}{2}$} 
	 \loigiai{
		Ta  có
		\begin{align*}
			\lim\dfrac{3n-\sqrt{4n^2+n+1}}{n+\sqrt{n^2-2n-2}} 
			&=\lim\dfrac{3n-n\sqrt{4+\dfrac{1}{n}+\dfrac{1}{n^2}}}{n+n\sqrt{1-\dfrac{2}{n}-\dfrac{2}{n^2}}}\\ 
			&=\lim\dfrac{3-\sqrt{4+\dfrac{1}{n}+\dfrac{1}{n^2}}}{1+\sqrt{1-\dfrac{2}{n}-\dfrac{2}{n^2}}}\\ 
			 &=\dfrac{3-2}{1+1}\\
			&=\dfrac{1}{2}. 
		\end{align*}
	} 
\end{ex}

% nguồn: Ok_Toan11_TracNghiem_Lan2
\begin{ex}%[1D3N1-2]%[ma:17f89af4]
	 Giới hạn $\lim\dfrac{n-n^3}{n^2+n+1}$ bằng
	\choice 
		{$1$} 
		{\True $-\infty$} 
		{$-1$} 
		{$+\infty$} 
	 \loigiai{
		Ta có $\lim\dfrac{n-n^3}{n^2+n+1}=\lim\left(n \cdot  \dfrac{\dfrac{1}{n^2}-1}{1+\dfrac{1}{n}+\dfrac{1}{n^2}}\right)=-\infty$. \\
		 Vì $\lim n=+\infty$ và $\lim\dfrac{\dfrac{1}{n^2}-1}{1+\dfrac{1}{n}+\dfrac{1}{n^2}}=\dfrac{0-1}{1+0+0}=-1< 0$.\\
	 }
\end{ex}

% nguồn: Ok_Toan11_TracNghiem_Lan2
\begin{ex}%[1D3N1-2]%[ma:3f2a3d33]
	 Dãy số $\left(u_n \right)$ với $u_n=\dfrac{1+2+3+\cdots+n}{1\,011n^2+1\,012}$. Khi đó $\lim\left(u_n+1\right)$ bằng
	 \choice
		{$\dfrac{2\,019}{2\,022}$} 
		{\True  $\dfrac{2\,023}{2\,022}$}
		{$\dfrac{2\,022}{2\,023}$} 
		 {$\dfrac{2\,021}{2\,022}$}
	 \loigiai{
		Ta có $1+2+3+\cdots+n=\dfrac{n\left(n+1\right)}{2}$.\\ 
		 Nên
		\begin{align*}
			 u_n&=\dfrac{1+2+3+\cdots+n}{1\,011n^2+1\,012}\\
			 &=\dfrac{\dfrac{n\left(n+1\right)}{2}}{1\,011n^2+1\,012}\\
			 &=\dfrac{n^2+n}{2\,022n^2+2\,024}.
		\end{align*}
		 Do đó $u_n+1=\dfrac{2\,023n^2+n+2\,024}{2\,022n^2+2\,024}$.\\
		Suy  ra
		\begin{align*}
			 \lim\left(u_n+1\right)
			&=\lim\dfrac{2\,023n^2+n+2\,024}{2\,022n^2+2\,024}\\ 
			&=\lim\dfrac{2\,023+\dfrac{1}{n}+\dfrac{2\,024}{n^2}}{2\,022+\dfrac{2\,024}{n^2}}\\ 
			&=\dfrac{2\,023}{2\,022}. 
		\end{align*}
	 }
\end{ex}

% nguồn: Ok_Toan11_TracNghiem_Lan2
\begin{ex}%[1D3N1-2]%[ma:0b39e321]
	Giá trị của $B=\lim\dfrac{4n^2+3n+1}{\left(3n-1\right)^2}$  bằng
	 \choice
		 {\True $\dfrac{4}{9}$}
		{$\dfrac{4}{3}$} 
		 {$0$}
		 {$4$}
	 \loigiai{
		Ta có $B=\lim\dfrac{4n^2+3n+1}{\left(3n-1\right)^2}=\lim\dfrac{n^2\left(4+\dfrac{3}{n}+\dfrac{1}{n^2}  \right)}{n^2\left(3-\dfrac{1}{n} \right)^2}=\lim\dfrac{4+\dfrac{3}{n}+\dfrac{1}{n^2}}{\left(3-\dfrac{1}{n} \right)^2}=\dfrac{4+0+0}{\left(3-0\right)^2}=\dfrac{4}{9}$.
	} 
\end{ex}

% nguồn: Ok_Toan11_TracNghiem_Lan2
\begin{ex}%[1D3N1-2]%[ma:42655ce4]
	Giới hạn  $\lim\dfrac{2n^2+n+5}{2n+1}$ bằng
	\choice 
		{$1$} 
		{$0$} 
		 {$-\infty$}
		 {\True $+\infty$}
	\loigiai{ 
		Ta có $\lim\dfrac{2n^2+n+5}{2n+1}=\lim\left(n \cdot \dfrac{2+\dfrac{1}{n}+\dfrac{5}{n^2}}{2+\dfrac{1}{n}}\right)=+\infty$.  \\
		 Vì $\lim n=+\infty$ và $\lim\dfrac{2+\dfrac{1}{n}+\dfrac{5}{n^2}}{2+\dfrac{1}{n}}=\dfrac{2+0+0}{2+0}=1> 0$.
	 }
\end{ex}

% nguồn: Ok_Toan11_TracNghiem_Lan2
\begin{ex}%[1D3N1-2]%[ma:10ad5e3c]
	Giới hạn $\lim\dfrac{2n+1}{n^2+3n+2}$ bằng 
	\choice 
		 {$2$}
		{$+\infty$} 
		 {\True $0$}
		{$-\infty$} 
	 \loigiai{
		 Ta có $\lim\dfrac{2n+1}{n^2+3n+2}=\lim\dfrac{\dfrac{2}{n}+\dfrac{1}{n^2}}{1+\dfrac{3}{n}+\dfrac{2}{n^2}}=\dfrac{0+0}{1+0+0}=0$.\\
	 }
\end{ex}

% nguồn: Ok_Toan11_TracNghiem_Lan2
\begin{ex}%[1D3N1-2]%[ma:343b4f9a]
	Giới hạn $\lim\dfrac{\sqrt{1+5+\ldots+(4n-3)}}{2n-1}$ bằng 
	 \choice
		{$1$} 
		{$+\infty$} 
		{\True $\dfrac{\sqrt{2}}{2}$} 
		 {$0$}
	\loigiai{Ta có $1+5+\ldots+(4n-3)=\dfrac{n(1+4n-3)}{2}=\dfrac{n(4n-2)}{2}=n(2n-1)$.\\ 
	 Suy ra $\lim\dfrac{\sqrt{1+5+\ldots+\left(4n-3\right)}}{2n-1}=\lim\dfrac{\sqrt{n(2n-1)}}{2n-1}=\lim\dfrac{\sqrt{2-\dfrac{1}{n}}}{2-\dfrac{1}{n}}=\dfrac{\sqrt{2}}{2}$.}
\end{ex}

% nguồn: Ok_Toan11_TracNghiem_Lan2
\begin{ex}%[1D3N1-2]%[ma:0832639c]
	 Tính giới hạn $\lim\dfrac{4n+2018}{2n+1}$.
	 \choice
		{$\dfrac{1}{2}$} 
		 {$4$}
		 {\True $2$}
		 {$2018$}
	\loigiai{ 
	Ta có  $\lim\dfrac{4n+2018}{2n+1}=\lim\dfrac{4+\dfrac{2018}{n}}{2+\dfrac{1}{n}}=\dfrac{4+0}{2+0}=2$.}
\end{ex}

% nguồn: Ok_Toan11_TracNghiem_Lan2
\begin{ex}%[1D3N1-2]%[ma:bd710a40]
	$\lim\dfrac{\sqrt{4n^2+1}-\sqrt{n+2}}{2n-3}$ bằng 
	\choice 
		 {$\dfrac{3}{2}$}
		{$2$} 
		{\True  $1$}
		 {$+\infty$}
	\loigiai{ 
		Ta có  $\lim\dfrac{\sqrt{4n^2+1}-\sqrt{n+2}}{2n-3}=\lim\dfrac{\sqrt{4+\dfrac{1}{n^2}}-\sqrt{\dfrac{1}{n}+\dfrac{2}{n^2}}}{2-\dfrac{3}{n}}=\dfrac{2-0}{2}=1$.
	 }
\end{ex}

% nguồn: Ok_Toan11_TracNghiem_Lan2
\begin{ex}%[1D3N1-2]%[ma:765f119e]
	Giới hạn của dãy số $(u_n)$ với  $u_n=\dfrac{n(n+1)}{2n^2+n+3}$ bằng
	 \choice
		 {\True $\dfrac{1}{2}$}
		{$1$} 
		{$3$} 
		 {$+\infty$}
	 \loigiai{
		 Ta có $\lim u_n=\lim\dfrac{n(n+1)}{2n^2+n+3}=\lim\dfrac{n^2\left(1+\dfrac{1}{n}\right)}{n^2\left(2+\dfrac{1}{n}+\dfrac{3}{n^2}\right)}=\lim\dfrac{1+\dfrac{1}{n}}{2+\dfrac{1}{n}+\dfrac{3}{n^2}}=\dfrac{1+0}{2+0+0}=\dfrac{1}{2}$.\\
	 }
\end{ex}

% nguồn: Ok_Toan11_TracNghiem_Lan2
\begin{ex}%[1D3N1-2]%[ma:d0634789]
	Tính giới hạn  của $\lim\dfrac{2n-1}{3-n}$.
	 \choice
		{$-2$} 
		{$\dfrac{2}{3}$} 
		{$1$} 
		 {\True $-\dfrac{1}{3}$}
	\loigiai{ 
	 Ta có $\lim\dfrac{2n-1}{3-n}=\lim\dfrac{2-\dfrac{1}{n}}{\dfrac{3}{n}-1}=\dfrac{2-0}{0-1}=-\dfrac{1}{3}$.}
\end{ex}

% nguồn: Ok_Toan11_TracNghiem_Lan2
\begin{ex}%[1D3N1-2]%[ma:f7ca2f62]
	Giá trị của $\lim\dfrac{2-n}{n+1}$ bằng 
	\choice 
		{$1$} 
		 {$2$}
		{\True  $-1$}
		 {$0$}
	 \loigiai{
		 Ta có $\lim\dfrac{2-n}{n+1}=\lim\dfrac{\dfrac{2}{n}-1}{1+\dfrac{1}{n}}=\dfrac{0-1}{1+0}=-1$.
	} 
\end{ex}

% nguồn: Ok_Toan11_TracNghiem_Lan2
\begin{ex}%[1D3N1-2]%[ma:1ca27dbd]
	Giới hạn  $L=\lim\dfrac{(2n-3)(3n+1)}{6n^2+n+1}$ bằng
	 \choice
		{\True  $1$}
		{$0$} 
		 {$-\infty$}
		{$+\infty$} 
	 \loigiai{
		Ta  có $L=\lim\dfrac{(2n-3)(3n+1)}{6n^2+n+1}=\lim\dfrac{n\left(2-\dfrac{3}{n}\right) \cdot n\left(3+\dfrac{1}{n}\right)}{n^2\left(6+\dfrac{1}{n}+\dfrac{1}{n^2}\right)}=\lim\dfrac{\left(2-\dfrac{3}{n}\right) \cdot\left(3+\dfrac{1}{n}\right)}{6+\dfrac{1}{n}+\dfrac{1}{n^2}}=\dfrac{2\cdot 3}{6}=1$.\\
	 }
\end{ex}

% nguồn: Ok_Toan11_TracNghiem_Lan2
\begin{ex}%[1D3N1-2]%[ma:dfacae7a]
	Dãy số nào sau đây có giới hạn khác $0$? 
	\choice 
		{$\dfrac{1}{n}$} 
		 {$\dfrac{1}{\sqrt n}$}
		 {\True $\dfrac{n+1}{n}$}
		{$\dfrac{\sin n}{\sqrt  n}$}
	 \loigiai{
	Ta có $\lim\dfrac{n+1}{n}=\lim1+\lim\dfrac{1}{n}=1$.} 
\end{ex}

% nguồn: Ok_Toan11_TracNghiem_Lan2
\begin{ex}%[1D3N1-2]%[ma:49a8775d]
	 Tính giới hạn $\lim\dfrac{n^2-3n^3}{2n^3+5n-2}$.
	 \choice
		 {$\dfrac{1}{5}$}
		 {$0$}
		{\True $-\dfrac{3}{2}$} 
		{$\dfrac{1}{2}$} 
	\loigiai{ 
	 Ta có $\lim\dfrac{n^2-3n^3}{2n^3+5n-2}=\lim\dfrac{n^3\left(\dfrac{1}{n}-3\right)}{n^3\left(2+\dfrac{5}{n^2}-\dfrac{2}{n^3}\right)}=\lim\dfrac{\dfrac{1}{n}-3}{2+\dfrac{5}{n^2}-\dfrac{2}{n^3}}=\dfrac{0-3}{2+0-0}=-\dfrac{3}{2}$.}
\end{ex}

% nguồn: Ok_Toan11_TracNghiem_Lan2
\begin{ex}%[1D3N1-2]%[ma:aad7f24d]
	Giới hạn $\lim\dfrac{8n^5-2n^3+1}{2n^2-4n^5+2\,019}$  bằng
	\choice 
		{\True  $-2$}
		{$4$} 
		{$+\infty$} 
		{$0$} 
	\loigiai{ 
		Ta có $\lim\dfrac{8n^5-2n^3+1}{2n^2-4n^5+2\,019}=\lim\dfrac{8-\dfrac{2}{n^2}+\dfrac{1}{n^5}}{\dfrac{2}{n^3}-4+\dfrac{2\,019}{n^5}}=\dfrac{8-0+0}{0-4+0}=-2$. 
	} 
\end{ex}

% nguồn: Ok_Toan11_TracNghiem_Lan2
\begin{ex}%[1D3N1-2]%[ma:83e4c5ca]
	Giá trị  của $\lim\dfrac{2n^2-3}{1-2n^2}$ bằng
	 \choice
		{$-3$} 
		{$2$} 
		{\True $-1$} 
		{$0$} 
	 \loigiai{
		Ta  có
	 $\lim\dfrac{2n^2-3}{1-2n^2}=\lim\dfrac{2-\dfrac{3}{n^2}}{\dfrac{1}{n^2}-2}=\dfrac{2-0}{0-2}=-1$.}
\end{ex}

% nguồn: Ok_Toan11_TracNghiem_Lan2
\begin{ex}%[1D3N1-2]%[ma:008230f7]
	Tính giới hạn  $\lim\dfrac{\sqrt{3n^2+1}+n}{1-2n^2}$.
	 \choice
		 {$-2$}
		{$-\dfrac{3}{2}$} 
		{$+\infty$} 
		{\True $0$} 
	 \loigiai{
		Ta có 
		\begin{align*}
			\lim\dfrac{\sqrt{3n^2+1}+n}{1-2n^2} 
			&=\lim\dfrac{n\sqrt{3+\dfrac{1}{n^2}}+n}{1-2n^2}\\ 
			&=\lim\dfrac{\dfrac{1}{n}  \sqrt{3+\dfrac{1}{n^2}}+\dfrac{1}{n}}{\dfrac{1}{n^2}-2}\\
			&=0. 
		\end{align*}
	 }
\end{ex}

% nguồn: Ok_Toan11_TracNghiem_Lan2
\begin{ex}%[1D3N1-2]%[ma:ccc92513]
	Tính $\lim\dfrac{2n+1}{1+n}$ được kết  quả là
	\choice 
		{\True $2$} 
		 {$0$}
		 {$\dfrac{1}{2}$}
		{$1$} 
	 \loigiai{
	 Ta có $\lim\dfrac{2n+1}{1+n}=\lim\dfrac{n\left(2+\dfrac{1}{n}\right)}{n\left(\dfrac{1}{n}+1\right)}=\lim\dfrac{2+\dfrac{1}{n}}{\dfrac{1}{n}+1}=\dfrac{2+0}{0+1}=2$.}
\end{ex}

% nguồn: Ok_TracNghiemLop11_Lan1
\begin{ex}%[1D3N1-2]%[ma:aa3597e3]
	Tính giới  hạn $L=\lim \dfrac{{n^2}+n+5}{2{n^2}+1}$.
	 \choice
		{$L=\dfrac{3}{2}$} 
		 {\True $L=\dfrac{1}{2}$}
		{$L=2$} 
		 {$L=1$}
	\loigiai{ 
		Ta có 
		\[L=\lim \dfrac{{n^2}+n+5}{2{n^2}+1}=\lim  \dfrac{1+\dfrac{1}{n}+\dfrac{5}{{n^2}}}{2+\dfrac{1}{{n^2}}}=\dfrac{1}{2}.\]
		Giải nhanh:  Ta sử dụng đánh giá $\dfrac{{n^2}+n+5}{2{n^2}+1}\sim \dfrac{{n^2}}{2{n^2}}=\dfrac{1}{2}.$
	} 
\end{ex}

% nguồn: Ok_TracNghiemLop11_Lan1
\begin{ex}%[1D3N1-2]%[ma:ca50d2f3]
	Kết quả của  giới hạn $\lim \dfrac{{{\pi}^n}+{3^n}+{2^{2n}}}{3{{\pi}^n}-{3^n}+{2^{2n+2}}}$ là
	\choice 
		{$1$} 
		 {$\dfrac{1}{3}$}
		 {$+\infty$}
		{\True  $\dfrac{1}{4}$}
	 \loigiai{
		Ta có $$\lim \dfrac{{{\pi}^n}+{3^n}+{2^{2n}}}{3{{\pi}^n}-{3^n}+{2^{2n+2}}}=\lim \dfrac{{{\left(\dfrac{\pi}{4} \right)}^n}+{{\left(\dfrac{3}{4}  \right)}^n}+1}{3.{{\left(\dfrac{\pi}{4} \right)}^n}-3.{{\left(\dfrac{3}{4} \right)}^n}+4}=\dfrac{1}{4}.$$
		Giải nhanh: Ta sử dụng đánh giá  $\dfrac{{{\pi}^n}+{3^n}+{2^{2n}}}{3{{\pi}^n}-{3^n}+{2^{2n+2}}}=\dfrac{{{\pi}^n}+{3^n}+{4^n}}{3{{\pi}^n}-{3^n}+{{4.4}^n}}\sim \dfrac{{4^n}}{{{4.4}^n}}=\dfrac{1}{4}$.
	 }
\end{ex}

% nguồn: Ok_TracNghiemLop11_Lan1
\begin{ex}%[1D3N1-2]%[ma:e5004611]
	 Tính giới hạn $\displaystyle\lim\dfrac{2n-1}{n-1}$.
	 \choice
		{$-2$} 
		 {$1$}
		{\True $2$} 
		{$-1$} 
	 \loigiai{
		Ta có $\lim\dfrac{2n-1}{n-1}=\lim\dfrac{2-\dfrac{1}{n}}{1-\dfrac{1}{n}} = 2$. 
	 }
\end{ex}

% nguồn: Ok_TracNghiemLop11_Lan1
\begin{ex}%[1D3N1-2]%[ma:a18f9fc1]
	Giá trị của giới hạn  $\lim \dfrac{-3}{4{n^2}-2n+1}$ là
	\choice 
		{$-\dfrac{3}{4}$} 
		 {$-\infty$}
		 {\True $0$}
		 {$-1$}
	\loigiai{ 
		Ta có $\lim \dfrac{-3}{4{n^2}-2n+1}=\lim \dfrac{\dfrac{-3}{{n^2}}}{4-\dfrac{2}{n}+\dfrac{1}{{n^2}}}=\dfrac{0}{4}=0$. 
	} 
\end{ex}

% nguồn: Ok_TracNghiemLop11_Lan1
\begin{ex}%[1D3N1-2]%[ma:95cfc216]
	Giá trị  của giới hạn $\lim \dfrac{3n^3-2n+1}{4n^4+2n+1}$ bằng
	\choice 
		 {$+\infty$}
		{$\dfrac{3}{4}$} 
		{\True  $0$}
		{$\dfrac{2}{7}$} 
	 \loigiai{
		Ta  có
		\begin{eqnarray*}
			\lim \dfrac{3n^3-2n+1}{4n^4+2n+1}&=&\lim\dfrac{n^3\left(3-\dfrac{2}{n^2}+\dfrac{1}{n^3}\right)}{n^4\left(4+\dfrac{2}{n^3}+\dfrac{1}{n^4}\right)}\\ 
			&=& \lim\left[\dfrac{1}{n}\cdot \dfrac{3-\dfrac{2}{n^2}+\dfrac{1}{n^3}}{4+\dfrac{2}{n^3}+\dfrac{1}{n^4}}\right]  \\
			&=&0\cdot \dfrac{3-0+0}{4+0+0}=0. 
		\end{eqnarray*}
	} 
\end{ex}

% nguồn: Ok_TracNghiemLop11_Lan1
\begin{ex}%[1D3N1-2]%[ma:ae000f06]
	Giá trị của giới  hạn $\lim \dfrac{3{n^3}-2n+1}{4{n^4}+2n+1}$ là
	\choice 
		{$+\infty$} 
		 {\True $0$}
		 {$\dfrac{2}{7}$}
		 {$\dfrac{3}{4}$}
	 \loigiai{
		Ta có $\lim  \dfrac{3{n^3}-2n+1}{4{n^4}+2n+1}=\lim \dfrac{\dfrac{3}{n}-\dfrac{2}{{n^2}}+\dfrac{1}{{n^4}}}{4+\dfrac{2}{{n^3}}+\dfrac{1}{{n^4}}}=\dfrac0{4}=0$.
	} 
\end{ex}

% nguồn: Ok_TracNghiemLop11_Lan1
\begin{ex}%[1D3N1-2]%[ma:d6b04f05]
	Cho hai dãy số $\left({u_n} \right)$ và $\left({v_n} \right)$  có ${u_n}=\dfrac{1}{n+1}$ và ${v_n}=\dfrac{1}{n+2}$. Khi đó $\lim \dfrac{{v_n}}{{u_n}}$ có giá trị bằng
	\choice 
		 {\True $1$}
		 {$2$}
		{$0$} 
		 {$3$}
	\loigiai{ 
		Ta có $\lim \dfrac{{v_n}}{{u_n}}=\lim  \dfrac{n+1}{n+2}=\lim \dfrac{1+\dfrac{1}{n}}{1+\dfrac{2}{n}}=\dfrac{1}{1}=1$.\\
		Giải nhanh: Ta sử dụng đánh giá $\dfrac{n+1}{n+2}\sim \dfrac{n}{n}=1$. 
	 }
\end{ex}

% nguồn: Ok_TracNghiemLop11_Lan1
\begin{ex}%[1D3N1-2]%[ma:1ff58e0d]
	 Giá trị của $\lim\dfrac{2n+1}{n+1}$ bằng
	 \choice
		{$0$} 
		{$3$} 
		 {$1$}
		 {\True $2$}
	\loigiai{ 
		 Ta có
		\[\lim\dfrac{2n+1}{n+1} 
		=\lim\dfrac{2+\dfrac{1}{n}}{1+\dfrac{1}{n}} 
		 =2.\]
	 }
\end{ex}

% nguồn: Ok_TracNghiemLop11_Lan1
\begin{ex}%[1D3N1-2]%[ma:601ec2eb]
	Tính $L = \lim \dfrac{2n+1}{n^2 +  n + 3}$.
	\choice 
		{$L = 2$} 
		 {$L = \dfrac{2}{3}$}
		{\True $L  = 0$}
		{$L  = \dfrac{1}{3}$}
	\loigiai{Ta có $L = \lim \dfrac{2n+1}{n^2 + n + 3} = \lim \dfrac{\tfrac{2}{n}+\tfrac{1}{n^2}}{1 + \tfrac{1}{n} + \tfrac{3}{n^2}} = 0$. 
	} 
\end{ex}

% nguồn: Ok_TracNghiemLop11_Lan1
\begin{ex}%[1D3N1-2]%[ma:affddf91]
	 Tính $\lim \dfrac{2n-3}{2-n}$.
	\choice 
		 {$-\dfrac{3}{2}$}
		{$1$} 
		{$3$} 
		 {\True $-2$}
	\loigiai{ 
		 Ta có $\lim \dfrac{n\left(2-\dfrac{3}{n}\right)}{n\left(\dfrac{2}{n}-1\right)}=\lim \dfrac{2-\dfrac{3}{n}}{\dfrac{1}{n}-1}=\dfrac{2}{-1}=-2$.
	 }
\end{ex}

% nguồn: Ok_TracNghiemLop11_Lan1
\begin{ex}%[1D3N1-2]%[ma:11470b24]
	Dãy  số nào có giới hạn bằng $0$?
	 \choice
		{$\left(\dfrac{5}{4}\right)^n$} 
		 {$\left(\dfrac{\pi}{3}\right)^n$}
		{\True  $\left(-\dfrac{\pi}{4}\right)^n$}
		{$\left(-\dfrac{4}{3}\right)^n$} 
	 \loigiai{
		Vì $\Big\vert -\dfrac{\pi}{4} \Big\vert \approx 0{,}78<1$  nên $\displaystyle\lim\limits \left(-\dfrac{\pi}{4}\right)^n =0$.
	 }
\end{ex}

% nguồn: Ok_TracNghiemLop11_Lan1
\begin{ex}%[1D3N1-2]%[ma:67990b5d]
	Tính $L=\displaystyle \lim  \dfrac{n-1}{n^3 +3}$
	\choice{$ L=2  $}
		{$L=3  $}
		{$L= 1 $} 
		{\True  $L= 0 $}
	\loigiai{$L=\displaystyle \lim  \dfrac{n-1}{n^3 +3}=\displaystyle \lim \dfrac{\dfrac{1}{n^2}-\dfrac{1}{n^3}}{1 +\dfrac{1}{n^3}}=\dfrac{0}{1}=0.$
	 }
\end{ex}

% nguồn: Ok_TracNghiemLop11_Lan1
\begin{ex}%[1D3N1-2]%[ma:dbfff29d]
	Tìm giới hạn sau  $\lim\limits_{n \to +\infty} \dfrac{4n -1}{-n -1}$.
	 \choice
		 {\True $-4$}
		{$7$} 
		{$-7$} 
		 {$4$}
	\loigiai{ 
		Chia cả tử và mẫu  cho $n$ ta được
		$$\lim\limits_{n \to +\infty} \dfrac{4n  -1}{-n -1} = -4.$$
		Vậy $\lim\limits_{n  \to +\infty} \dfrac{4n -1}{-n -1} = -4$.
	 }
\end{ex}

% nguồn: Ok_TracNghiemLop11_Lan1
\begin{ex}%[1D3N1-2]%[ma:989779cf]
	Trong các giới hạn sau, giới hạn nào bằng  $0$?
	 \choice
		{$ \lim(n^3-3n+1)  $}
		{\True $  \lim\dfrac{n^2+n}{n^3+1} $}
		{$  \lim\dfrac{n^2+n+1}{4n+1} $}
		{$ \lim\dfrac{2^n-3^n}{3^n+2} $} 
	 \loigiai{
		Ta  có $ \lim\dfrac{n^2+n}{n^3+1}=\lim\dfrac{\dfrac{1}{n}+\dfrac{1}{n^2}}{1+\dfrac{1}{n^3}}=0. $
	 }
\end{ex}

% nguồn: Ok_Toan11_TracNghiem_Lan2
\begin{ex}%[1D3H1-2]%[ma:5132a1c1]
	$\lim\dfrac {2n+1}{n-1}$  bằng
	\choice 
		{\True $2$} 
		 {$+\infty$}
		 {$-\infty$}
		{$1$} 
	 \loigiai{$\lim\dfrac {2n+1}{n-1}=\lim\dfrac {2+\dfrac {1}{n}}{1-\dfrac {1}{n}}=\dfrac {2+0}{1-0}=2$.
	 }
\end{ex}

% nguồn: Ok_Toan11_TracNghiem_Lan2
\begin{ex}%[1D3H1-2]%[ma:8c6954df]
	Cho $a\in \mathbb{R}$ sao cho giới hạn $\lim\dfrac{an^2+a^2n+1}{\left(n+1\right)^2}=a^2-a+1$.  Khi đó khẳng định nào sau đây là đúng?
	\choice 
		{\True $0<a<2$} 
		{$0<a<\dfrac{1}{2}$} 
		{$-1<a<0$} 
		{$1<a<3$} 
	 \loigiai{
		Ta  có $\lim\dfrac{an^2+a^2n+1}{\left(n+1\right)^2}=\lim\dfrac{an^2+a^2n+1}{n^2+2n+1}=\lim\dfrac{a+\dfrac{a^2}{n}+\dfrac{1}{n^2}}{1+\dfrac{2}{n}+\dfrac{1}{n^2}}=\dfrac{a+0+0}{1+0+0}=a$.
		\\ 
		Suy ra  $a^2-a+1=a\Rightarrow a^2-2a+1=0\Rightarrow a=1$.
	 }
\end{ex}

% nguồn: Ok_Toan11_TracNghiem_Lan2
\begin{ex}%[1D3H1-2]%[ma:a67c882a]
	 Tìm $\lim\dfrac{8n^5-2n^3+1}{4n^5+2n^2+1}$.
	\choice 
		{\True $2$} 
		 {$8$}
		{$1$} 
		{$4$} 
	 \loigiai{
	Ta có $\lim\dfrac{8n^5-2n^3+1}{4n^5+2n^2+1}=\lim\dfrac{n^5\left(8-\dfrac{2}{n^2}+\dfrac{1}{n^5}\right)}{n^5\left(4+\dfrac{2}{n^3}+\dfrac{1}{n^5}\right)}=\lim\dfrac{8-\dfrac{2}{n^2}+\dfrac{1}{n^5}}{4+\dfrac{2}{n^3}+\dfrac{1}{n^5}}=\dfrac{8-0+0}{4+0+0}=\dfrac{8}{4}=2$.} 
\end{ex}

% nguồn: Ok_Toan11_TracNghiem_Lan2
\begin{ex}%[1D3H1-2]%[ma:42b6bc5e]
	 Giá trị $\lim\dfrac{1+19n}{18n+19}$ bằng
	\choice 
		 {\True $\dfrac{19}{18}$}
		 {$\dfrac{1}{18}$}
		 {$+\infty$}
		{$\dfrac{1}{19}$} 
	\loigiai{Ta có  $\lim\dfrac{1+19n}{18n+19}=\lim\dfrac{\dfrac{1}{n}+19}{18+\dfrac{19}{n}}=\dfrac{19}{18}$.
	} 
\end{ex}

% nguồn: Ok_Toan11_TracNghiem_Lan2
\begin{ex}%[1D3H1-2]%[ma:e48b0bf9]
	Biết $\lim\dfrac{2n^4+n+1}{n-n^3+3n^4}=\dfrac{a}{b}$ với $\dfrac{a}{b}$  là phân số tối giản. Tính $a+2b$.
	\choice 
		 {\True $8$}
		{$7$} 
		 {$6$}
		 {$9$}
	\loigiai{Ta  có $\lim\dfrac{2n^4+n+1}{n-n^3+3n^4}=\lim\dfrac{n^4\left(2+\dfrac{1}{n^3}+\dfrac{1}{n^4}\right)}{n^4\left(\dfrac{1}{n^4}-\dfrac{1}{n}+3\right)}=\lim\dfrac{2+\dfrac{1}{n^3}+\dfrac{1}{n^4}}{\dfrac{1}{n^4}-\dfrac{1}{n}+3}=\dfrac{2}{3}$.\\
		Suy ra $a=2$; $b=3$.\\ 
		Vậy  $a+2b=8$.
	 }
\end{ex}

% nguồn: Ok_Toan11_TracNghiem_Lan2
\begin{ex}%[1D3H1-2]%[ma:e929b00c]
	Tính giới hạn $\lim\left[\dfrac {1}{1\cdot 2}+\dfrac {1}{2\cdot 3}+\dfrac {1}{3\cdot 4}+\cdots  +\dfrac {1}{n\left(n+1 \right)}\right]$.
	\choice 
		 {$0$}
		{$2$} 
		{\True  $1$}
		{$\dfrac  {3}{2}$}
	\loigiai{Ta  có
		\begin{align*}
			&\dfrac {1}{1\cdot  2}+\dfrac {1}{2\cdot 3}+\dfrac {1}{3\cdot 4}+\cdots +\dfrac {1}{n(n+1)}\\
			=&\dfrac {1}{1}-\dfrac {1}{2}+\dfrac {1}{2}-\dfrac  {1}{3}+\cdots +\dfrac {1}{n-1}-\dfrac {1}{n}+\dfrac {1}{n}-\dfrac {1}{n+1}\\
			 =&1-\dfrac {1}{n+1}.
		\end{align*}
		Vậy $\lim\left[\dfrac {1}{1\cdot 2}+\dfrac {1}{2\cdot 3}+\dfrac {1}{3\cdot  4}+\cdots +\dfrac {1}{n(n+1)}\right]=\lim\left(1-\dfrac {1}{n+1} \right)=1$.
	 }
\end{ex}

% nguồn: Ok_Toan11_TracNghiem_Lan2
\begin{ex}%[1D3H1-2]%[ma:746ea974]
	 Tính $\lim\dfrac {1+19n}{18n+19}$.
	\choice 
		{$+\infty  $}
		 {$\dfrac {1}{19}$}
		{$\dfrac  {1}{18}$}
		{\True  $\dfrac {19}{18}$}
	\loigiai{Ta có $\lim\dfrac {1+19n}{18n+19}=\lim\dfrac {\dfrac  {1}{n}+19}{18+\dfrac {19}{n}}=\dfrac {0+19}{18+0}=\dfrac {19}{18}$.
	 }
\end{ex}

% nguồn: Ok_Toan11_TracNghiem_Lan2
\begin{ex}%[1D3H1-2]%[ma:933f8ea1]
	Tính giới hạn $L=\lim\dfrac {2n+1}{2+n-n^2}$. 
	\choice 
		{$L=-\infty $} 
		 {$L=-2$}
		 {$L=1$}
		{\True  $L=0$}
	\loigiai{Ta có $L=\lim\dfrac {2n+1}{2+n-n^2}=\lim\dfrac  {\dfrac {2}{n}+\dfrac {1}{n^2}}{\dfrac {2}{n^2}+\dfrac {1}{n}-1}=\dfrac {0+0}{0+0-1}=0$.
	 }
\end{ex}

% nguồn: Ok_Toan11_TracNghiem_Lan2
\begin{ex}%[1D3H1-2]%[ma:480b9baa]
	Tìm $\lim\left(\dfrac{1}{n^2}+\dfrac{2}{n^2}+\ldots+\dfrac{n}{n^2} \right)$. 
	 \choice
		 {$+\infty$}
		{\True  $\dfrac{1}{2}$}
		{$\dfrac{1}{n}$} 
		{$0$} 
	\loigiai{  Ta có
		$\lim\left(\dfrac{1}{n^2}+\dfrac{2}{n^2}+\ldots+\dfrac{n}{n^2} \right)=\lim\left(\dfrac{1+2+\ldots+n}{n^2} \right)=\lim\dfrac{n\left(n+1\right)}{2n^2}=\lim\dfrac{1+\dfrac{1}{n}}{2} =\dfrac{1}{2}$. 
	} 
\end{ex}

% nguồn: Ok_Toan11_TracNghiem_Lan2
\begin{ex}%[1D3H1-2]%[ma:d12a1b0c]
	$\lim\left(\dfrac{1}{n^2}+\dfrac{2}{n^2}+\dfrac{3}{n^2}+\ldots+\dfrac{n}{n^2} \right)$  bằng
	 \choice
		 {$1$}
		{$0$} 
		 {$\dfrac{1}{3}$}
		 {\True $\dfrac{1}{2}$}
	\loigiai{ Ta có 
		\begin{eqnarray*}
			 \lim\left(\dfrac{1}{n^2}+\dfrac{2}{n^2}+\dfrac{3}{n^2}+\ldots+\dfrac{n}{n^2} \right)&=&\lim\dfrac{1+2+3+\ldots+n}{n^2} \\&=&\lim\dfrac{n(n+1)}{2n^2}\\&=& \lim\left(\dfrac{1}{2}+\dfrac{1}{2n} \right)=\dfrac{1}{2}.
		\end{eqnarray*}
	} 
\end{ex}

% nguồn: Ok_Toan11_TracNghiem_Lan2
\begin{ex}%[1D3H1-2]%[ma:5d55204d]
	Tìm $ a$  để $\lim\dfrac {an^2-3n}{9n^2+5}=\dfrac {2}{3}$.
	 \choice
		{$a=4$} 
		{\True $a=6$} 
		 {$a=8$}
		 {$a=9$}
	 \loigiai{Ta có
		\begin{eqnarray*}
			 & & \lim\dfrac {an^2-3n}{9n^2+5}=\dfrac {2}{3}\\
			 &\Leftrightarrow & \lim\dfrac {a-\dfrac {3}{n}}{9+\dfrac {5}{n^2}}=\dfrac {2}{3}\\
			&\Leftrightarrow &  \dfrac {a}{9}=\dfrac {2}{3}\\
			&\Leftrightarrow & a=6 
		\end{eqnarray*}
		 Vậy $a=6$.
	} 
\end{ex}

% nguồn: Ok_Toan11_TracNghiem_Lan2
\begin{ex}%[1D3H1-2]%[ma:2cb8c6da]
	Tính  giới hạn $\lim\dfrac{3\cdot 2^{n+1}-2\cdot 3^{n+1}}{4+3^{n}}$.
	\choice 
		 {$\dfrac{3}{2}$}
		 {$0$}
		{$\dfrac{6}{5}$} 
		 {\True $-6$}
	\loigiai{ 
		 Ta có $\lim\dfrac{3\cdot 2^{n+1}-2\cdot 3^{n+1}}{4+3^{n}}=\lim\dfrac{6\cdot\left(\dfrac{2}{3}\right)^{n}-6}{4\cdot\left(\dfrac{1}{3}\right)^{n}+1}=-6$.
	 }
\end{ex}

% nguồn: Ok_Toan11_TracNghiem_Lan2
\begin{ex}%[1D3H1-2]%[ma:5eb96da6]
	Tính  $L=\lim\dfrac {n-1}{n^3+3}$.
	\choice 
		 {$L=2$}
		{$L=3$} 
		{\True  $L=0$}
		 {$L=1$}
	\loigiai{Ta có $L=\lim\dfrac {n-1}{n^3+3}=\lim\dfrac {n^3\left( \dfrac {1}{n^2}-\dfrac {1}{n^3} \right)}{n^3\left(  1+\dfrac {3}{n^3} \right)}=\lim\dfrac {\dfrac {1}{n^2}-\dfrac {1}{n^3}}{1+\dfrac {3}{n^3}}=\dfrac {0-0}{1+0}=0$.
	} 
\end{ex}

% nguồn: Ok_Toan11_TracNghiem_Lan2
\begin{ex}%[1D3H1-2]%[ma:9f4f9583]
	Dãy số $(u_n)$  nào sau đây có giới hạn bằng $\dfrac {1}{5}$?
	\choice 
		{ $u_n=\dfrac  {1-2n^2}{5n+5}$}
		{  $u_n=\dfrac {1-2n}{5n+5n^2}$}
		{ $u_n=\dfrac  {n^2-2n}{5n+5n^2}$}
		{ $u_n=\dfrac  {1-2n}{5n+5}$}
	\loigiai{Ta có $\lim\dfrac {n^2-2n}{5n+5n^2}$$=\lim\dfrac {n^2\left( 1-\dfrac {2}{n} \right)}{n^2\left( \dfrac {5}{n}+5  \right)}=\lim\dfrac {1-\dfrac{2}{n} }{\dfrac {5}{n}+5}=\dfrac {1-0}{0+5}=\dfrac {1}{5}$.
	 }
\end{ex}

% nguồn: Ok_Toan11_TracNghiem_Lan2
\begin{ex}%[1D3H1-2]%[ma:7d999ac5]
	Kết  quả của $I=\lim\dfrac{2n+2020}{3n+2021}$.
	 \choice
		 {$I=\dfrac{3}{2}$}
		 {\True $I=\dfrac{2}{3}$}
		 {$I=\dfrac{2020}{2021}$}
		 {$I=1$}
	 \loigiai{
		Ta có $I=\lim\dfrac{2n+2020}{3n+2021}=\lim\dfrac{2+\dfrac{2020}{n}}{3+\dfrac{2021}{n}}=\dfrac{2+0}{3+0}=\dfrac{2}{3}$. 
	} 
\end{ex}

% nguồn: Ok_Toan11_TracNghiem_Lan2
\begin{ex}%[1D3H1-2]%[ma:e2780cd8]
	$\lim\left(3^{n}-4^{n}\right)$ là 
	\choice 
		{$+\infty$} 
		{\True $-\infty$} 
		 {$\dfrac{4}{3}$}
		{1} 
	 \loigiai{
		 Ta có $\lim\left(3^{n}-4^{n}\right)=\lim 4^{n}\left[\left(\dfrac{3}{4}\right)^{n}-1\right]=-\infty$.
	 }
\end{ex}

% nguồn: Ok_Toan11_TracNghiem_Lan2
\begin{ex}%[1D3H1-2]%[ma:0859f9f9]
	Tính  $\lim\dfrac{\sqrt[3]{n^3+2023}}{2n+2024}$.
	 \choice
		{$\dfrac{\sqrt[3]{2023}}{2024}$} 
		{\True $\dfrac{1}{2}$} 
		{$+\infty$} 
		 {$0$}
	\loigiai{Ta có $\lim\dfrac{\sqrt[3]{n^3+2023}}{2n+2024}=\lim\dfrac{\sqrt[3]{1+\dfrac{2023}{n^3}}}{2+\dfrac{2024}{n}}=\dfrac{1}{2}$. 
	} 
\end{ex}

% nguồn: Ok_Toan11_TracNghiem_Lan2
\begin{ex}%[1D3H1-2]%[ma:94be6129]
	$\lim {\dfrac{100^{n+1}+3\cdot 99^{n}}{10^{2n}-2\cdot 98^{n+1}}}$ là 
	 \choice
		 {$+\infty$}
		 {\True $100$}
		 {$\dfrac{1}{100}$}
		{$0$} 
	 \loigiai{
		Ta có 
		$\lim\dfrac{100^{n+1}+3\cdot 99^{n}}{10^{2n}-2\cdot  98^{n+1}}=\lim\dfrac{100+3\cdot\left(\dfrac{99}{100}\right)^{n}}{1-2\cdot\left(\dfrac{98}{100}\right)^{n}}=100$.
	} 
\end{ex}

% nguồn: Ok_Toan11_TracNghiem_Lan2
\begin{ex}%[1D3H1-2]%[ma:96371b3e]
	Dãy số  nào sau đây có giới hạn bằng $0$?
	\choice 
		{$u_n=\dfrac {n^2-2}{5n+3n^2}$} 
		{$u_n=\dfrac  {n^2-2n}{5n+3n^2}$}
		{\True $ u_n=\dfrac  {1-2n}{5n+3n^2}$}
		{$u_n=\dfrac {1-2n^2}{5n+3n^2}$} 
	\loigiai{\begin{itemize}
		\item $\lim\dfrac {n^2-2}{5n+3n^2}=\lim\dfrac {1-\dfrac {2}{n^2}}{\dfrac  {5}{n}+3}=\dfrac {1-0}{0+3}=\dfrac {1}{3}$.
		\item $\lim\dfrac {n^2-2n}{5n+3n^2}=\lim\dfrac {1-\dfrac  {2}{n}}{\dfrac {5}{n}+3}=\dfrac {1-0}{0+3}=\dfrac {1}{3}$.
		\item $\lim\dfrac {1-2n}{5n+3n^2}=\lim\dfrac {\dfrac {1}{n^2}-\dfrac  {2}{n}}{\dfrac {5}{n}+3}=\dfrac {0-0}{0+3}=0$.
		 \item $\lim\dfrac {1-2n^2}{5n+3n^2}=\lim\dfrac {\dfrac {1}{n^2}-2}{\dfrac {5}{n}+3}=\dfrac {0-2}{0+3}=-\dfrac {2}{3}$.
	\end{itemize}
 }
\end{ex}

% nguồn: Ok_Toan11_TracNghiem_Lan2
\begin{ex}%[1D3H1-2]%[ma:1497a30a]
	Tính $A=\lim\left(3+\dfrac  {1}{n^2} \right)$.
	 \choice
		{\True  $A=3$}
		 {$A=-\infty$}
		{$A=+\infty$} 
		{$A=0$} 
	\loigiai{Ta có $A=\lim\left(3+\dfrac {1}{n^2} \right)=3+0=3$. 
	 }
\end{ex}

% nguồn: Ok_Toan11_TracNghiem_Lan2
\begin{ex}%[1D3H1-2]%[ma:9c06ba32]
	 Tính $\lim\sqrt{\dfrac{1^2+2^2+3^2+\ldots+n^2}{2n\left(n+7\right)\left(6n+5\right)}}$.
	 \choice
		{\True $\dfrac{1}{6}$} 
		{$\dfrac{1}{2\sqrt{6}}$} 
		 {$\dfrac{1}{2}$}
		{$+\infty$} 
	 \loigiai{
		Ta  có $1^2+2^2+3^2+\ldots+n^2=\dfrac{n\left(n+1\right)\left(2n+1\right)}{6}$.\\
		Khi  đó
		\begin{eqnarray*}
			\lim\sqrt{\dfrac{1^2+2^2+3^2+\ldots+n^2}{2n\left(n+7\right)\left(6n+5\right)}}&=&\lim\sqrt{\dfrac{n\left(n+1\right)\left(2n+1\right)}{12n\left(n+7\right)\left(6n+5\right)}}\\ 
			&=&\lim\sqrt{\dfrac{\left(1+\dfrac{1}{n} \right)\left(2+\dfrac{1}{n} \right)}{12\left(1+\dfrac{7}{n} \right)\left(6+\dfrac{5}{n} \right)}}\\&=&\sqrt{\dfrac{\left(1+0 \right)\left(2+0 \right)}{12\left(1+0 \right)\left(6+0  \right)}}\\&=&\dfrac{1}{6}.
		\end{eqnarray*}
	 }
\end{ex}

% nguồn: Ok_Toan11_TracNghiem_Lan2
\begin{ex}%[1D3H1-2]%[ma:e5abc1c4]
	$\lim\dfrac{1}{2n+7}$ bằng 
	\choice 
		 {$\dfrac{1}{7}$}
		{$+\infty $} 
		{$\dfrac{1}{2}$} 
		 {\True $0$}
	\loigiai{ 
	Ta  có $\lim\dfrac{1}{2n+7}$ $=\lim\dfrac{\dfrac{1}{n}}{2+\dfrac{7}{n}}=0$ .}
\end{ex}

% nguồn: Ok_Toan11_TracNghiem_Lan2
\begin{ex}%[1D3H1-2]%[ma:4864f90c]
	Giá trị của tham số $a$ để $\lim\dfrac{2n^2-n+4}{an^2+n+3}=\dfrac{4}{3}$ là 
	 \choice
		{$a=\dfrac{8}{3}$} 
		{$a=\dfrac{2}{3}$} 
		{\True  $a=\dfrac{3}{2}$}
		{$a=6$} 
	 \loigiai{Ta có $\lim\dfrac{2n^2-n+4}{an^2+n+3}=\lim\dfrac{2-\dfrac{1}{n}+\dfrac{4}{n^2}}{a+\dfrac{1}{n}+\dfrac{3}{n^2}}=\dfrac{2}{a}$.\\
		Theo đề bài thì  $\dfrac{2}{a}=\dfrac{4}{3}\Leftrightarrow a=\dfrac{3}{2}$.
	 }
\end{ex}

% nguồn: Ok_Toan11_TracNghiem_Lan2
\begin{ex}%[1D3H1-2]%[ma:bc3dd48d]
	Tính giá trị của $\lim\dfrac {\cos  n+\sin n}{n^2+1}$
	\choice 
		 {$1$}
		{\True  $0$}
		 {$+\infty$}
		{$-\infty$} 
	\loigiai{Ta  có $0<\left|\dfrac {\cos n+\sin n}{n^2+1} \right|\le \dfrac {|\cos n|+|\sin n|}{n^2+1}<\dfrac {2}{n^2+1}$ và $\lim\dfrac {2}{n^2+1}=0$.\\
		Suy  ra $\lim\dfrac {\cos n+\sin n}{n^2+1}=0$.
	 }
\end{ex}

% nguồn: Ok_Toan11_TracNghiem_Lan2
\begin{ex}%[1D3H1-2]%[ma:57102262]
	Giới hạn $\lim\dfrac{3\sqrt{5n^2+n}}{2(3n+1)}=\dfrac{a\sqrt{5}}{b}$  (với $a,b\in\mathbb{N}^{*}$ và $\dfrac{a}{b}$ là phân số tối giản). Tổng $T=a+b$ bằng
	 \choice
		{$9$} 
		 {\True $3$}
		{$8$} 
		 {$5$}
	 \loigiai{Ta có $$
		\lim\dfrac{3\sqrt{5n^2+n}}{2(3n+1)}=\lim\dfrac{3\sqrt{5+\dfrac{1}{n}}}{2\left(3+\dfrac{1}{n}\right)}=\dfrac{1\sqrt{5}}{2}.$$ 
		Suy ra $a=1$,  $b=2$.\\
		Vậy $T=3$. 
	} 
\end{ex}

% nguồn: Ok_Toan11_TracNghiem_Lan2
\begin{ex}%[1D3H1-2]%[ma:03c308c4]
	Biết  $\lim\dfrac{4n+1}{2n+5}=a$ và $\lim\dfrac{3^n+4^n}{4^{n+1}+3}=b$. Tổng $a+4b$ bằng
	 \choice
		{$5$} 
		{$4$} 
		{$2$} 
		{\True $3$} 
	\loigiai{Ta có \begin{itemize}
		\item  $\lim\dfrac{4n+1}{2n+5}=\lim\dfrac{4+\dfrac{1}{n}}{2+\dfrac{5}{n}}=2\Rightarrow a=2$,
		 \item $\lim\dfrac{3^n+4^n}{4^{n+1}+3}=\lim\dfrac{3^n+4^n}{4.4^n+3}=\lim\dfrac{\left(\dfrac{3}{4}\right)^n+1}{4+3.\left(\dfrac{1}{4}\right)^n}=\dfrac{1}{4}\Rightarrow b=\dfrac{1}{4}$.
	\end{itemize}
	 Vậy $a+4b=3$.
 }
\end{ex}

% nguồn: Ok_Toan11_TracNghiem_Lan2
\begin{ex}%[1D3H1-2]%[ma:b53d347e]
	 ${\lim\dfrac{2{n^2}-3}{{n^6}+5{n^5}}}$ bằng
	\choice 
		{$2$} 
		{\True  $0$}
		 {$\dfrac {-3}{5}$}
		{$-3$} 
	\loigiai{ 
		Ta có $\lim\dfrac  {2n^2-3}{n^6+5n^5}=\lim\dfrac {\dfrac {2}{n^4}-\dfrac {3}{n^6}}{1+\dfrac {5}{n}}=\dfrac {0-0}{1+0}=0$.
	} 
\end{ex}

% nguồn: Ok_Toan11_TracNghiem_Lan2
\begin{ex}%[1D3H1-2]%[ma:28479b7e]
	Biết $\lim\dfrac{8n+1}{an-2}=4$ với $a\ne 0$ là tham số.  Khi đó $a-a^2$ bằng
	 \choice
		 {$-4$}
		 {$-6$}
		 {$2$}
		 {\True $-2$}
	 \loigiai{
		Ta có $\lim\dfrac{8n+1}{an-2}=\lim\dfrac{n\left(8+\dfrac{1}{n} \right)}{n\left(a-\dfrac{2}{n} \right)}=\lim\dfrac{8+\dfrac{1}{n}}{a-\dfrac{2}{n}}=\dfrac{8+0}{a-0}=\dfrac{8}{a}=4\Rightarrow a=2$.\\ 
		Khi  đó $a-a^2=2-2^2=-2$.
	 }
\end{ex}

% nguồn: Ok_Toan11_TracNghiem_Lan2
\begin{ex}%[1D3H1-2]%[ma:85d89443]
	Tìm $\lim u_n$  biết $u_n=\dfrac{n\sqrt{1+3+5+\ldots+\left(2n-1\right)}}{2n^2+1}$.
	\choice 
		 {\True $\dfrac{1}{2}$}
		 {$+\infty$}
		{$1$} 
		{$-\infty$} 
	\loigiai{  Ta có $1+3+5+\ldots+\left(2n-1\right)=\dfrac{n}{2}\cdot2n=n^2$.\\
		 Suy ra
		\begin{eqnarray*}
			\lim  u_n&=&\lim\dfrac{n\sqrt{1+3+5+\ldots+\left(2n-1\right)}}{2n^2+1}\\
			&=&\lim\dfrac{n\sqrt{n^2}}{2n^2+1}\\&=&\lim\dfrac{n^2}{2n^2+1}\\&=&\lim\dfrac{1}{2+\dfrac{1}{n^2}}\\&=&\dfrac{1}{2+0}\\&=&\dfrac{1}{2}. 
		\end{eqnarray*}
	} 
\end{ex}

% nguồn: Ok_Toan11_TracNghiem_Lan2
\begin{ex}%[1D3H1-2]%[ma:6adcabe4]
	$\lim\dfrac {2n+1}{n-1}$  bằng
	\choice 
		{$1$} 
		{\True $2$} 
		{$-1$} 
		{$-2$} 
	\loigiai{$\lim\dfrac {2n+1}{n-1}=\lim\dfrac {n\left( 2+\dfrac {1}{n} \right)}{n\left(  1-\dfrac {1}{n} \right)}=\lim\dfrac {2+\dfrac {1}{n}}{1-\dfrac {1}{n}}=\dfrac {2+0}{1-0}=2$.
	 }
\end{ex}

% nguồn: Ok_Toan11_TracNghiem_Lan2
\begin{ex}%[1D3H1-2]%[ma:10902407]
	Gọi $S$ là tập hợp các tham  số nguyên $a$ thỏa mãn $\lim\left(\dfrac{3n+2}{n+2}+a^2-4a\right)=0$. Tổng các phần tử của $S$ bằng
	 \choice
		{\True  $4$}
		{$3$} 
		{$5$} 
		 {$2$}
	 \loigiai{
		Ta  có
		\begin{eqnarray*}
			\lim\left(\dfrac{3n+2}{n+2}+a^2-4a\right) 
			&=&\lim\dfrac{\left(a^2-4a+3\right)n+2+2a^2-8a}{n+2}\\&=&\lim\dfrac{a^2-4a+3+\dfrac{2+2a^2-8a}{n}}{1+\dfrac{2}{n}}  =a^2-4a+3.
		\end{eqnarray*}
		Theo giả  thiết $\lim\left(\dfrac{3n+2}{n+2}+a^2-4a\right)=0\Leftrightarrow a^2-4a+3=0\Leftrightarrow \hoac{&a=3\\&a=1.}$\\
		 Vậy $S=\left\{1; 3\right\}$ suy ra tổng các phần tử của $S$ là $1+3=4$.
	} 
\end{ex}

% nguồn: Ok_Toan11_TracNghiem_Lan2
\begin{ex}%[1D3H1-2]%[ma:cb84fb87]
	Tính $\lim\dfrac{\left(2n+1\right)^6}{\left(n+2\right)^4\left(2n-1\right)^2}$. 
	\choice 
		 {$\dfrac{1}{16}$}
		{$15$} 
		{$8$} 
		 {\True $16$}
	\loigiai{ 
		Ta có \begin{align*}
		\lim\dfrac{\left(2n+1\right)^6}{\left(n+2\right)^4\left(2n-1\right)^2}=\lim\dfrac{n^6\left(2+\dfrac{1}{n}  \right)^6}{n^6\left(1+\dfrac{2}{n} \right)^4\left(2-\dfrac{1}{n} \right)^2}=\lim\dfrac{\left(2+\dfrac{1}{n} \right)^6}{\left(1+\dfrac{2}{n} \right)^4\left(2-\dfrac{1}{n} \right)^2}=2^4=16.
	\end{align*}
} 
\end{ex}

% nguồn: Ok_Toan11_TracNghiem_Lan2
\begin{ex}%[1D3H1-2]%[ma:bc1ccd8b]
	Giới hạn nào  có kết quả bằng $0$?
	\choice 
		{\True  $\lim\dfrac{1+2\cdot 2\,017^{n}}{2\,016^{n}+2\,018^{n}}$}
		{$\lim\dfrac{1+2\cdot  2\,018^{n}}{2\,016^{n}+2\,017^{n+1}}$}
		 {$\lim\dfrac{1+2\cdot 2\,018^{n}}{2\,017^{n}+2\,018^{n}}$}
		{$\lim\dfrac{2\cdot  2\,018^{n+1}-2\,018}{2\,016^{n}+2\,018^{n}}$}
	 \loigiai{
		Ta có $\lim\dfrac{1+2\cdot  2\,017^{n}}{2\,016^{n}+2\,018^{n}}=\lim\dfrac{\left(\dfrac{1}{2\,018}\right)^{n}+2\cdot\left(\dfrac{2\,017}{2\,018}\right)^{n}}{\left(\dfrac{2\,016}{2\,018}\right)^{n}+1}=0$.
	} 
\end{ex}

% nguồn: Ok_Toan11_TracNghiem_Lan2
\begin{ex}%[1D3H1-2]%[ma:60cd4773]
	Tính  $\lim\dfrac{\sqrt{4n^2+2}+n}{n-1}$.
	 \choice
		{$2$} 
		{\True  $3$}
		{$+\infty$} 
		{$5$} 
	 \loigiai{Ta có $\lim\dfrac{\sqrt{4n^2+2}+n}{n-1}=\lim\dfrac{\sqrt{4+\dfrac{2}{n^2}}+1}{1-\dfrac{1}{n}}=3$.
	} 
\end{ex}

% nguồn: Ok_Toan11_TracNghiem_Lan2
\begin{ex}%[1D3H1-2]%[ma:c3442e05]
	Biết $\lim\dfrac {an+4}{4n+3}=-2$  tìm $2a+1$.
	\choice 
		{$2a+1=-7$} 
		{$2a+1=-8$} 
		{\True  $2a+1=-15$}
		 {$2a+1=-17$}
	\loigiai{Ta có \begin{eqnarray*}
		& & \lim\dfrac  {an+4}{4n+3}=-2\\
		&\Leftrightarrow  & \lim\dfrac {a+\dfrac {4}{n}}{4+\dfrac {3}{n}}=-2\\
		&\Leftrightarrow &  \dfrac {a}{4}=-2\\
		&\Leftrightarrow &  a=-8.
	\end{eqnarray*}
	Suy ra  $2a+1=-15$.
 }
\end{ex}

% nguồn: Ok_Toan11_TracNghiem_Lan2
\begin{ex}%[1D3H1-2]%[ma:682a23e2]
	${\lim\dfrac{2\,018}{n}}$ bằng 
	\choice 
		{$-\infty $} 
		{\True  $0$}
		{$1$} 
		 {$+\infty $}
	\loigiai{${\lim\dfrac{2\,018}{n}}=0$. 
	} 
\end{ex}

% nguồn: Ok_Toan11_TracNghiem_Lan2
\begin{ex}%[1D3H1-2]%[ma:30393094]
	Tính giới hạn $J=\lim\dfrac {(n+1)(2n+3)}{n^3-2}$. 
	\choice 
		{$J=-\dfrac {3}{2}$} 
		{$J=2$} 
		{\True  $J=0$}
		{$J=-2$} 
	\loigiai{Ta có $J=\lim\dfrac {(n+1)(2n+3)}{n^3-2}=\lim\dfrac {2n^2+5n+3}{n^3-2}=\lim\dfrac {\dfrac  {2}{n}+\dfrac {5}{n^2}+\dfrac {3}{n^3}}{1-\dfrac {2}{n^3}}=\dfrac {0+0+0}{1-0}=0$.
	 }
\end{ex}

% nguồn: Ok_Toan11_TracNghiem_Lan2
\begin{ex}%[1D3H1-2]%[ma:8f97a723]
	 $\lim\dfrac{1}{5n+2}$ bằng
	\choice 
		{$\dfrac{1}{5}$} 
		{\True $0$} 
		 {$\dfrac{1}{2}$}
		{$+\infty $} 
	 \loigiai{
	 $\lim\dfrac{1}{5n+2}=\lim\dfrac{1}{n}\cdot\left(\dfrac{1}{5+\dfrac{2}{n}}\right)=0\cdot\dfrac{1}{5}=0$.}
\end{ex}

% nguồn: Ok_Toan11_TracNghiem_Lan2
\begin{ex}%[1D3H1-2]%[ma:2999c929]
	Tính  $\lim\dfrac{3\sqrt{n+5}-2\sqrt{n}}{n^2+5}$.
	\choice 
		 {\True $0$}
		{$+\infty$} 
		{$-\infty$} 
		{$1$} 
	 \loigiai{
	Ta có  $\lim\dfrac{3\sqrt{n+5}-2\sqrt{n}}{n^2+5}=\lim\dfrac{3\sqrt{\dfrac{1}{n^3}+\dfrac{5}{n^4}}-2\sqrt{\dfrac{1}{n^3}}}{1+\dfrac{5}{n^2}}=0$}
\end{ex}

% nguồn: Ok_Toan11_TracNghiem_Lan2
\begin{ex}%[1D3H1-2]%[ma:6fcea606]
	$\lim\dfrac{2n^4-2n+2}{4n^4+2n+5}$  bằng
	 \choice
		{$\dfrac{2}{11}$} 
		{\True  $\dfrac{1}{2}$}
		{$+\infty $} 
		 {$0$}
	\loigiai{ 
	Ta có $\lim\dfrac{2n^4-2n+2}{4n^4+2n+5}=\lim\dfrac{2-\dfrac{2}{n^3}+\dfrac{2}{n^4}}{4+\dfrac{2}{n^3}+\dfrac{5}{n^4}}=\dfrac{2-0+0}{4+0+0}=\dfrac{1}{2}$.} 
\end{ex}

% nguồn: Ok_Toan11_TracNghiem_Lan2
\begin{ex}%[1D3H1-2]%[ma:5d9182ce]
	Tính $I=\lim\dfrac{2n-3}{2{n^2}+3n+1}$. 
	 \choice
		{$I=-\infty  $}
		{\True  $I=0$}
		{$I=+\infty  $}
		 {$I=1$}
	\loigiai{$I=\lim\dfrac {2n-3}{2n^2+3n+1}=\lim\dfrac {n^2\left( \dfrac {2}{n}-\dfrac  {3}{n^2} \right)}{n^2\left( 2+\dfrac {3}{n}+\dfrac {1}{n^2} \right)}=\lim\dfrac {\dfrac {2}{n}-\dfrac {3}{n^2}}{2+\dfrac {3}{n}+\dfrac {1}{n^2}}=\dfrac {0-0}{2+0+0}=0$.
	} 
\end{ex}

% nguồn: Ok_Toan11_TracNghiem_Lan2
\begin{ex}%[1D3H1-2]%[ma:ba640fcf]
	Với $n$ là số nguyên dương, đặt $S_n=\dfrac {1}{1\sqrt {2}+2\sqrt {1}}+\dfrac {1}{2\sqrt {3}+3\sqrt {2}}+\cdots +\dfrac {1}{n\sqrt {n+1}+( n+1)\sqrt {n}}$. Khi đó $\lim S_n$ bằng 
	\choice 
		{$\dfrac {1}{\sqrt  {2}+1}$}
		 {$\dfrac {1}{\sqrt {2}-1}$}
		{\True $1$} 
		{$\dfrac {1}{\sqrt  {2}+2}$}
	 \loigiai{
		Ta có $\dfrac {1}{n\sqrt {n+1}+(n+1)\sqrt {n}}=\dfrac {\sqrt {n+1}-\sqrt  {n}}{\sqrt {n}\sqrt {n+1}}=\dfrac {1}{\sqrt {n}}-\dfrac {1}{\sqrt {n+1}}$.\\
		Suy  ra
		\begin{align*}
			S_n&=\dfrac {1}{1\sqrt {2}+2\sqrt {1}}+\dfrac {1}{2\sqrt {3}+3\sqrt  {2}}+\cdots +\dfrac {1}{n\sqrt {n+1}+( n+1)\sqrt {n}}\\
			 &=\dfrac{1}{\sqrt{1}}-\dfrac{1}{\sqrt{2}}+\dfrac{1}{\sqrt{2}}-\dfrac{1}{\sqrt{3}}+\cdots \dfrac{1}{\sqrt{n}}-\dfrac{1}{\sqrt{n+1}}\\
			 &=1-\dfrac{1}{\sqrt{n+1}}.
		\end{align*}
		Suy ra $\lim S_n=1$. 
	 }
\end{ex}

% nguồn: Ok_Toan11_TracNghiem_Lan2
\begin{ex}%[1D3H1-2]%[ma:f938739f]
	Tính  giới hạn $I=\lim\dfrac {(n+1)(4n-2)}{n^3+2}$.
	 \choice
		{\True  $I=0$}
		{$I=2$} 
		{$I=1$} 
		 {$I=3$}
	 \loigiai{Ta có
		\begin{align*}
			I&=\lim\dfrac {\left( n+1 \right)\left( 4n-2 \right)}{n^3+2}\\ 
			&=\lim\dfrac {n^2\left(  1+\dfrac {1}{n} \right)\left( 4-\dfrac {2}{n} \right)}{n^3\left( 1+\dfrac {2}{n^3} \right)}\\
			 &=\lim\left[\dfrac {1}{n}\cdot \dfrac {\left( 1+\dfrac {1}{n} \right)\left( 4-\dfrac {2}{n} \right)}{1+\dfrac {2}{n^3}}\right]\\
			&=0\cdot \dfrac  {\left( 1+0 \right)\left( 4-0 \right)}{1+0}\\
			&=0. 
		\end{align*}
		 Vậy $I=0$.
	 }
\end{ex}

% nguồn: Ok_Toan11_TracNghiem_Lan2
\begin{ex}%[1D3H1-2]%[ma:0d962df2]
	Giả sử $\lim\dfrac{\sqrt{9n^2-4}-\sqrt[3]{n^3+n^2}}{5n+11}=\dfrac{a}{b}$ với $a,b\in\mathbb{N}$ và $\dfrac{a}{b}$ tối giản. Giá trị của $T=2a-3b$  bằng
	 \choice
		 {$-9$}
		{$-1$} 
		 {$21$}
		{\True  $-11$}
	 \loigiai{Ta có $\lim\dfrac{\sqrt{9n^2-4}-\sqrt[3]{n^3+n^2}}{5n+11}=\lim\dfrac{\sqrt{9-\dfrac{4}{n^2}}-\sqrt[3]{1+\dfrac{1}{n}}}{5+\dfrac{11}{n}}=\dfrac{2}{5}$.\\
		Suy ra $a=2,b=5$. Do đó $T=2(2)-3(5)=-11$. 
	} 
\end{ex}

% nguồn: Ok_Toan11_TracNghiem_Lan2
\begin{ex}%[1D3H1-2]%[ma:66b0a8f7]
	 $\lim\dfrac{1+3+5+\ldots+2n+1}{3n^2+4}$ bằng
	 \choice
		 {$\dfrac{2}{3}$}
		{$0$} 
		{\True $\dfrac{1}{3}$} 
		 {$+\infty$}
	\loigiai{ 
		 Ta có $1+3+5+\ldots+\left(2n+1\right)=\dfrac{\left(1+2n+1\right)\left(n+1\right)}{2}=\left(n+1\right)^2$. \\
		Suy  ra
		$\lim\dfrac{1+3+5+\ldots+\left(2n+1\right)}{3n^2+4}=\lim\dfrac{\left(n+1\right)^2}{3n^2+4}=\lim\dfrac{1+\dfrac{2}{n}+\dfrac{1}{n^2}}{3+\dfrac{4}{n^2}}=\dfrac{1+0+0}{3+0}=\dfrac{1}{3}$. 
	} 
\end{ex}

% nguồn: Ok_Toan11_TracNghiem_Lan2
\begin{ex}%[1D3H1-2]%[ma:fa61dfc9]
	Giá trị của  $\lim\dfrac {2n^2-3n+1}{n^2+2n}$ bằng
	 \choice
		{$3$} 
		 {\True $2$}
		 {$1$}
		{$-\dfrac  {3}{2}$}
	\loigiai{Ta có $\lim\dfrac {2n^2-3n+1}{n^2+2n}=\lim\dfrac {2-\dfrac {3}{n}+\dfrac {1}{n^2}}{1+\dfrac  {2}{n}}=\dfrac {2-0+0}{1+0}=2$.
	} 
\end{ex}

% nguồn: Ok_Toan11_TracNghiem_Lan2
\begin{ex}%[1D3H1-2]%[ma:71d26ca0]
	Giới hạn  $\lim\dfrac{1^2+2^2+3^2+4^2+\ldots+n^2}{n^3+2n+7}$ có giá trị bằng?
	\choice 
		 {$\dfrac{2}{3}$}
		{$\dfrac{1}{6}$} 
		 {$0$}
		{\True  $\dfrac{1}{3}$}
	\loigiai{ 
		Ta có  kết quả quen thuộc $1^2+2^2+3^2+\ldots+n^2$ $=\dfrac{n\left(n+1\right)\left(2n+1\right)}{6}$.
		Do  đó
		\begin{eqnarray*}
			\lim\dfrac{1^2+2^2+3^2+4^2+\ldots+n^2}{n^3+2n+7}&=&\lim\dfrac{n\left(n+1\right)\left(2n+1\right)}{6\left(n^3+2n+7\right)}\\ &=&\lim\dfrac{\left(1+\dfrac{1}{n} \right)\left(2+\dfrac{1}{n} \right)}{6\left(1+\dfrac{2}{n^2}+\dfrac{7}{n^3} \right)}\\&=&\dfrac{1\cdot2}{6}=\dfrac{1}{3}. 
		\end{eqnarray*}
	 }
\end{ex}

% nguồn: Ok_Toan11_TracNghiem_Lan2
\begin{ex}%[1D3H1-2]%[ma:e9c028de]
	Dãy số $(u_n)$ với $u_n=\dfrac{1+2+3+\ldots+n}{1\,011n^2+1\,012}$. Khi đó $\lim\left(u_n+1\right)$  bằng
	\choice 
		 {$\dfrac{2\,019}{2\,022}$}
		{\True $\dfrac{2\,023}{2\,022}$} 
		{$\dfrac{2\,022}{2\,023}$} 
		 {$\dfrac{2\,021}{2\,022}$}
	 \loigiai{
		Ta  có $1+2+3+\ldots+n=\dfrac{n\left(n+1\right)}{2}$.\\
		Nên  $u_n=\dfrac{1+2+3+\ldots+n}{1\,011n^2+1\,012}=\dfrac{\dfrac{n\left(n+1\right)}{2}}{1\,011n^2+1\,012}=\dfrac{n^2+n}{2\,022n^2+2\,024}$.\\
		Do  đó $u_n+1=\dfrac{2\,023n^2+n+2\,024}{2\,022n^2+2\,024}$.\\
		Suy ra $\lim\left(u_n+1\right)$  $=\lim\dfrac{2\,023n^2+n+2\,024}{2\,022n^2+2\,024}$ $=\lim\dfrac{2\,023+\dfrac{1}{n}+\dfrac{2\,024}{n^2}}{2\,022+\dfrac{2\,024}{n^2}}=\dfrac{2\,023+0+0}{2\,022+0}=\dfrac{2\,023}{2\,022}$.
	 }
\end{ex}

% nguồn: Ok_Toan11_TracNghiem_Lan2
\begin{ex}%[1D3H1-2]%[ma:15f2dd65]
	 $\lim\dfrac{1}{5n+3}$ bằng
	 \choice
		{\True $0$} 
		 {$\dfrac{1}{3}$}
		{$+\infty  $}
		{$\dfrac{1}{5}$} 
	 \loigiai{
	Ta  có $\lim\dfrac{1}{5n+3}=\lim\dfrac{\dfrac{1}{n}}{5+\dfrac{3}{n}}=0$ .}
\end{ex}

% nguồn: Ok_Toan11_TracNghiem_Lan2
\begin{ex}%[1D3H1-2]%[ma:488bc57f]
	$\lim\dfrac{\left(2n-1\right)^2\left(5+3n^3\right)}{\left(3n+4\right)^3\left(1-n\right)^2}$ bằng 
	\choice 
		 {$\dfrac{2}{3}$}
		{$\dfrac{2}{9}$} 
		{$\dfrac{4}{3}$} 
		{\True  $\dfrac{4}{9}$}
	 \loigiai{
		 Ta có\\
		 $\lim\dfrac{\left(2n-1\right)^2\left(5+3n^3\right)}{\left(3n+4\right)^3\left(1-n\right)^2}=\lim\dfrac{\dfrac{\left(2n-1\right)^2}{n^2}\cdot\dfrac{\left(5+3n^3\right)}{n^3}}{\dfrac{\left(3n+4\right)^3}{n^3}\cdot \dfrac{\left(1-n\right)^2}{n^2}}=\lim\dfrac{\left(2-\dfrac{1}{n} \right)^2.\left(\dfrac{5}{n^3}+3\right)}{\left(3+\dfrac{4}{n} \right)^3. \left(\dfrac{1}{n}-1\right)^2}=\dfrac{2^2\cdot3}{3^3}=\dfrac{4}{9}$.\\
		Vậy $\lim\dfrac{\left(2n-1\right)^2\left(5+3n^3\right)}{\left(3n+4\right)^3\left(1-n\right)^2}=\dfrac{4}{9}$. 
	} 
\end{ex}

% nguồn: Ok_Toan11_TracNghiem_Lan2
\begin{ex}%[1D3H1-2]%[ma:a6d3df6d]
	Biết $\lim\dfrac{2n^2+3n^3}{n^3+2}=\dfrac{a}{b}$,  với $\dfrac{a}{b};(a;b\in\mathbb{N})$ là phân số tối giản. Tính $a+b$.
	\choice 
		{$3$} 
		{\True  $4$}
		{$5$} 
		{$6$} 
	 \loigiai{Ta có $\lim\dfrac{2n^2+3n^3}{n^3+2}=\lim\dfrac{n^3\left(\dfrac{2}{n}+3\right)}{n^3\left(1+\dfrac{2}{n^3}\right)}=\lim\dfrac{\dfrac{2}{n}+3}{1+\dfrac{2}{n^3}}=3$.\\
		Auy ra $a=3$; $b=1$.\\ 
		Vậy  $a+b=4$.
	} 
\end{ex}

% nguồn: Ok_Toan11_TracNghiem_Lan2
\begin{ex}%[1D3H1-2]%[ma:85fa6d7c]
	Tính  $\lim\dfrac{\sqrt{n+5}-2n}{n-1}$.
	 \choice
		{$0$} 
		{$-1$} 
		 {$-\infty$}
		 {\True $-2$}
	\loigiai{Ta có  $\lim\dfrac{\sqrt{n+5}-2n}{n-1}=\lim\dfrac{\sqrt{\dfrac{1}{n}+\dfrac{5}{n^2}}-2}{1-\dfrac{1}{n}}=-2$.
	} 
\end{ex}

% nguồn: Ok_Toan11_TracNghiem_Lan2
\begin{ex}%[1D3H1-2]%[ma:e2974cf6]
	Kết quả đúng của $\lim\dfrac{-n^2+2n+1}{\sqrt{3n^4+2}}$ là 
	\choice 
		 {$\dfrac{1}{2}$}
		{$-\dfrac{2}{3}$} 
		{\True  $-\dfrac{\sqrt 3}{3}$}
		 {$-\dfrac{1}{2}$}
	 \loigiai{
		Ta có  $\lim\dfrac{-n^2+2n+1}{\sqrt{3n^4+2}}=\lim\dfrac{-1+\dfrac{2}{n}+\dfrac{1}{n^2}}{\sqrt{3+\dfrac{2}{n^4}}}=\dfrac{-1+0+0}{\sqrt{3+0}}=-\dfrac{\sqrt 3}{3}$.
	 }
\end{ex}

% nguồn: Ok_Toan11_TracNghiem_Lan2
\begin{ex}%[1D3H1-2]%[ma:17a66e64]
	Biết $\lim\dfrac{2n^3+n^2-4}{an^3+2}=\dfrac{1}{2}$ với $a \ne 0$ là tham số. Khi  đó $a-a^2$ bằng
	 \choice
		{\True $-12$} 
		 {$-2$}
		{$0$} 
		 {$-6$}
	 \loigiai{
		Ta có $\lim\dfrac{2n^3+n^2-4}{an^3+2}=\lim\dfrac{n^3\left(2+\dfrac{1}{n}-\dfrac{4}{n^3} \right)}{n^3\left(a+\dfrac{2}{n^3} \right)}=\lim\dfrac{2+\dfrac{1}{n}-\dfrac{4}{n^3}}{a+\dfrac{2}{n^3}}=\dfrac{2+0-0}{a+0}=\dfrac{2}{a}=\dfrac{1}{2}$.\\ 
		Suy ra $a=4$.  Khi đó $a-a^2=4-4^2=-12$.
	} 
\end{ex}

% nguồn: Ok_Toan11_TracNghiem_Lan2
\begin{ex}%[1D3H1-2]%[ma:22e6bf7f]
	$\lim\dfrac  {3n+5}{2n-4}$ bằng
	\choice 
		{\True $\dfrac {3}{2}$} 
		{$-\dfrac  {5}{4}$}
		{$3$} 
		 {$-4$}
	\loigiai{ 
		 Ta có $\lim\dfrac {3n+5}{2n-4}=\lim\dfrac {3+\dfrac {5}{n}}{2-\dfrac {4}{n}}=\dfrac {3+0}{2-0}=\dfrac {3}{2}$.
	 }
\end{ex}

% nguồn: Ok_Toan11_TracNghiem_Lan2
\begin{ex}%[1D3H1-2]%[ma:989f7929]
	Tính $\lim\dfrac{\sqrt[3]{8n^3+n^2-1}+4n+3}{3n-7}$. 
	 \choice
		{\True $2$} 
		{$\dfrac{2}{3}$} 
		 {$\dfrac{4}{3}$}
		{$4$} 
	 \loigiai{Ta có $\lim\dfrac{\sqrt[3]{8n^3+n^2-1}+4n+3}{3n-7}=\lim\dfrac{\sqrt[3]{8+\dfrac{1}{n}-\dfrac{1}{n^3}}+4+\dfrac{3}{n}}{3-\dfrac{7}{n}}=2$.
	} 
\end{ex}

% nguồn: Ok_Toan11_TracNghiem_Lan2
\begin{ex}%[1D3H1-2]%[ma:f7d9024f]
	Trong các  giới hạn sau, giới hạn nào bằng $-1$?
	\choice 
		 {$\lim\dfrac {n-2\,021}{n-2\,022}$}
		{$\lim\dfrac {n-2\,022}{2\,022n-1}$} 
		{$\lim\dfrac  {2n-2\,022}{2\,022n-1}$}
		{\True $\lim\dfrac {n^2-2\,022}{2\,022n-n^2}$} 
	\loigiai{Ta có $\lim\dfrac {n^2-2\,022}{2\,022n-n^2}=\lim\dfrac {1-\dfrac {2\,022}{n^2}}{\dfrac  {2\,022}{n}-1}=\dfrac {1-0}{0-1}=-1$.
	 }
\end{ex}

% nguồn: Ok_Toan11_TracNghiem_Lan2
\begin{ex}%[1D3H1-2]%[ma:2f98e4c9]
	Tính $\lim\dfrac{2^{n}+1}{2\cdot  2^{n}+3}$.
	 \choice
		 {$2$}
		{$0$} 
		{$1$} 
		{\True  $\dfrac{1}{2}$}
	\loigiai{ 
		Ta có  $\lim\dfrac{2^{n}+1}{2\cdot 2^{n}+3}=\lim\dfrac{1+\left(\dfrac{1}{2}\right)^{n}}{2+3\cdot\left(\dfrac{1}{2}\right)^{n}}=\dfrac{1+0}{2+0}=\dfrac{1}{2}$.
	} 
\end{ex}

% nguồn: Ok_Toan11_TracNghiem_Lan2
\begin{ex}%[1D3H1-2]%[ma:9398aac5]
	 Tính $\lim\dfrac{\sqrt[3]{n^6+4}-\sqrt[3]{27n^6-3n^2+1}}{n^2-2024}$.
	\choice 
		{$0$} 
		{$+\infty$} 
		 {$-3$}
		 {\True $-2$}
	\loigiai{Ta có $\lim\dfrac{\sqrt[3]{n^6+4}-\sqrt[3]{27n^6-3n^2+1}}{n^2-2024}=\lim\dfrac{\sqrt[3]{1+\dfrac{4}{n^6}}-\sqrt[3]{27-\dfrac{3}{n^4}+\dfrac{1}{n^6}}}{1-\dfrac{2024}{n^2}}=-2$. 
	} 
\end{ex}

% nguồn: Ok_Toan11_TracNghiem_Lan2
\begin{ex}%[1D3H1-2]%[ma:235dd96d]
	Cho dãy số $(u_n)$ xác định bởi $u_n=\dfrac{1}{n^2}+\dfrac{3}{n^2}+\ldots+\dfrac{2n-1}{n^2}$ với $n\in \mathbb{N}^{*}$ Giá  trị của $\lim u_n$ bằng
	 \choice
		 {$0$}
		{$+\infty$} 
		{$-\infty$} 
		{\True $1$} 
	\loigiai{ 
		Ta có  $1+3+\ldots+\left(2n-1\right)=n^2$. \\
		Suy ra 
		$\dfrac{1}{n^2}+\dfrac{3}{n^2}+\ldots+\dfrac{2n-1}{n^2}=\dfrac{1+3+\ldots+\left(2n-1\right)}{n^2}=\dfrac{n^2}{n^2}=1$.\\ 
		Vậy  $\lim u_n=1$.
	} 
\end{ex}

% nguồn: Ok_Toan11_TracNghiem_Lan2
\begin{ex}%[1D3H1-2]%[ma:510258fe]
	Cho dãy số $(u_n)$ với $u_n=\dfrac{1}{1\cdot 3}+\dfrac{1}{3\cdot 5}+\ldots+\dfrac{1}{\left(2n-1\right)\cdot\left(2n+1\right)}$. Tính $\lim  u_n$.
	 \choice
		{\True $\dfrac{1}{2}$} 
		 {$0$}
		{$1$} 
		{$\dfrac{1}{4}$} 
	 \loigiai{
		Ta có 
		\begin{eqnarray*}
			u_n&=&\dfrac{1}{1\cdot  3}+\dfrac{1}{3\cdot 5}+\ldots+\dfrac{1}{\left(2n-1\right)\cdot\left(2n+1\right)}\\
			 &=&\dfrac{1}{2} \left(\dfrac{1}{1}-\dfrac{1}{3}+\dfrac{1}{3}-\dfrac{1}{5}+\ldots+\dfrac{1}{2n-1}-\dfrac{1}{2n+1} \right)\\
			&=&\dfrac{1}{2} \left(\dfrac{1}{1}-\dfrac{1}{2n+1} \right)=\dfrac{n}{2n+1} 
		\end{eqnarray*}
		 Suy ra $\lim u_n=\lim\dfrac{n}{2n+1}=\dfrac{1}{2}$.
	 }
\end{ex}

% nguồn: Ok_Toan11_TracNghiem_Lan2
\begin{ex}%[1D3H1-2]%[ma:eceb4684]
	Tính $\lim\dfrac{\sqrt{n^2-4}-\sqrt[3]{8n^3+n^2}}{n}$. 
	 \choice
		 {$0$}
		{$-2$} 
		 {$-\infty$}
		{\True  $-1$}
	\loigiai{Ta có  $\lim\dfrac{\sqrt{n^2-4}-\sqrt[3]{8n^3+n^2}}{n}=\lim\dfrac{\sqrt{1-\dfrac{4}{n^2}}-\sqrt[3]{8+\dfrac{1}{n}}}{1}=-1$.
	} 
\end{ex}

% nguồn: Ok_Toan11_TracNghiem_Lan2
\begin{ex}%[1D3H1-2]%[ma:94e54805]
	Tính giới hạn $\lim\dfrac{3^{n+1}-4\cdot 5^n}{2\cdot 5^n+5\cdot  4^{n+1}}$.
	\choice 
		{$-\dfrac{1}{7}$} 
		{$2$} 
		{\True $-2$} 
		{$-\dfrac{1}{5}$} 
	\loigiai{Ta có $\lim\dfrac{3^{n+1}-4\cdot 5^n}{2\cdot 5^n+5\cdot 4^{n+1}}=\lim\dfrac{3\cdot 3^n-4\cdot 5^n}{2\cdot  5^n+20\cdot 4^n}=\lim\dfrac{3\cdot\left(\dfrac{3}{5}\right)^n-4}{2+20\cdot\left(\dfrac{4}{5}\right)^n}=-2$.
	 }
\end{ex}

% nguồn: Ok_Toan11_TracNghiem_Lan2
\begin{ex}%[1D3H1-2]%[ma:488df496]
	Dãy số $(u_n)$ với $u_n=\dfrac{\left(3n-1\right)\left(3-n\right)^2}{\left(4n-5\right)^3}$ có giới hạn  bằng phân số tối giản $\dfrac{a}{b}$. Tính $a\cdot b$.
	\choice 
		 {\True $192$}
		{$68$} 
		 {$32$}
		 {$128$}
	 \loigiai{
		 Ta có $\lim\dfrac{\left(3n-1\right)\left(3-n\right)^2}{\left(4n-5\right)^3}=\lim\dfrac{\left(3-\dfrac{1}{n} \right)\left(\dfrac{3}{n}-1\right)^2}{\left(4-\dfrac{5}{n} \right)^3}=\dfrac{\left(3-0 \right)\left(0-1\right)^2}{\left(4-0 \right)^3}=\dfrac{3}{64}=\dfrac{a}{b}$.\\ Do đó $a\cdot b=192$.
	 }
\end{ex}

% nguồn: Ok_Toan11_TracNghiem_Lan2
\begin{ex}%[1D3H1-2]%[ma:e5082af6]
	 Tìm $I=\lim\dfrac{7{n^2}-2{n^3}+1}{3{n^3}+2{n^2}+1}$.
	\choice 
		{$\dfrac  {7}{3}$}
		{\True  $-\dfrac {2}{3}$}
		 {$0$}
		{$1$} 
	 \loigiai{
		Ta có $I=\lim\dfrac {7n^2-2n^3+1}{3n^3+2n^2+1}=\lim\dfrac {\dfrac {7}{n}-2+\dfrac {1}{n^3}}{3+\dfrac {2}{n}+\dfrac {1}{n^3}}=\dfrac {0-2+0}{3+0+0}=-\dfrac {2}{3}$. 
	 }
\end{ex}

% nguồn: Ok_TracNghiemLop11_Lan1
\begin{ex}%[1D3H1-2]%[ma:6e8448bd]
	Tính $\lim  \left(\dfrac{n^3+n^2-n}{1-2n^2-4n^3}\right)$.
	\choice 
		 {$\dfrac{7}{4}$}
		{$\dfrac{1}{4}$} 
		{\True$-\dfrac{1}{4}$} 
		 {$-\dfrac{9}{4}$}
	 \loigiai{$\lim \left(\dfrac{n^3+n^2-n}{1-2n^2-4n^3}\right)=\lim \left(\dfrac{1+\frac1n-\frac1{n^2}}{\frac1{n^3}-2\cdot\frac1n-4}\right)=-\dfrac14$.}
\end{ex}

% nguồn: Ok_TracNghiemLop11_Lan1
\begin{ex}%[1D3H1-2]%[ma:6133d96c]
	Cho dãy số $\left({u_n} \right)$ với  ${u_n}=\dfrac{an+4}{5n+3}$
	trong đó $a$ là tham số thực. Để dãy số $\left({u_n} \right)$ có giới hạn bằng $2$, giá trị của $a$  là
	 \choice
		 {\True $a=10$}
		{$a=8$} 
		 {$a=6$}
		{$a=4$} 
	 \loigiai{
		Ta có $\lim {u_n}=\lim \dfrac{an+4}{5n+3}=\lim  \dfrac{a+\dfrac{4}{n}}{5+\dfrac{3}{n}}=\dfrac{a}{5}$. Khi đó
		 $\lim {u_n}=2\Leftrightarrow \dfrac{a}{5}=2\Leftrightarrow a=10$.
	} 
\end{ex}

% nguồn: Ok_TracNghiemLop11_Lan1
\begin{ex}%[1D3H1-2]%[ma:1f2a5708]
	Cho $\lim\limits_{n \to +\infty} \dfrac{an+1}{2n-4}=-5$ giá  trị của $a$ là
	\choice{\True  $-10$}
		 {$1$}
		{$0$} 
		 {$6$}
	\loigiai{ 
		Ta có $\lim\limits_{n \to +\infty} \dfrac{an+1}{2n-4}=-5\Leftrightarrow \lim\limits_{n \to +\infty} \dfrac{a+\dfrac{1}{n}}{2-\dfrac{4}{n}}=-5\Leftrightarrow \dfrac{a}{2}=-5\Leftrightarrow a=-10$. 
	 }
\end{ex}

% nguồn: Ok_TracNghiemLop11_Lan1
\begin{ex}%[1D3H1-2]%[ma:f0c1a02b]
	$\lim\limits_{n \to +\infty} \dfrac{\sqrt{n^2+n-1}+n}{2n+1}$ bằng  bao nhiêu?
	 \choice{$-\infty$}
		{$+\infty$} 
		{$0$} 
		{\True  $1$}
	 \loigiai{
		Ta  có
		\begin{eqnarray*}
			\lim\limits_{n \to +\infty} \dfrac{\sqrt{n^2+n-1}+n}{2n+1}&=&\lim\limits_{n  \to +\infty} \dfrac{n\left(\sqrt{1+\dfrac{1}{n}-\dfrac{1}{n^2}}+1\right)}{n\left(2+\dfrac{1}{n}\right)}\\
			 &=&\lim\limits_{n \to +\infty} \dfrac{\sqrt{1+\dfrac{1}{n}-\dfrac{1}{n^2}}+1}{2+\dfrac{1}{n}}\\
			&=&1. 
		\end{eqnarray*}
	} 
\end{ex}

% nguồn: Ok_TracNghiemLop11_Lan1
\begin{ex}%[1D3H1-2]%[ma:625ca517]
	Kết quả của giới hạn $\lim \dfrac{{2^{n+1}}+3n+10}{3{n^2}-n+2}$ 
	 là
	\choice 
		{\True  $+\infty$}
		{$\dfrac{2}{3}$} 
		 {$\dfrac{3}{2}$}
		{$-\infty$} 
	\loigiai{ 
		Ta  có
		\[{2^n}=\sum\limits_{k=0}^n{C_n^k}\Rightarrow {2^n}\ge C_n^3=\dfrac{n\left(n-1\right)\left(n-2\right)}{6}\sim \dfrac{{n^3}}{6}\Rightarrow \left\{\begin{aligned}
		& \dfrac{n}{{2^n}}\to 0\\  & \dfrac{{2^n}}{{n^2}}\to+\infty \\
	\end{aligned} \right..\]
	Khi đó: $\lim \dfrac{{2^{n+1}}+3n+10}{3{n^2}-n+2}=\lim \dfrac{{2^n}}{{n^2}}\cdot\dfrac{2+3\cdot\dfrac{n}{{2^n}}+10\cdot{{\left(\dfrac{1}{2} \right)}^n}}{3-\dfrac{1}{n}+\dfrac{2}{{n^2}}}=+\infty$,\\ 
	vì $\left\{\begin{aligned}
	&  \lim \dfrac{{2^n}}{{n^2}}=+\infty \\ & \lim \dfrac{2+3\cdot \dfrac{n}{{2^n}}+10\cdot{{\left(\dfrac{1}{2} \right)}^n}}{3-\dfrac{1}{n}+\dfrac{2}{{n^2}}}=\dfrac{2}{3} > 0\\
	\end{aligned} \right.$.
	} 
\end{ex}

% nguồn: Ok_TracNghiemLop11_Lan1
\begin{ex}%[1D3H1-2]%[ma:2ccd1c49]
	Tìm $\lim  \dfrac{2-5^{n+2}}{3^n+2\cdot5^n}$.
	 \choice
		{$-\dfrac{5}{2}$} 
		{$\dfrac{5}{2}$} 
		{\True  $-\dfrac{25}{2}$}
		 {$-\dfrac{1}{2}$}
	\loigiai{ 
		Ta có $\lim \dfrac{2-5^{n+2}}{3^n+2\cdot5^n}=\lim \dfrac{2-25\cdot5^n}{3^n+2\cdot5^n}=\lim  \dfrac{2{\left(\dfrac{1}{5}\right)}^n-25}{{\left(\dfrac{3}{5}\right)}^n+2}=\dfrac{-25}{2}$. \\
	 }
\end{ex}

% nguồn: Ok_TracNghiemLop11_Lan1
\begin{ex}%[1D3H1-2]%[ma:02fe1c20]
	Tính giới hạn $\lim \dfrac{n^3 - 2n + 1}{4n  - 5n^3}$.
	\choice 
		 {\True $-\dfrac{1}{5}$}
		{$\dfrac{1}{4}$} 
		{$+\infty$} 
		{$0$} 
	 \loigiai{
	$\lim \dfrac{n^3 - 2n + 1}{4n - 5n^3} = \lim \dfrac{n^3 \left( 1 - \dfrac{2}{n^2} + \dfrac{1}{n^3} \right)}{n^3 \left( \dfrac{4}{n^2} - 5 \right)} = \lim \dfrac{1  - \dfrac{2}{n^2} + \dfrac{1}{n^3}}{\dfrac{4}{n^2} - 5} = -\dfrac{1}{5}$.
	 }
\end{ex}

% nguồn: Ok_TracNghiemLop11_Lan1
\begin{ex}%[1D3H1-2]%[ma:1d62ba92]
	Dãy số $(u_n)$ với $u_n=\dfrac{\sqrt{2n^3+n}+3n-1}{\sqrt{6n^3+2n^2}+n}$ có giới hạn bằng $\sqrt{\dfrac{a}{b}}$, $a>0$, $b>0$ và $\text{ƯCLN}(a,b=1)$. Hãy tính giá  trị của $a^2+b^2$.
	 \choice
		{$5$} 
		 {$40$}
		 {$9$}
		 {\True $10$}
	\loigiai{ 
		Ta có $$\lim u_n=\lim  \dfrac{\sqrt{2n^3+n}+3n-1}{\sqrt{6n^3+2n^2}+n}=\lim \dfrac{\sqrt{2+\dfrac{1}{n^2}}+\dfrac{3}{\sqrt{n}}-\dfrac{1}{n\sqrt{n}}}{\sqrt{6+\dfrac{2}{n}}+\dfrac{1}{\sqrt{n}}}=\sqrt{\dfrac{1}{3}}.$$
	Suy ra $a=1$, $b=3 \Rightarrow  a^2+b^2=10$.}
\end{ex}

% nguồn: Ok_TracNghiemLop11_Lan1
\begin{ex}%[1D3H1-2]%[ma:9a78cc74]
	Cho dãy số $(u_n)$ có $u_n=n\cdot \sqrt{\dfrac{2n+2}{n^4+n^2-1}}$. Khi đó, $\lim \limits_{n\to +\infty  } u_n$ có giá trị là
	 \choice
		{$+\infty  $}
		 {$2$}
		{$-\infty $} 
		 {\True $0$}
	\loigiai{ 
		Ta  có
		\begin{align*}
			 \lim \limits_{n\to +\infty } u_n&=\lim \limits_{n\to +\infty } n\cdot \sqrt{\dfrac{2n+2}{n^4+n^2-1}}\\
			&=\lim \limits_{n\to +\infty } \sqrt{\dfrac{2n^3+2n^2}{n^4+n^2-1}}\\ 
			&=\lim  \limits_{n\to +\infty } \sqrt{\dfrac{\dfrac{2}{n}+\dfrac{2}{n^2}}{1+\dfrac{1}{n^2}-\dfrac{1}{n^4}}}\\
			 &=\sqrt{\dfrac{0+0}{1+0-0}}=0.\
		\end{align*}
	} 
\end{ex}

% nguồn: Ok_TracNghiemLop11_Lan1
\begin{ex}%[1D3H1-2]%[ma:3de6b941]
	 Tính $\lim \dfrac{4^{n}-5^n}{2^n+3\cdot 5^n}$.
	\choice{$\dfrac{4}{3}$} 
		 {$\dfrac{1}{3}$}
		{$1$} 
		{\True $-\dfrac{1}{3}$} 
	\loigiai{ 
		Ta có $\lim \dfrac{4^{n}-5^n}{2^n+3\cdot 5^n}=\lim \dfrac{\left(\frac{4}{5}\right)^n-1}{\left(\frac{2}{5}\right)^n+3}=\dfrac{-1}{3}.$ 
	 }
\end{ex}

% nguồn: Ok_Toan11_TracNghiem_Lan2
\begin{ex}%[1D3V1-2]%[ma:a81209e4]
	Tìm giới hạn $\lim\dfrac{1+2+2^2+\cdots+2^n}{1+3+3^2+\cdots+3^n}$. 
	 \choice
		 {\True $0$}
		 {$2$}
		{$1$} 
		 {$\dfrac{3}{2}$}
	 \loigiai{
		Ta có \allowdisplaybreaks\begin{eqnarray*}
		\lim\dfrac{1+2+2^2+\cdots+2^n}{1+3+3^2+\cdots+3^n}&=&\lim\dfrac{\dfrac{1-2^{n+1}}{-1}}{\dfrac{1-3^{n+1}}{-2}}\\ 
		&=&\lim\dfrac{\left(1-2^{n+1}\right)\cdot 2}{1-3^{n+1}}\\ 
		&=&\lim\dfrac{\left(\left(\dfrac{1}{3}\right)^{n+1}-\left(\dfrac{2}{3}\right)^{n+1}\right)\cdot  2}{\left(\dfrac{1}{3}\right)^{n+1}-1}=0.
	\end{eqnarray*}
 }
\end{ex}

% nguồn: Ok_Toan11_TracNghiem_Lan2
\begin{ex}%[1D3V1-2]%[ma:9ccf3d63]
	Cho dãy số $(u_n)$ xác định bởi $u_n=\dfrac{1}{n^2}+\dfrac{3}{n^2}+\cdots+\dfrac{2n-1}{n^2}$ với $n\in\mathbb{N}^{*}$. Giá trị  của $\lim u_n$ bằng
	\choice 
		 {$0$}
		 {$+\infty$}
		 {$-\infty$}
		 {\True $1$}
	 \loigiai{Ta có $1+3+\cdots+(2n-1)=n^2$.\\
		 Do đó $u_n=\dfrac{1}{n^2}+\dfrac{3}{n^2}+\cdots+\dfrac{2n-1}{n^2}=\dfrac{1+3+\cdots+(2n-1)}{n^2}=\dfrac{n^2}{n^2}=1$.\\
		Vậy $\lim  u_n=1$.
	 }
\end{ex}

% nguồn: Ok_Toan11_TracNghiem_Lan2
\begin{ex}%[1D3V1-2]%[ma:f8f673f8]
	Cho dãy số $\left(u_n \right)$ thỏa mãn $u_n=\sqrt{n+2018}-\sqrt{n+2017}, \forall n\in \mathbb{N}^{*}$. Khẳng định  nào sau đây sai?
	\choice 
		{\True Dãy số $\left(u_n \right)$ là dãy  tăng}
		 {$\lim\limits_{n\to+\infty} u_n=0$}
		{$0< u_n < \dfrac{1}{2\sqrt{2018}},  \forall n\in \mathbb{N}^{*}$}
		 {$\lim\limits_{n\to+\infty} \dfrac{u_{n+1}}{u_n}=1$}
	\loigiai{Ta có $u_n=\sqrt{n+2018}-\sqrt{n+2017}=\dfrac{1}{\sqrt{n+2018}+\sqrt{n+2017}}$.\\Suy ra $\dfrac{u_{n+1}}{u_n}=\dfrac{\sqrt{n+2018}+\sqrt{n+2017}}{\sqrt{n+2019}+\sqrt{n+2018}} < 1$ với mọi $n\in  \mathbb{N}^{*}$. Do đó, dãy số $\left(u_n \right)$ giảm.\\
		 Chú ý
		\begin{itemize}
			\item  $\lim\limits_{n\to+\infty} u_n=\lim\limits_{n\to+\infty} \dfrac{1}{\sqrt{n+2018}+\sqrt{n+2017}}=0$.
			\item $\lim\limits_{n\to+\infty} \dfrac{u_{n+1}}{u_n}=\lim\limits_{n\to+\infty}  \dfrac{\sqrt{n+2018}+\sqrt{n+2017}}{\sqrt{n+2019}+\sqrt{n+2018}}=1$.
			\item $0< u_n=\dfrac{1}{\sqrt{n+2018}+\sqrt{n+2017}} < \dfrac{1}{2\sqrt{n+2017}}  \le \dfrac{1}{2\sqrt{2018}}$.
		\end{itemize}
	 }
\end{ex}

% nguồn: Ok_Toan11_TracNghiem_Lan2
\begin{ex}%[1D3V1-2]%[ma:b413c6a1]
	Tính giới hạn $\lim\left[\left(1-\dfrac{1}{2^2} \right)\left(1-\dfrac{1}{3^2} \right)\ldots \left(1-\dfrac{1}{n^2} \right)\right]$. 
	\choice 
		{$1$} 
		{\True  $\dfrac{1}{2}$}
		 {$\dfrac{1}{4}$}
		{$\dfrac{3}{2}$} 
	\loigiai{ 
		Xét dãy số $(u_n)$, với $u_n=\left(1-\dfrac{1}{2^2} \right)\left(1-\dfrac{1}{3^2}  \right)\ldots \left(1-\dfrac{1}{n^2} \right)$, $n\ge 2$, $n\in \mathbb{N}$. Ta có
		\begin{align*}
			 &u_2=1-\dfrac{1}{2^2}=\dfrac{3}{4}=\dfrac{2+1}{2\cdot 2};\\
			&u_3=\left(1-\dfrac{1}{2^2} \right)\cdot\left(1-\dfrac{1}{3^2} \right)=\dfrac{3}{4}\cdot\dfrac{8}{9}=\dfrac{4}{6}=\dfrac{3+1}{2\cdot  3};\\
			 &u_4=\left(1-\dfrac{1}{2^2} \right)\cdot\left(1-\dfrac{1}{3^2} \right)\cdot\left(1-\dfrac{1}{4^2} \right)=\dfrac{3}{4}\cdot\dfrac{8}{9}\cdot\dfrac{15}{16}=\dfrac{5}{8}=\dfrac{4+1}{2\cdot4}\\
			&\ldots  \ldots\\
			 &u_n=\dfrac{n+1}{2n}.
		\end{align*}
		Dễ dàng chứng minh bằng phương pháp quy nạp để khẳng định $u_n=\dfrac{n+1}{2n},  \forall n\ge 2$.\\
		Khi đó $\lim\left[\left(1-\dfrac{1}{2^2} \right)\left(1-\dfrac{1}{3^2} \right)\ldots  \left(1-\dfrac{1}{n^2} \right)\right]=\lim\dfrac{n+1}{2n}=\dfrac{1}{2}$.
	} 
\end{ex}

% nguồn: Ok_Toan11_TracNghiem_Lan2
\begin{ex}%[1D3V1-2]%[ma:5ead5d49]
	Giới  hạn $\lim\left(\dfrac{1}{n^2}+\dfrac{2}{n^2}+\dfrac{3}{n^2}+\cdots+\dfrac{n}{n^2}\right)$ bằng
	 \choice
		 {$1$}
		{$0$} 
		 {$\dfrac{1}{3}$}
		{\True $\dfrac{1}{2}$} 
	 \loigiai{
		Ta có \allowdisplaybreaks\begin{eqnarray*}
		\lim\left(\dfrac{1}{n^2}+\dfrac{2}{n^2}+\dfrac{3}{n^2}+\cdots+\dfrac{n}{n^2}\right)&=&\lim\left(\dfrac{1+2+3+\cdots+n}{n^2}\right)\\ 
		 &=&\lim\left(\dfrac{n(n+1)}{2n^2}\right)\\
		&=&\lim\left(\dfrac{1}{2}+\dfrac{1}{2n}\right)=\dfrac{1}{2}. 
	\end{eqnarray*}
 }
\end{ex}

% nguồn: Ok_Toan11_TracNghiem_Lan2
\begin{ex}%[1D3V1-2]%[ma:82f48033]
	Cho dãy số $(u_n)\colon u_n=\dfrac{1}{n^2+n},n\in\mathbb{N}^{*}$ và  $(v_n)\colon v_n=\sum\limits_{k=1}^nu_n$. Tính $\lim v_n$
	 \choice
		 {$0$}
		 {\True $1$}
		 {$+\infty$}
		{$-\infty$} 
	\loigiai{ 
		Ta  có $u_n=\dfrac{1}{n^2+n}=\dfrac{1}{n(n+1)}=\dfrac{1}{n}-\dfrac{1}{n+1}$.\\
		Suy ra \allowdisplaybreaks\begin{eqnarray*}
		v_n&=&\sum\limits_{k=1}^nu_n=u_1+u_2+u_3+\cdots+u_n\\ 
		 &=&\left(1-\dfrac{1}{2}\right)+\left(\dfrac{1}{2}-\dfrac{1}{3}\right)+\left(\dfrac{1}{3}-\dfrac{1}{4}\right)+\cdots+\left(\dfrac{1}{n}-\dfrac{1}{n+1}\right)\\
		&=&1-\dfrac{1}{n+1}=\dfrac{n}{n+1}. 
	\end{eqnarray*}
	Vậy  $\lim v_n=\lim\dfrac{n}{n+1}=\lim\dfrac{1}{1+\dfrac{1}{n}}=1$.
 }
\end{ex}

% nguồn: Ok_Toan11_TracNghiem_Lan2
\begin{ex}%[1D3V1-2]%[ma:5366c458]
	Giá trị  $\lim\left(\dfrac{1}{n^2}+\dfrac{2}{n^2}+\cdots+\dfrac{n}{n^2}\right)$ bằng
	\choice 
		 {$1$}
		{\True  $\dfrac{1}{2}$}
		{$\dfrac{1}{3}$} 
		 {$0$}
	\loigiai{Ta có $1$; $2$; $3$;  \ldots; $n$ là một cấp số cộng với $u_1=1$; $d=1$ suy ra $1+2+\cdots+n=\dfrac{n(n+1)}{2}$.\\
		 Do đó $$\lim\left(\dfrac{1}{n^2}+\dfrac{2}{n^2}+\cdots+\dfrac{n}{n^2}\right)=\lim\left(\dfrac{\dfrac{n(n+1)}{2}}{n^2}\right)=\lim\dfrac{n+1}{2n}=\lim\left(\dfrac{1}{2}+\dfrac{1}{2n}\right)=\dfrac{1}{2}.$$
	 }
\end{ex}

% nguồn: Ok_Toan11_TracNghiem_Lan2
\begin{ex}%[1D3V1-2]%[ma:b6d83dc2]
	Tính giới hạn $\lim\left[\dfrac{1}{1\cdot 2}+\dfrac{1}{2\cdot  3}+\dfrac{1}{3\cdot 4}+\cdots+\dfrac{1}{n(n+1)}\right]$.
	\choice 
		 {$0$}
		{$2$} 
		{\True  $1$}
		{$\dfrac{3}{2}$} 
	\loigiai{Ta có 
		\allowdisplaybreaks\begin{eqnarray*}
		\dfrac{1}{1\cdot  2}+\dfrac{1}{2\cdot 3}+\dfrac{1}{3\cdot 4}+\cdots+\dfrac{1}{n(n+1)}&=&\dfrac{1}{1}-\dfrac{1}{2}+\dfrac{1}{2}-\dfrac{1}{3}+\cdots+\dfrac{1}{n-1}-\dfrac{1}{n}+\dfrac{1}{n}-\dfrac{1}{n+1}\\
		 &=&1-\dfrac{1}{n+1}.
	\end{eqnarray*}
	Vậy $\lim\left[\dfrac{1}{1\cdot  2}+\dfrac{1}{2\cdot 3}+\dfrac{1}{3\cdot 4}+\cdots+\dfrac{1}{n(n+1)}\right]=\lim\left(1-\dfrac{1}{n+1}\right)=1$.
 }
\end{ex}

% nguồn: Ok_Toan11_TracNghiem_Lan2
\begin{ex}%[1D3V1-2]%[ma:bda7d77e]
	Với $n$ là số tự nhiên lớn hơn $2$, đặt $S_n=\dfrac{1}{\mathrm{C}_3^3}+\dfrac{1}{\mathrm{C}_4^3}+\dfrac{1}{\mathrm{C}_5^4}+\cdots+\dfrac{1}{\mathrm{C}_n^3}$. Tính  $\lim S_n$.
	\choice 
		{$1$} 
		 {\True $\dfrac{3}{2}$}
		 {$3$}
		 {$\dfrac{1}{3}$}
	 \loigiai{Ta có $\mathrm{C}_n^3=\dfrac{n!}{3!(n-3)!}=\dfrac{(n-3)!(n-2)(n-1)n}{(n-3)!\times6}=\dfrac{n(n-1)(n-2)}{6}$.\\
		Suy  ra $\dfrac{1}{\mathrm{C}_n^3}=\dfrac{6}{n(n-1)(n-2)}$.\\
		 Như thế $S_n=\dfrac{6}{1.2.3}+\dfrac{6}{2.3.4}+\dfrac{6}{3.4.5}+\cdots+\dfrac{6}{n(n-1)(n-2)}$.\\
		Nhận  xét
		\begin{itemize}
			\item $\dfrac{2}{1\cdot 2\cdot 3}=\dfrac{1}{1\cdot  2}-\dfrac{1}{2\cdot3}$;
			\item  $\dfrac{2}{2\cdot 3\cdot 4}=\dfrac{1}{2\cdot 3}-\dfrac{1}{3\cdot 4}$;
			\item $\vdots$ 
			\item $\dfrac{2}{(n-2)(n-1)n}=\dfrac{1}{(n-2)(n-1)}-\dfrac{1}{(n-1)n}$. 
		\end{itemize}
		Do đó \allowdisplaybreaks\begin{eqnarray*}
		S_n&=&3\left(\dfrac{1}{1\cdot 2}-\dfrac{1}{2\cdot  3}+\dfrac{1}{2\cdot 3}-\dfrac{1}{3\cdot 4}+\cdots+\dfrac{1}{n-2}-\dfrac{1}{n-1}+\dfrac{1}{n-1}-\dfrac{1}{n}\right)\\
		 &=&3\left(\dfrac{1}{2}-\dfrac{1}{n}\right)=3\left(\dfrac{n-2}{2n}\right)=\dfrac{3n-6}{2n}.
	\end{eqnarray*}
	Vậy $\lim  S_n=\lim\left(\dfrac{3n-6}{2n}\right)=\lim\left(\dfrac{3-\dfrac{6}{n}}{2}\right)=\dfrac{3}{2}$.
} 
\end{ex}

% nguồn: Ok_Toan11_TracNghiem_Lan2
\begin{ex}%[1D3V1-2]%[ma:221320f4]
	Cho dãy số $\left(u_n \right)$ thoả mãn $\heva{&{u_1=1}  \\
	&{u_{n+1}=\dfrac{2}{3} u_n+4, \forall  n\in \mathbb{N}^{*}}
	}$.  Tìm $\lim u_n$
	\choice 
		 {$\lim u_n=1$}
		 {$\lim u_n=4$}
		{\True$\lim  u_n=12$}
		{$\lim u_n=3$} 
	\loigiai{Đặt $v_n=u_n-12, \forall n\in \mathbb{N}^{*}$. Khi đó $v_{n+1}=u_{n+1}-12=\dfrac{2}{3} u_n+4-12=\dfrac{2}{3} (u_n-12)=\dfrac{2}{3} v_n, \forall  n\in \mathbb{N}^{*}$.\\
		Suy ra dãy số $\left(v_n \right)$ là cấp số nhân với công bội $q=\dfrac{2}{3}$ và  số hạng đầu $v_1=-11$.\\
		Suy ra $v_n=-11\left(\dfrac{2}{3}  \right)^{n-1}, \forall n\in\mathbb{N}^{*}$.\\
	Từ đó $u_n=-11\left(\dfrac{2}{3} \right)^{n-1}+12, \forall n\in\mathbb{N}^{*}$. Vậy $\lim u_n=12$.} 
\end{ex}

% nguồn: Ok_Toan11_TracNghiem_Lan2
\begin{ex}%[1D3V1-2]%[ma:5455f73f]
	Cho cấp số cộng $\left(u_n \right)$ có số hạng đầu $u_1=2$ và công sai $d=3$. Tìm $\lim\dfrac{n}{u_n}$ 
	\choice 
		 {\True $L=\dfrac{1}{3}$}
		{$L=\dfrac{1}{2}$} 
		{$L=3$} 
		 {$L=2$}
	\loigiai{Ta có $u_n=u_1+\left(n-1\right)d=2+\left(n-1\right)3=3n-1$. Do  đó $\lim\dfrac{n}{u_n}=\lim\dfrac{n}{3n-1}=\lim\dfrac{1}{3-\dfrac{1}{n}}=\dfrac{1}{3}$.
	 }
\end{ex}

% nguồn: Ok_Toan11_TracNghiem_Lan2
\begin{ex}%[1D3V1-2]%[ma:e5cfd87c]
	Cho dãy số $(u_n)$ với $u_n=\dfrac{1}{1\cdot 3}+\dfrac{1}{3\cdot 5}+\cdots+\dfrac{1}{(2n-1)\cdot (2n+1)}$. Tính $\lim  u_n$.
	 \choice
		{\True $\dfrac{1}{2}$} 
		{$0$} 
		{$1$} 
		 {$\dfrac{1}{4}$}
	\loigiai{Ta có \allowdisplaybreaks\begin{eqnarray*}
		 u_n&=&\dfrac{1}{1\cdot 3}+\dfrac{1}{3\cdot 5}+\cdots+\dfrac{1}{(2n-1)\cdot (2n+1)}\\
		&=&\dfrac{1}{2}\left(\dfrac{1}{1}-\dfrac{1}{3}+\dfrac{1}{3}-\dfrac{1}{5}+\cdots+\dfrac{1}{2n-1}-\dfrac{1}{2n+1}\right)\\ 
		&=&\dfrac{1}{2}\left(\dfrac{1}{1}-\dfrac{1}{2n+1}\right)=\dfrac{n}{2n+1}. 
	\end{eqnarray*}
	 Vậy $\lim u_n=\lim\dfrac{n}{2n+1}=\dfrac{1}{2}$.
} 
\end{ex}

% nguồn: Ok_Toan11_TracNghiem_Lan2
\begin{ex}%[1D3V1-2]%[ma:cec787e3]
	Cho $L=\lim\dfrac{(a n+1)^3\left(n^2+3n+5\right)}{(2n-1)^2\left(n^2+3\right)^2}$. Tìm tất cả các giá trị của $a$ để  $L=\dfrac{1}{2}$.
	 \choice
		{$a=0$} 
		{\True  $a=1$}
		{$a \in \mathbb{R}$} 
		{$a  \in \varnothing$}
	 \loigiai{
		Ta có $\begin{aligned}
		L&=\lim\dfrac{(a  n+1)^3\left(n^2+3n+5\right)}{(2n-1)\left(n^2+3\right)^2}\\
		&=\lim\dfrac{n^3\left(a+\dfrac{1}{n}\right)^3\cdot n^2\left(1+\dfrac{3}{n}+\dfrac{5}{n^2}\right)}{n\left(2-\dfrac{1}{n}\right) \cdot  n^4\left(1+\dfrac{3}{n^2}\right)^2}\\
		&=\lim\dfrac{\left(a+\dfrac{1}{n}\right)^3\cdot\left(1+\dfrac{3}{n}+\dfrac{5}{n^2}\right)}{\left(2-\dfrac{1}{n}\right) \cdot\left(1+\dfrac{3}{n^2}\right)^2}\\ 
		&=\dfrac{a^3}{2}. 
	\end{aligned}$\\
	Từ đề bài  suy ra $\dfrac{a^3}{2}=\dfrac{1}{2} \Leftrightarrow a=1$.\\
	Vậy  $a=1$ là giá trị cần tìm.
} 
\end{ex}

% nguồn: Ok_Toan11_TracNghiem_Lan2
\begin{ex}%[1D3V1-2]%[ma:e2cd61fd]
	Biết  $\lim\dfrac{\sqrt{n^4-4}-\sqrt[3]{27n^6+5}+9n+1}{an^2+n+1}=-1$ với $a\neq0$. Mệnh đề nào dưới đây đúng?
	 \choice
		{$a\in(-7;-1)$} 
		{$a\in(-1;1)$} 
		{\True $a\in(1;7)$} 
		 {$a\in(7;12)$}
	 \loigiai{
		Ta có  $\lim\dfrac{\sqrt{n^4-4}-\sqrt[3]{27n^6+5}+9n+1}{an^2+n+1}=\lim\dfrac{\sqrt{1-\dfrac{4}{n^4}}-\sqrt[3]{27+\dfrac{5}{n^6}}+\dfrac{9}{n}+\dfrac{1}{n^2}}{a+\dfrac{1}{n}+\dfrac{1}{n^2}}=\dfrac{-2}{a}$.\\
	Suy ra $\dfrac{-2}{a}=-1\Rightarrow  a=2\in(1;7)$}
\end{ex}

% nguồn: Ok_Toan11_TracNghiem_Lan2
\begin{ex}%[1D3V1-2]%[ma:c82d6062]
	Cho  $a,b$ là các số thực thỏa mãn $\lim\dfrac{an^3+2bn^2+n+2023}{2n^2+3}=1$. Tổng $2a+b$ bằng
	\choice 
		{\True  $1$}
		 {$4$}
		{$2$} 
		 {$3$}
	\loigiai{Từ $\lim\dfrac{an^3+2bn^2+n+2023}{2n^2+3}=1\Rightarrow a=0$ (vì  nếu $a\neq0$ thì bậc cao nhất của tử lớn hơn bậc cao nhất của mẫu thì giới hạn là vô cực).\\
		Lúc  đó $\lim\dfrac{an^3+2bn^2+n+2023}{2n^2+3}=\lim\dfrac{2bn^2+n+2023}{2n^2+3}=\lim\dfrac{2b+\dfrac{1}{n}+\dfrac{2023}{n^2}}{2+\dfrac{3}{n^2}}=b=1$.\\
		Vậy  $2a+b=1$.
	} 
\end{ex}

% nguồn: Ok_Toan11_TracNghiem_Lan2
\begin{ex}%[1D3V1-2]%[ma:d52b3e72]
	Trong các dãy số $\left(u_n \right)$ cho dưới đây, dãy số nào có giới  hạn khác $1$?
	 \choice
		{$u_n=\dfrac{n\left(n-2018\right)^{2017}}{\left(n-2017\right)^{2018}}$} 
		{$u_n=n\left(\sqrt{n^2+2020}-\sqrt{4n^2+2017}  \right)$}
		 {\True $u_n=\dfrac{2}{1.3}+\dfrac{2}{3.5}+\cdots+\dfrac{2}{\left(2n+1\right)\left(2n+3\right)}$}
		{$\heva{&{u_1=2018} \\ 
		&{u_{n+1}=\dfrac{1}{2} \left(u_n+1\right),  n\ge 1}
		}$} 
	\loigiai{\begin{itemize}
		\item Với phương án: 
		\[u_n=\dfrac{n\left(n-2018\right)^{2017}}{\left(n-2017\right)^{2018}} \to \dfrac{n.n^{2017}}{n^{2018}}  \to 1.
		 \]
		\item Với phương  án:
		\[u_n=n\left(\sqrt{n^2+2020}-\sqrt{4n^2+2017}  \right)\to n\left(\sqrt{n^2}-\sqrt{4n^2} \right)\to n.\left(-n\right)\to-\infty.
		 \]
		\item Với phương án: 
		\[u_n=\left(1-\dfrac{1}{3} \right)+\left(\dfrac{1}{3}-\dfrac{1}{5}  \right)+\cdots+\left(\dfrac{1}{2n+1}-\dfrac{1}{2n+3} \right)=1-\dfrac{1}{2n+3} \to \dfrac{1}{2}.
		 \]
		\item Với phương án: 
		\[u_{n+1}=\dfrac{1}{2} \left(u_n+1\right)\Leftrightarrow  u_{n+1}-1=\dfrac{1}{2} \left(u_n-1\right).
		 \]
		Đặt $v_n=u_n-1$ ta có $\heva{&{v_1=2017}  \\
		&{v_{n+1}=\dfrac{1}{2}.v_n,  n\ge 1}
		}$. Suy ra dãy $\left(v_n \right)$ là một cấp số nhân có số hạng đầu bằng $2017$ công bội bằng $\dfrac{1}{2}$ nên 
		\[v_n=2017.\left(\dfrac{1}{2}  \right)^{n-1} \left(n\ge 1\right).
		 \]
		Suy ra $u_n=2017.\left(\dfrac{1}{2} \right)^{n-1}+1\left(n\ge 1\right)$ do đó  $\lim u_n=1$.
	\end{itemize}
	Ở phương án D, ta có thể chứng minh $u_n > 1$ với mọi $n\ge 1$ và $\left(u_n \right)$ là dãy giảm nên $\left(u_n \right)$ sẽ có giới hạn. Gọi $\lim u_n=a$. Khi đó từ $u_{n+1}=\dfrac{1}{2} \left(u_n+1\right), n\ge 1$  suy ra $a=\dfrac{1}{2} \left(a+1\right)\Leftrightarrow a=1$ do đó $\lim u_n=1$.
 }
\end{ex}

% nguồn: Ok_Toan11_TracNghiem_Lan2
\begin{ex}%[1D3V1-2]%[ma:7c8c1c54]
	Đặt $f(n)=\left(n^2+n+1\right)^2+1$, xét dãy số $\left(u_n  \right)$ sao cho $u_n=\dfrac{f(1).f(3).f(5)\cdots f\left(2n-1\right)}{f(2).f(4).f(6)\cdots f(2n)}$. Tìm $\lim n\sqrt{u_n}$
	 \choice
		 {$\lim n\sqrt{u_n}=\dfrac{1}{\sqrt{3}}$}
		{$\lim n\sqrt{u_n}=\sqrt{3}$} 
		{\True$\lim  n\sqrt{u_n}=\dfrac{1}{\sqrt{2}}$}
		{$\lim n\sqrt{u_n}=\sqrt{2}$} 
	\loigiai{Ta có  $f(n)=\left(n^2+n+1\right)^2+1=\left(n^2+1\right)\left[\left(n+1\right)^2+1\right]$.\\
		Do đó  $u_n=\dfrac{\left(1^2+1\right)\left(2^2+1\right)\left(3^2+1\right)\left(4^2+1\right)\cdots\left[\left(2n-1\right)^2+1\right]\left[4n^2+1\right]}{\left(2^2+1\right)\left(3^2+1\right)\left(4^2+1\right)\left(5^2+1\right)\cdots\left[4n^2+1\right]\left[\left(2n+1\right)^2+1\right]}$
		\[\Rightarrow u_n=\dfrac{2}{\left(2n+1\right)^2+1} \Rightarrow n\sqrt{u(n)}=\sqrt{\dfrac{2n^2}{\left(2n+1\right)^2+1}}. 
		 \]
		 \[\lim n\sqrt{u(n)}=\lim\sqrt{\dfrac{2n^2}{\left(2n+1\right)^2+1}}=\lim\sqrt{\dfrac{2}{\left(2+\dfrac{1}{n} \right)^2+\dfrac{1}{n^2}}}=\dfrac{1}{\sqrt{2}}.
		 \]
	} 
\end{ex}

% nguồn: Ok_Toan11_TracNghiem_Lan2
\begin{ex}%[1D3V1-2]%[ma:25ee8ad2]
	Có tất cả bao nhiêu giá trị  nguyên của tham số $a$ thuộc khoảng $(0;2\,019)$ để \[\lim\sqrt{\dfrac{9^{n}+3^{n+1}}{5^{n}+9^{n+a}}}\leq \dfrac{1}{2\,187}\]
	\choice 
		 {$2\,018$}
		 {\True $2\,012$}
		{$2\,019$} 
		{$2\,011$} 
	\loigiai{ 
		 Ta có $\lim\sqrt{\dfrac{9^{n}+3^{n+1}}{5^{n}+9^{n+a}}}=\lim\sqrt{\dfrac{1+3\left(\dfrac{1}{3}\right)^{n}}{\left(\dfrac{5}{9}\right)^{n}+9^{a}}}=\dfrac{1}{3^{a}}\leq \dfrac{1}{2\,187}\Leftrightarrow \dfrac{1}{3^{a}}\leq \dfrac{1}{3^{7}}\Leftrightarrow a\geq 7$.\\
		Do $a$  nguyên thuộc khoảng $(0;2\,019)$ nên $a\in\{7;8;\ldots;2\,018\}$.
	} 
\end{ex}

% nguồn: Ok_Toan11_TracNghiem_Lan2
\begin{ex}%[1D3V1-2]%[ma:2fa1c646]
	Biết rằng $\lim\dfrac{\sqrt[3]{an^3+5n^2-7}}{\sqrt{3n^2-n+2}}=b\sqrt{3}$. Tính giá trị của biểu thức $P=\dfrac{a}{b^3}$. 
	\choice 
		 {$P=\dfrac{1}{27}$}
		{$P=3$} 
		{\True $P=27$} 
		 {$P=\dfrac{1}{3}$}
	\loigiai{Ta có  $\lim\dfrac{\sqrt[3]{an^3+5n^2-7}}{\sqrt{3n^2-n+2}}=\lim\dfrac{\sqrt[3]{a+\dfrac{5}{n}-\dfrac{7}{n^3}}}{\sqrt{3-\dfrac{1}{n}+\dfrac{2}{n^2}}}=\dfrac{\sqrt[3]a}{\sqrt{3}}$.\\
		Theo đề bài  thì $\dfrac{\sqrt[3]a}{\sqrt{3}}=b\sqrt{3}\Rightarrow\sqrt[3]a=3b\Rightarrow a=27b^3$.\\
		 Vậy $P=27$.
	 }
\end{ex}

% nguồn: Ok_Toan11_TracNghiem_Lan2
\begin{ex}%[1D3V1-2]%[ma:720b18e7]
	Tìm  $\lim u_n$ biết $u_n=\dfrac{1}{2^2-1}+\dfrac{1}{3^2-1}+\cdots+\dfrac{1}{n^2-1}$.
	\choice 
		 {\True $\dfrac{3}{4}$}
		 {$\dfrac{3}{5}$}
		{$\dfrac{2}{3}$} 
		 {$\dfrac{4}{3}$}
	 \loigiai{
		Ta có \allowdisplaybreaks\begin{eqnarray*}
		u_n&=&\dfrac{1}{2^2-1}+\dfrac{1}{3^2-1}+\cdots+\dfrac{1}{n^2-1}\\ 
		&=&\dfrac{1}{1\cdot 3}+\dfrac{1}{2\cdot  4}+\dfrac{1}{3\cdot 5}+\cdots+\dfrac{1}{(n-1)(n+1)}\\
		 &=&\dfrac{1}{2}\left(\dfrac{1}{1}-\dfrac{1}{3}+\dfrac{1}{2}-\dfrac{1}{4}+\dfrac{1}{3}-\dfrac{1}{5}+\cdots+\dfrac{1}{n-1}-\dfrac{1}{n+1}\right)\\
		&=&\dfrac{1}{2}\left(\dfrac{1}{1}+\dfrac{1}{2}-\dfrac{1}{n+1}\right)\\ 
		&=&\dfrac{3}{4}-\dfrac{1}{2(n+1)}. 
	\end{eqnarray*}
Vậy $\lim  u_n=\lim\left[\dfrac{3}{4}-\dfrac{1}{2(n+1)}\right]=\dfrac{3}{4}$.}
\end{ex}

% nguồn: Ok_TracNghiemLop11_Lan1
\begin{ex}%[1D3V1-2]%[ma:c8fe7475]
	Tìm tất cả các giá trị thực  của tham số $b$ để biểu thức $A=\lim\limits\dfrac{9-b^2n^2}{11n^2+3}<0$.
	\choice 
		{$b\le  0$}
		{\True $b\ne  0$}
		{$b<0$} 
		 {$b>0$}
	\loigiai{Ta  có $A=\lim\limits\dfrac{9-b^2n^2}{11n^2+3}=\lim\limits\dfrac{n^2\left(\dfrac{9}{n^2}-b^2\right)}{n^2\left(11+\dfrac{3}{n^2}\right)}=\lim\limits\dfrac{\dfrac{9}{n^2}-b^2}{11+\dfrac{3}{n^2}}=-\dfrac{b^2}{11}$.\\
	Yêu cầu bài toán xảy ra  khi $-\dfrac{b^2}{11}<0\Leftrightarrow b^2>0\Leftrightarrow b\ne0$.}
\end{ex}

% nguồn: Ok_Toan11_TracNghiem_Lan2
\begin{ex}%[1D3C1-2]%[ma:10979c63]
	Đặt $f(n)=\left(n^2+n+1\right)^2+1$, xét dãy số $(u_n)$ sao cho $u_n=\dfrac{f(1)\cdot f(3)\cdot f(5)\ldots f(2n-1)}{f(2)\cdot f(4)\cdot f(6)\ldots  f(2n)}$. Tìm $\lim n\sqrt{u_n}$.
	 \choice
		 {$\lim n\sqrt{u_n}=\dfrac{1}{\sqrt{3}}$}
		 {$\lim n\sqrt{u_n}=\sqrt{3}$}
		{\True  $\lim n\sqrt{u_n}=\dfrac{1}{\sqrt{2}}$}
		{$\lim n\sqrt{u_n}=\sqrt{2}$} 
	\loigiai{ 
		Ta có $f(n)=\left(n^2+n+1\right)^2+1=\left(n^2+1\right)\left[(n+1)^2+1\right]$.\\ 
		Do đó $u_n=\dfrac{\left(1^2+1\right)\left(2^2+1\right)\left(3^2+1\right)\left(4^2+1\right)\ldots\left[(2n-1)^2+1\right]\left[4n^2+1\right]}{\left(2^2+1\right)\left(3^2+1\right)\left(4^2+1\right)\left(5^2+1\right)\ldots\left[4n^2+1\right]\left[(2n+1)^2+1\right]}=\dfrac{2}{(2n+1)^2+1}$.\\ 
		Suy ra $n\sqrt{u(n)}=\sqrt{\dfrac{2n^2}{(2n+1)^2+1}}$.\\ 
		Vậy 
	 $\lim n\sqrt{u(n)}=\lim\sqrt{\dfrac{2n^2}{(2n+1)^2+1}}=\lim\sqrt{\dfrac{2}{\left(2+\dfrac{1}{n}\right)^2+\dfrac{1}{n^2}}}=\dfrac{1}{\sqrt{2}}$.}
\end{ex}

% nguồn: Ok_Toan11_TracNghiem_Lan2
\begin{ex}%[1D3C1-2]%[ma:8181de9b]
	Có bao nhiêu giá trị nguyên của  tham số $a$ thuộc khoảng $\left(0;2018\right)$ để có $\lim\sqrt{\dfrac{9^n+3^{n+1}}{5^n+9^{n+a}}} \le \dfrac{1}{2187}$?
	\choice 
		 {\True $2011$}
		{$2016$} 
		{$2019$} 
		{$2009$} 
	\loigiai{Do $\dfrac{9^n+3^{n+1}}{5^n+9^{n+a}}  > 0$ với $\forall n$ nên $\lim\sqrt{\dfrac{9^n+3^{n+1}}{5^n+9^{n+a}}}=\sqrt{\lim\dfrac{9^n+3^{n+1}}{5^n+9^{n+a}}}=\sqrt{\lim\dfrac{1+3.\left(\dfrac{1}{3} \right)^n}{\left(\dfrac{5}{9} \right)^n+9^a}}=\sqrt{\dfrac{1}{9^a}}=\dfrac{1}{3^a}$.\\
		Theo đề bài ta  có $\lim\sqrt{\dfrac{9^n+3^{n+1}}{5^n+9^{n+a}}} \le \dfrac{1}{2187}
		\Leftrightarrow \dfrac{1}{3^a} \le \dfrac{1}{2187} 
		\Leftrightarrow a\ge 7$.\\ 
		Do $a$  là số nguyên thuộc khoảng $\left(0;2018\right)$ nên có $a\in \left\{7;8;9;\ldots;2017\right\}
		\Rightarrow$ có $2011$ giá trị  của $a$.
	} 
\end{ex}

% nguồn: Ok_Toan11_TracNghiem_Lan2
\begin{ex}%[1D3C1-2]%[ma:e3e634f4]
	Với $n$ là  số nguyên dương, đặt $$S_n=\dfrac{1}{1\sqrt{2}+2\sqrt{1}}+\dfrac{1}{2\sqrt{3}+3\sqrt{2}}+\cdots+\dfrac{1}{n\sqrt{n+1}+(n+1)\sqrt{n}}.$$ Khi đó $\lim S_n$ bằng
	 \choice
		{$\dfrac{1}{\sqrt{2}+1}$} 
		 {$\dfrac{1}{\sqrt{2}-1}$}
		{\True $1$} 
		{$\dfrac{1}{\sqrt{2}+2}$} 
	\loigiai{ 
		Ta có \allowdisplaybreaks\begin{eqnarray*}
		\dfrac{1}{n\sqrt{n+1}+(n+1)\sqrt{n}}&=&\dfrac{1}{\sqrt{n}\sqrt{n+1}(\sqrt{n+1}+\sqrt{n})}\\ 
		&=&\dfrac{\sqrt{n+1}-\sqrt{n}}{\sqrt{n}\sqrt{n+1}}\\ 
		 &=&\dfrac{1}{\sqrt{n}}-\dfrac{1}{\sqrt{n+1}}.
	\end{eqnarray*}
	Suy ra \allowdisplaybreaks\begin{eqnarray*}
	 S_n&=&\dfrac{1}{1\sqrt{2}+2\sqrt{1}}+\dfrac{1}{2\sqrt{3}+3\sqrt{2}}+\cdots+\dfrac{1}{n\sqrt{n+1}+(n+1)\sqrt{n}}\\
	 &=&\dfrac{1}{1}-\dfrac{1}{\sqrt{2}}+\dfrac{1}{\sqrt{2}}-\dfrac{1}{\sqrt{3}}+\cdots\dfrac{1}{\sqrt{n}}-\dfrac{1}{\sqrt{n+1}}\\
	 &=&1-\dfrac{1}{\sqrt{n+1}}.
	\end{eqnarray*}
 Vậy $\lim S_n=1$.}
\end{ex}

% nguồn: Ok_Toan11_TracNghiem_Lan2
\begin{ex}%[1D3C1-2]%[ma:8df61f61]
	Cho dãy số $\left(u_n \right)$ được xác  định như sau $u_1=2016;u_{n-1}=n^2\left(u_{n-1}-u_n \right)$ với mọi $n\in \mathbb{N}^{*}, n\ge 2$ tìm giới hạn của dãy số $\left(u_n \right)$
	\choice 
		 {$1011$}
		{$1010$} 
		{\True$1008$} 
		 {$1009$}
	 \loigiai{Ta có $u_{n-1}=n^2\left(u_{n-1}-u_n \right)
		 \Leftrightarrow u_{n-1} \left(n^2-1\right)=n^2u_n
		\Leftrightarrow u_n=\dfrac{n-1}{n}.\dfrac{n+1}{n}.u_{n-1}$. Khi đó ta  có:
		 \[u_2=\dfrac{1}{2}\cdot\dfrac{3}{2}.u_1\]
		\[u_3=\dfrac{2}{3}\cdot  \dfrac{4}{3}.u_2\]
		 {\dots}
		\[u_n=\dfrac{n-1}{n}\cdot \dfrac{n+1}{n}.u_{n-1} 
		\] 
		Nhân theo vế  các đẳng thức trên ta có $u_n=\dfrac{n+1}{2n}\cdot u_1=\dfrac{n+1}{n}.1008$. Vậy $\lim u_n=1008$.
	} 
\end{ex}

% nguồn: Ok_Toan11_TracNghiem_Lan2
\begin{ex}%[1D3C1-2]%[ma:9c7809b1]
	Với $n$ là số tự nhiên lớn hơn  $2$ đặt $S_n=\dfrac{1}{C_3^3}+\dfrac{1}{C_4^3}+\dfrac{1}{C_5^4}+\cdots+\dfrac{1}{C_n^3}$. Tính $\lim S_n$
	\choice 
		{$1$} 
		{\True $\dfrac{3}{2}$} 
		 {$3$}
		 {$\dfrac{1}{3}$}
	\loigiai{Ta  có $C_n^3=\dfrac{n!}{3!\left(n-3\right)!}=\dfrac{\left(n-3\right)!\left(n-2\right)\left(n-1\right)n}{\left(n-3\right)!\times 6}=\dfrac{n\left(n-1\right)\left(n-2\right)}{6}
		 \Rightarrow \dfrac{1}{C_n^3}=\dfrac{6}{n\left(n-1\right)\left(n-2\right)}$.\\
		 Vậy ta có $S_n=\dfrac{6}{1\cdot 2\cdot 3}+\dfrac{6}{2\cdot 3\cdot 4}+\dfrac{6}{3\cdot 4\cdot 5}+\cdots+\dfrac{6}{n\left(n-1\right)\left(n-2\right)}$.\\
		 Nhận xét $\dfrac{2}{1\cdot 2\cdot 3}=\dfrac{1}{1\cdot 2}-\dfrac{1}{2\cdot 3}$; $\dfrac{2}{2\cdot 3\cdot 4}=\dfrac{1}{2\cdot 3}-\dfrac{1}{3\cdot 4}$ $\ldots$ $\dfrac{2}{\left(n-2\right)\left(n-1\right)n}=\dfrac{1}{\left(n-2\right)\left(n-1\right)}-\dfrac{1}{\left(n-1\right)n}$.
		\[\Rightarrow S_n=3\left(\dfrac{1}{1.2}-\dfrac{1}{2.3}+\dfrac{1}{2.3}-\dfrac{1}{3.4}+\cdots+\dfrac{1}{n-2}-\dfrac{1}{n-1}+\dfrac{1}{n-1}-\dfrac{1}{n}  \right)=3\left(\dfrac{1}{2}-\dfrac{1}{n} \right)=3\left(\dfrac{n-2}{2n} \right)=\dfrac{3n-6}{2n}\cdot
		\] 
		Vậy $\lim S_n=\lim\left(\dfrac{3n-6}{2n}  \right)=\lim\left(\dfrac{3-\dfrac{6}{n}}{2} \right)=\dfrac{3}{2}$.
	} 
\end{ex}

% nguồn: Ok_Toan11_TracNghiem_Lan2
\begin{ex}%[1D3C1-2]%[ma:413bb072]
	Dãy số $\left(u_n  \right)$ nào sau đây có giới hạn khác số $1$ khi $n$ dần đến vô cùng?
	\choice 
		{\True  $u_n=\dfrac{\left(2017-n\right)^{2018}}{n\left(2018-n\right)^{2017}}$}
		{$u_n=n\left(\sqrt{n^2+2018}-\sqrt{n^2+2016}  \right)$}
		{$\heva{&{u_1=2017} \\ 
		&{u_{n+1}=\dfrac{1}{2} \left(u_n+1\right),  n=1, 2, 3\ldots}
		}$} 
		 {$u_n=\dfrac{1}{1.2}+\dfrac{1}{2.3}+\dfrac{1}{3.4}+\cdots+\dfrac{1}{n\left(n+1\right)}$}
	\loigiai{Ta tính giới hạn của  các dãy số trong từng đáp án
		\begin{itemize}
			\item $\lim u_n=\lim\dfrac{\left(2017-n\right)^{2018}}{n\left(2018-n\right)^{2017}}=\lim\left[\dfrac{2017-n}{n}.\left(\dfrac{2017-n}{2018-n} \right)^{2017} \right]$ 
			\[=\lim\left[\left(\dfrac{2017}{n}-1\right)\left(\dfrac{\dfrac{2017}{n}-1}{\dfrac{2018}{n}-1} \right)^{2017} \right]=-1. 
			\] 
			\item $\lim u_n=\lim n\left(\sqrt{n^2+2018}-\sqrt{n^2+2016}  \right)=\lim\dfrac{n\left(n^2+2018-n^2-2016\right)}{\sqrt{n^2+2018}+\sqrt{n^2+2016}}$
			\[=\lim\dfrac{2n}{\sqrt{n^2+2018}+\sqrt{n^2+2016}}=\lim\dfrac{2}{\sqrt{1+\dfrac{2018}{n^2}}+\sqrt{1+\dfrac{2016}{n^2}}}=1. 
			 \]
			\item \textbf{Cách  1:} Ta có $u_{n+1}-1=\dfrac{1}{2} \left(u_n-1\right)
			\Rightarrow u_n-1=\dfrac{1}{2}  \left(u_{n-1}-1\right)=\ldots=\dfrac{1}{2^{n-1}} \left(u_1-1\right)$
			 \[\Rightarrow u_n=\dfrac{2016}{2^{n-1}}+1\Leftrightarrow u_n=4032.\left(\dfrac{1}{2} \right)^n+1\Rightarrow \lim u_n=1.
			 \]
			\item \textbf{Cách  2:}\\
			Bước 1: Ta chứng minh $\left(u_n \right)$ giảm và bị chặn dưới bởi $1$. Thật vậy bằng quy nạp ta có $u_1=2017>  1$.\\
			Giả sử  $u_n > 1\Rightarrow u_{n+1}=\dfrac{1}{2} \left(u_n+1\right) > \dfrac{1}{2} \left(1+1\right)=1$.\\
			Vậy $u_n > 1\forall  n\in \mathbb{N}^{*}$. Hơn nữa $u_{n+1}-u_n=\dfrac{1}{2} \left(1-u_n \right) < 0$ nên $\left(u_n \right)$ là dãy giảm.
			Suy ra $\left(u_n \right)$ có  giới hạn $\lim u_n=a$.\\
			Bước 2: Ta  có ${\mathop{a=\lim}\limits_{}} u_n=\lim u_{n+1}=\lim\dfrac{1}{2} \left(u_n+1\right)=\dfrac{1}{2} \lim u_n+\dfrac{1}{2}=\dfrac{1}{2} a+\dfrac{1}{2}$
			 \[\Rightarrow a=1.
			 \]
			 \item Ta có $u_n=\dfrac{1}{1.2}+\dfrac{1}{2.3}+\dfrac{1}{3.4}+\cdots+\dfrac{1}{n\left(n+1\right)}=1-\dfrac{1}{2}+\dfrac{1}{2}-\dfrac{1}{3}+\cdots+\dfrac{1}{n}-\dfrac{1}{n+1}=1-\dfrac{1}{n+1}=\dfrac{n}{n+1}$.
			\[\Rightarrow  \lim u_n=\lim\dfrac{n}{n+1}=1.
			 \]
		\end{itemize}
	} 
\end{ex}

% nguồn: Ok_Toan11_TracNghiem_Lan2
\begin{ex}%[1D3C1-2]%[ma:55c28e81]
	Cho dãy số  $\left(u_n \right)$ thỏa mãn $\heva{&{u_1=2} \\
	&{3\sqrt{4u_{n+1}+1}=\sqrt{4u_n+1}+4, \left(n\in  \mathbb{N}^{*} \right)}
	 }$. Tính $\lim u_n$
	\choice 
		 {$\dfrac{1}{3}$}
		 {\True $\dfrac{3}{4}$}
		{$\dfrac{1}{2}$} 
		{$\dfrac{2}{3}$} 
	 \loigiai{
		Chứng minh $\left(u_n  \right)$ là dãy giảm, tức là chứng minh: $u_{n+1} \le u_n, \forall n\in \mathbb{N}^{*}$.\\
		Với $n=1$ ta có: $3\sqrt{4u_2+1}=\sqrt{4u_1+1}+4\Leftrightarrow u_2=\dfrac{10}{9} \le u_1$-Giả sử mệnh đề đúng với $n=k$ tức là: $u_{k+1} \le u_k,  \forall n\in \mathbb{N}^{*}$.\
		Ta cần chứng minh mệnh đề đúng với $n=k+1$ tức là chứng minh: $u_{k+2} \le u_{k+1}$. Ta có 
		\[3\sqrt{4u_{k+2}+1}=\sqrt{4u_{k+1}+1}+4\le \sqrt{4u_k+1}+4=33\sqrt{4u_{k+1}+1}  \Leftrightarrow u_{k+2} \le u_{k+1}.
		 \]
		Vậy theo nguyên lý quy nạp suy ra $u_{n+1} \le u_n, \forall n\in \mathbb{N}^{*}$ tức $\left(u_n \right)$ là dãy  giảm.\\
		Tương tự, dùng quy  nạp ta dễ dàng chứng minh được $\dfrac{3}{4} < u_n \le 2$ tức dãy $\left(u_n \right)$ bị chặn. Từ đó suy ra dãy số có giới hạn.\\
		Đặt $x=\lim u_n$. Khi $n\to+\infty$ thì  $u_{n+1} \to x$ và
		\[3\sqrt{4x+1}=\sqrt{4x+1}+4\Leftrightarrow 36x+9=4x+1+16+8\sqrt{4x+1} \Leftrightarrow \sqrt{4x+1}=4x-1\Leftrightarrow  x=\dfrac{3}{4}.
		\] 
		Vậy $\lim  u_n=\dfrac{3}{4}$.
	 }
\end{ex}

% nguồn: Ok_Toan11_TracNghiem_Lan2
\begin{ex}%[1D3C1-2]%[ma:16198b15]
	Cho dãy số $(u_n)$ thỏa mãn: $u_1=1$; $u_{n+1}=\sqrt{\dfrac{2}{3}  u_n^2+a}$, $\forall n\in \mathbb{N}^{*}$. Biết rằng $\lim\left(u_1^2+u_2^2+\cdots+u_n^2-2n\right)=b$. Giá trị của biểu thức $T=ab$ là
	 \choice
		{\True  $-2$}
		 {$-1$}
		{$1$} 
		 {$2$}
	\loigiai{Ta có $\forall n\in \mathbb{N}^{*}$ 
		\[u_{n+1}=\sqrt{\dfrac{2}{3} u_n^2+a}  \Rightarrow u_{n+1}^2-3a=\dfrac{2}{3} \left(u_n^2-3a\right).
		\] 
		 Đặt $v_n=u_n^2-3a$ thì $\left(v_n \right)$ là cấp số nhân với $v_1=1-3a$ và công bội $q=\dfrac{2}{3}$.\\
		Do đó $v_n=\left(\dfrac{2}{3} \right)^{n-1}  \left(1-3a\right)\Rightarrow u_n^2=v_n+3a=\left(\dfrac{2}{3} \right)^{n-1} \left(1-3a\right)+3a$.\\
		Suy ra $u_1^2+u_2^2+\cdots+u_n^2-2n=\left(1-3a\right)\dfrac{1-\left(\dfrac{2}{3}  \right)^n}{1-\dfrac{2}{3}}-2n+3na=3\left(1-3a\right)\left(1-\left(\dfrac{2}{3} \right)^n \right)-n\left(3a-2\right)$.\\
		Vì  $\lim\left(u_1^2+u_2^2+\cdots+u_n^2-2n\right)=b$\\ nên $\lim\left(3\left(1-3a\right)\left(1-\left(\dfrac{2}{3} \right)^n \right)-n\left(3a-2\right)\right)=b\Leftrightarrow \heva{&{3a-2=0} \\
		&{b=3\left(1-3a\right)} 
		}  \Leftrightarrow \heva{&{a=\dfrac{2}{3}} \\
		&{b=-3}. 
		 }$\\
		suy ra $T=ab=-2$. 
	 }
\end{ex}

% nguồn: Ok_Toan11_TracNghiem_Lan2
\begin{ex}%[1D3C1-2]%[ma:f5ab0bf4]
	Tìm  $L=\lim\left(\dfrac{1}{1}+\dfrac{1}{1+2}+\cdots+\dfrac{1}{1+2+\cdots+n}\right)$.
	 \choice
		{$L=\dfrac{5}{2}$} 
		{$L=+\infty$} 
		{\True  $L=2$}
		{$L=\dfrac{3}{2}$} 
	\loigiai{Ta có $1+2+3+\cdots+k$ là tổng của cấp số cộng có $u_1=1,d=1$ nên $$1+2+3+\cdots+k=\dfrac{(1+k)k}{2}.$$ Khi đó $\dfrac{1}{1+2+\cdots+k}=\dfrac{2}{k(k+1)}=\dfrac{2}{k}-\dfrac{2}{k+1}$,  $\forall k\in\mathbb{N}^{*}$.\\
		 Vậy
		$L=\lim\left(\dfrac{2}{1}-\dfrac{2}{2}+\dfrac{2}{2}-\dfrac{2}{3}+\dfrac{2}{3}-\dfrac{2}{4}+\cdots+\dfrac{2}{n}-\dfrac{2}{n+1}\right)=\lim\left(\dfrac{2}{1}-\dfrac{2}{n+1}\right)=2$. 
	} 
\end{ex}

% nguồn: Ok_Toan11_TracNghiem_Lan2
\begin{ex}%[1D3C1-2]%[ma:1a5242ba]
	Cho hai dãy số $\left(u_n \right), \left(v_n \right)$ đều tồn tại giới hạn hữu hạn. Biết rằng hai dãy số đồng thời thỏa mãn các hệ thức $u_{n+1}=4v_n-2,  v_{n+1}=u_n+1$ với mọi $\forall n\in {\rm Z}^{+}$. Giá trị của giới hạn $\lim\limits_{n\to+\infty} \left(u_n+2v_n \right)$ bằng
	\choice 
		 {\True $0$}
		 {$\dfrac{3}{2}$}
		 {$-1$}
		{$\dfrac{1}{2}$} 
	\loigiai{Giả sử $\heva{&{\lim  u_n=a} \\
		&{\lim v_n=b} 
		}$ ta có $\heva{&{\lim u_{n+1}=\lim\left(4v_n-2\right)} \\ 
		&{\lim  v_{n+1}=\lim\left(u_n+1\right)}
		} 
		\Rightarrow \heva{&{a=4b-2}  \\
		&{b=a+1} 
		 } \Rightarrow \heva{&{a=-\dfrac{2}{3}} \\
		&{b=\dfrac{1}{3}} 
		}$.\\ 
		Vậy $\lim\limits_{n\to+\infty} \left(u_n+2v_n \right)=a+2b=-\dfrac{2}{3}+2.\dfrac{1}{3}=0$. 
	 }
\end{ex}

% nguồn: Ok_Toan11_TracNghiem_Lan2
\begin{ex}%[1D3C1-2]%[ma:ef712e72]
	Cho dãy số $\left(u_n \right)$ biết  $\heva{&{u_1=-2} \\
	&{u_n=3u_{n-1}-1, \forall n\ge  2}
	}$ khi đó $L=\lim\dfrac{u_n}{3^n}$ 
	 \choice
		{Không xác  định}
		 {$L=+\infty$}
		 {\True $L=-\dfrac{5}{6}$}
		{$L=0$} 
	\loigiai{Đặt $u_n=v_n+\dfrac{1}{2}$ thay vào biểu thức truy hồi ta có $v_n+\dfrac{1}{2}=3\left(v_{n-1}+\dfrac{1}{2} \right)-1\Leftrightarrow v_n=3v_{n-1}, \forall n\ge 2$. Dễ thấy $\left(v_n \right)$ là cấp số nhân với $v_1=u_1-\dfrac{1}{2}=-2-\dfrac{1}{2}=-\dfrac{5}{2}$ công bội  $q=3$ suy ra $v_n=-\dfrac{5}{2}.3^{n-1}$. Do đó $u_n=v_n+\dfrac{1}{2}=-\dfrac{5}{2}.3^{n-1}+\dfrac{1}{2} \left(n\ge 1\right)$. Vậy $L=\lim\dfrac{u_n}{3^n}=\lim\left(-\dfrac{5}{6}+\dfrac{1}{2.3^n} \right)=-\dfrac{5}{6}$.
	 }
\end{ex}

% nguồn: Ok_Toan11_TracNghiem_Lan2
\begin{ex}%[1D3C1-2]%[ma:f9d22915]
	Cho dãy số $\left(u_n \right)$ như sau: $u_n=\dfrac{n}{1+n^2+n^4}$, $\forall n=1$, $2$, $\ldots$. Tính giới  hạn $\lim\limits_{n\to+\infty} \left(u_1+u_2+\cdots+u_n \right)$
	\choice 
		{$\dfrac{1}{4}$} 
		{$1$} 
		{\True  $\dfrac{1}{2}$}
		{$\dfrac{1}{3}$} 
	\loigiai{Ta có  $u_n=\dfrac{n}{\left(1+n^2\right)^2-n^2}=\dfrac{n}{\left(n^2+n+1\right)\left(n^2-n+1\right)}=\dfrac{1}{2} \left(\dfrac{1}{n^2-n+1}-\dfrac{1}{n^2+n+1} \right)$.\\
		Ta  có
		\allowdisplaybreaks 
		\begin{eqnarray*}
			u_1+u_2+\cdots+u_n&=&\dfrac{1}{2} \left(1-\dfrac{1}{3}+\dfrac{1}{3}-\dfrac{1}{7}+\dfrac{1}{7}-\dfrac{1}{13}+\dfrac{1}{13}-\dfrac{1}{21}+\cdots+\dfrac{1}{n^2-n+1}-\dfrac{1}{n^2+n+1} \right)\\ 
			&=&\dfrac{1}{2} \left(1-\dfrac{1}{n^2+n+1} \right)=\dfrac{1}{2} \dfrac{n^2+n}{n^2+n+1}. 
		\end{eqnarray*}
		Suy  ra $\lim\left(u_1+u_2+\cdots+u_n \right)=\dfrac{1}{2} \lim\dfrac{1+\dfrac{1}{n}}{1+\dfrac{1}{n}+\dfrac{1}{n^2}}=\dfrac{1}{2}$.
	} 
\end{ex}

% nguồn: Ok_Toan11_TracNghiem_Lan2
\begin{ex}%[1D3C1-2]%[ma:012bedab]
	Cho dãy số $\left(u_n \right)$ xác định bởi $u_1=0$  và $u_{n+1}=u_n+4n+3, \forall n\ge 1$. Biết
	 \[\lim\dfrac{\sqrt{u_n}+\sqrt{u_{4n}}+\sqrt{u_{4^2n}}+\cdots+\sqrt{u_{4^{2018} n}}}{\sqrt{u_n}+\sqrt{u_{2n}}+\sqrt{u_{2^2n}}+\cdots+\sqrt{u_{2^{2018} n}}}=\dfrac{a^{2019}+b}{c}
	\] 
	với $a$, $b$, $c$ là các số  nguyên dương và $b < 2019$. Tính giá trị $S=a+b-c$.
	 \choice
		{$S=-1$} 
		{\True  $S=0$}
		{$S=2017$} 
		 {$S=2018$}
	\loigiai{Ta có 
		\[\heva{u_2=u_1+4.1+3 \\ u_3=u_2+4\cdot 2+3  \\ \ldots \\ u_n=u_{n-1}+4\left(n-1\right)+3}
		\] 
		Cộng vế theo vế và rút  gọn ta được\\
		$u_n=u_1+4.\left(1+2+\cdots+n-1\right)+3\left(n-1\right)=4\dfrac{n\left(n-1\right)}{2}+3\left(n-1\right)=2n^2+n-3$ với mọi $n\ge  1$. Suy ra
		\[\heva{u_{2n}=2(2n)^2+2n-3 \\ u_{2^2n}=2(2^2n)^2+2^2n-3 \\ \ldots \\ u_{2^{2018} n}=2\left(2^{2018} n\right)^2+2^{2018} n-3} 
		 \]
		 Và
		\[\heva{&u_{4n}=2(4n)^2+4n-3 \\ &u_{4^2n}=2(4^2n)^2+4^2n-3 \\ &\ldots  \\ &u_{4^{2018} n}=2\left(4^{2018} n\right)^2+4^{2018} n-3}
		 \]
		Do đó  $\lim\dfrac{\sqrt{u_n}+\sqrt{u_{4n}}+\sqrt{u_{4^2n}}+\cdots+\sqrt{u_{4^{2018} n}}}{\sqrt{u_n}+\sqrt{u_{2n}}+\sqrt{u_{2^2n}}+\cdots+\sqrt{u_{2^{2018} n}}}$
		\[=\lim\dfrac{\sqrt{2+\dfrac{1}{n}-\dfrac{3}{n^2}}+\sqrt{2.4^2+\dfrac{4}{n}-\dfrac{3}{n^2}}+\cdots+\sqrt{2\left(4^{2018}  \right)^2+\dfrac{4^{2018}}{n}-\dfrac{3}{n^2}}}{\sqrt{2+\dfrac{1}{n}-\dfrac{3}{n^2}}+\sqrt{2.2^2+\dfrac{2}{n}-\dfrac{3}{n^2}}+\cdots+\sqrt{2\left(2^{2018} \right)^2+\dfrac{2^{2018}}{n}-\dfrac{3}{n^2}}}
		 \]
		\[=\dfrac{\sqrt{2} \left(1+4+4^2+\cdots+4^{2018} \right)}{\sqrt{2}  \left(1+2+2^2+\cdots+2^{2018} \right)}=\dfrac{1\dfrac{1-4^{2019}}{1-4}}{\dfrac{1-2^{2019}}{1-2}}=\dfrac{1}{3} \dfrac{4^{2019}-1}{2^{2019}-1}=\dfrac{2^{2019}+1}{3}.
		\] 
		Vì $2^{2019} > 2019$ cho nên  sự xác định ở trên là duy nhất nên $\heva{&{a=2} \\
		&{b=1} \\ 
		 &{c=3}
		}$.\\ 
		 Vậy $S=a+b-c=0$.
	 }
\end{ex}

% nguồn: Ok_TracNghiemLop11_Lan1
\begin{ex}%[1D3B1-2]%[ma:08ec40b2]
	Tính giới hạn  $L=\lim \dfrac{n+1}{2n-3}$.
	 \choice{\True $L=\dfrac{1}{2}$}
		 {$L=1$}
		 {$L=2$}
		{$L=\dfrac{1}{6}$} 
	 \loigiai{
		Ta có $\lim \dfrac{n+1}{2n-3}= \lim  \dfrac{1+\dfrac{1}{n}}{2-\dfrac{3}{n}}=\dfrac{1}{2}$.
	 }
\end{ex}


%%%%%%%%%%%%%%%% ĐÚNG/SAI %%%%%%%%%%%%%%%%

% nguồn: Ok_dungsai_lop11_Lan1
\begin{ex}%[1D3N1-2]%[ma:bd2ff02c]
 Mỗi kết quả giới  hạn dưới đây đúng hay sai?
  \choiceTF
  {$\lim\limits_{n  \to +\infty} \dfrac{2n+3}{2n-1}=2$}
  {\True $\lim\limits_{n  \to +\infty} \dfrac{2n^2+3}{n^2+1}=2$}
  {$\lim\limits_{n \to +\infty} \dfrac{2n^2+n+3}{3-n^2}=\dfrac{2}{3}$} 
  {\True  $\lim\limits_{n \to +\infty} \dfrac{2n+3}{n^2+2}=0$}
 \loigiai{ 
  \begin{itemchoice}
   \itemch $\lim\limits_{n \to +\infty} \dfrac{2n+3}{2n-1}=\lim\limits_{n \to +\infty} \dfrac{n\left(2+\dfrac{3}{n}\right)}{n\left(2-\dfrac{1}{n}\right)}=\lim\limits_{n  \to +\infty} \dfrac{2+\dfrac{3}{n}}{2-\dfrac{1}{n}}=1$.
   \itemch $\lim\limits_{n \to +\infty} \dfrac{2n^2+3}{n^2+1}=\lim\limits_{n \to +\infty} \dfrac{n^2\left(2+\dfrac{3}{n^2}\right)}{n^2\left(1+\dfrac{1}{n^2}\right)}=\lim\limits_{n  \to +\infty} \dfrac{2+\dfrac{3}{n^2}}{1+\dfrac{1}{n^2}}=2$.
   \itemch $\lim\limits_{n \to +\infty} \dfrac{2n^2+n+3}{3-n^2}=\lim\limits_{n \to +\infty} \dfrac{n^2\left(2+\dfrac{1}{n}+\dfrac{3}{n^2}\right)}{n^2\left(\dfrac{3}{n^2}-1\right)}=\lim\limits_{n  \to +\infty} \dfrac{2+\dfrac{1}{n}+\dfrac{3}{n^2}}{\dfrac{3}{n^2}-1}=-2$.
   \itemch $\lim\limits_{n \to +\infty}  \dfrac{2n+3}{n^2+2}=\lim\limits_{n \to +\infty} \dfrac{n^2\left(\dfrac{2}{n}+\dfrac{3}{n^2}\right)}{n^2\left(1+\dfrac{2}{n^2}\right)}=\lim\limits_{n \to +\infty} \dfrac{\dfrac{2}{n}+\dfrac{3}{n^2}}{1+\dfrac{2}{n^2}}=0$.
  \end{itemchoice}
  }
\end{ex}

% nguồn: Ok_dungsai_lop11_Lan1
\begin{ex}%[1D3H1-2]%[ma:68b21c0f]
 Biết giới hạn $\lim \dfrac{2n^2+1}{3n^3-3n+3}=a$ và $\lim \dfrac{n\sqrt{n^2+1}}{\sqrt{4n^4-n^2+3}}=b$.  Khi đó:
  \choiceTF
  {Giá trị $a$ nhỏ hơn  $0$}
  {\True  Giá trị $b$ lớn hơn $0$}
  {\True Phương  trình lượng giác $\cos x=a$ có một nghiệm là $x=\dfrac{\pi}{2}$}
  {Cho cấp số cộng $\left(u_n \right)$ với công sai $d=b$ và $u_1=a$,  thì $u_3=\dfrac{3}{2}$}
  \loigiai{
  \begin{itemchoice}
   \itemch Sai. Ta có $\lim  \dfrac{2n^2+1}{3n^3-3n+3}=\lim \dfrac{n^3 \left(\dfrac{2}{n}+\dfrac{1}{n^3} \right)}{n^3 \left(3-\dfrac{3}{n^2}+\dfrac{3}{n^3} \right)}=\lim \dfrac{\dfrac{2}{n}+\dfrac{1}{n^3}}{3-\dfrac{3}{n^2}+\dfrac{3}{n^3}}=\dfrac0{3}=0$.
   \itemch Đúng. Ta có $\lim \dfrac{n\sqrt{n^2+1}}{\sqrt{4n^4-n^2+3}}=\lim \dfrac{n^2 \sqrt{1+\dfrac{1}{n^2}}}{n^2 \sqrt{4-\dfrac{1}{n^2}+\dfrac{3}{n^4}}}=\lim \dfrac{\sqrt{1+\dfrac{1}{n^2}}}{\sqrt{4-\dfrac{1}{n^2}+\dfrac{3}{n^4}}}=\dfrac{1}{2}$. 
   \itemch Đúng. Phương trình lượng giác $\cos x=0$ có một nghiệm  là $x=\dfrac{\pi}{2}$.
   \itemch Sai. Cho cấp số cộng $\left(u_n \right)$  với công sai $d=\dfrac{1}{2}$ và $u_1=0$, thì $u_3=0+2\cdot\dfrac{1}{2}=1$.
  \end{itemchoice}
 } 
\end{ex}

% nguồn: Ok_dungsai_lop11_Lan1
\begin{ex}%[1D3H1-2]%[ma:521dc5b4]
 Biết giới hạn $\lim  \left(-2n^3-5n+9\right)=a$ và $\lim \dfrac{4^n+3}{1+3\cdot 4^{n+1}}=b$. Khi đó:
  \choiceTF
  {Tích $a\cdot b=3$} 
  {\True Hàm  số $y=\sqrt{1-x}$ có tập xác định là $\mathscr{D}=\left(a;1\right]$}
  {\True Giá trị $b$ là số lớn  hơn $0$}
  {Phương trình lượng giác $\cos  x=b$ vô nghiệm}
  \loigiai{
  Ta có $\lim \left(-2n^3-5n+9\right)=\lim n^3 \left(-2-\dfrac{5}{n^2}+\dfrac{9}{n^3} \right)=-\infty$, do $\left\{\begin{array}{l}{\lim n^3=+\infty} \\ {\lim \left(-2-\dfrac{5}{n^2}+\dfrac{9}{n^3} \right)=-2}
 \end{array}\right.$
 và \[\lim \dfrac{4^n+3}{1+3\cdot 4^{n+1}}=\lim \dfrac{4^n+3}{1+12\cdot 4^n}=\lim \dfrac{4^n \left(1+\dfrac{3}{4^n} \right)}{4^n  \left(\dfrac{1}{4^n}+12\right)}=\lim \dfrac{1+\dfrac{3}{4^n}}{\dfrac{1}{4^n}+12}=\dfrac{1}{12}
  \]
 Do đó $a=-\infty,  b=\dfrac{1}{12}$
 \begin{itemchoice}
  \itemch Sai.  Tích $a\cdot b=-\infty$.
   \itemch Đúng. Hàm số $y=\sqrt{1-x}$ có tập xác định là $\mathscr{D}=\left(-\infty;1\right]$.
  \itemch Đúng. Giá trị $\dfrac{1}{12}$  là số lớn hơn $0$.
  \itemch Sai. Phương trình lượng giác $\cos x=\dfrac{1}{12}$ có  nghiệm.
 \end{itemchoice}
 } 
\end{ex}

% nguồn: Ok_dungsai_lop11_Lan1
\begin{ex}%[1D3H1-2]%[ma:b23b6817]
 Biết giới hạn $\lim \dfrac{-3n^3+1}{2n+5}=a$ và  $\lim \dfrac{(-1)^n \cdot 5^n}{2^n+5^{2n}}=b$. Khi đó:
  \choiceTF
  {\True $\lim \left(-3n^2+\dfrac{1}{n}  \right)=a$}
  {\True $x=b$ là hoành độ giao  điểm của đường thẳng $y=2x$ với trục hoành}
  {\True $\lim \left(\dfrac{1}{2024}  \right)^n=b$}
  {Cho cấp số cộng $\left(u_n \right)$ với công sai $d=\dfrac{1}{2}$ và $u_1=b$, thì $u_3=2$} 
  \loigiai{
  Ta có  $\lim \dfrac{-3n^3+1}{2n+5}=\lim \dfrac{n\left(-3n^2+\dfrac{1}{n} \right)}{n\left(2+\dfrac{5}{n} \right)}=\lim \dfrac{-3n^2+\dfrac{1}{n}}{2+\dfrac{5}{n}}=-\infty$.\\
  Vì $\heva{&{\lim \left(-3n^2+\dfrac{1}{n} \right)=-\infty}  \\
  &{\lim \left(2+\dfrac{5}{n} \right)=2} 
   }$
  nên $a=-\infty, b=0$. Do  đó:
  \begin{itemchoice}
   \itemch Đúng. $\lim \left(-3n^2+\dfrac{1}{n} \right)=-\infty=  a$.
   \itemch Đúng. Giao điểm của đường thẳng $y=2x$  với trục hoành có tọa độ $(0,0)$.
   \itemch Đúng. $\lim \left(\dfrac{1}{2024} \right)^n=0=b$. 
   \itemch Sai. Cấp số cộng $\left(u_n \right)$ với công sai $d=\dfrac{1}{2}$ và $u_1=b=0$,  thì $u_3=1$.
  \end{itemchoice}
 } 
\end{ex}

% nguồn: Ok_dungsai_lop11_Lan1
\begin{ex}%[1D3H1-2]%[ma:10f0bdf6]
 Biết giới hạn $\lim  \dfrac{2n+1}{-3n+2}=a$. Khi đó:
 \choiceTF 
  {Giá trị $a$ lớn hơn  $0$}
  {Ba số $-\dfrac{5}{3}$; $a$; $\dfrac{1}{3}$ tạo thành một cấp số cộng  với công sai bằng $2$}
  {Trên khoảng $\left(-\pi;\pi \right)$ phương trình lượng giác  $\sin x=a$ có $3$ nghiệm}
  {\True Cho cấp số nhân $\left(u_n \right)$ với công  bội $q=3$ và $u_1=a$, thì $u_3=-6$}
  \loigiai{
  \begin{itemchoice}
   \itemch Sai. Ta có $\lim \dfrac{2n+1}{-3n+2}=\lim \dfrac{n\left(2+\dfrac{1}{n} \right)}{n\left(-3+\dfrac{2}{n}  \right)}=\lim \dfrac{2+\dfrac{1}{n}}{-3+\dfrac{2}{n}}=\dfrac{-2}{3}$.
   \itemch Sai. Ba số $-\dfrac{5}{3};-\dfrac{2}{3};\dfrac{1}{3}$ tạo thành một cấp số cộng với công  sai bằng $1$.
   \itemch Sai. Trên khoảng $\left(-\pi;\pi \right)$ phương trình lượng  giác $\sin x=a$ có $2$ nghiệm.
   \itemch Đúng. Cho cấp số nhân $\left(u_n \right)$ với công bội  $q=3$ và $u_1=a$, thì $u_3=-6$.
  \end{itemchoice}
 } 
\end{ex}

% nguồn: Ok_dungsai_lop11_Lan1
\begin{ex}%[1D3H1-2]%[ma:a00f78f2]
 Biết giới hạn $\lim \dfrac{5n^3-2n+1}{n-2n^3}=a$. Khi  đó:
  \choiceTF
  {\True Giá trị $a$ nhỏ hơn  $0$}
  {\True $x=a$ là trục  đối xứng của parabol $(P)\colon y=x^2+5x+2$}
  {\True Phương  trình lượng giác $\sin x=a$ vô nghiệm}
  {Cho cấp số cộng $\left(u_n \right)$ với công sai $d=3$ và $u_1=a$, thì $u_3=6$} 
 \loigiai{ 
  \begin{itemchoice}
    \itemch Đúng. Ta có $\lim \dfrac{5n^3-2n+1}{n-2n^3}=\lim \dfrac{n^3 \left(5-\dfrac{2}{n^2}+\dfrac{1}{n^3} \right)}{n^3 \left(\dfrac{1}{n^2}-2\right)}=\lim \dfrac{5-\dfrac{2}{n^2}+\dfrac{1}{n^3}}{\dfrac{1}{n^2}-2}=-\dfrac{5}{2}$.
   \itemch Đúng. Parabol $(P)\colon y=x^2+5x+2$ nhận $x=-\dfrac{5}{2}$ làm trục  đối xứng.
   \itemch  Đúng. Phương trình lượng giác $\sin x=-\dfrac{5}{2}$ vô nghiệm.
   \itemch Sai. Cho cấp số  cộng $\left(u_n \right)$ với công sai $d=3$ và $u_1=a$, thì $$u_3=u_1+\left(3-1\right)d=-\dfrac{5}{2}+2\cdot 3=\dfrac{7}{2}.$$
  \end{itemchoice}
 } 
\end{ex}

% nguồn: Ok_dungsai_lop11_Lan1
\begin{ex}%[1D3H1-2]%[ma:37fb4fd5]
 Cho các dãy số $(a_n)$,  $(b_n)$, $(c_n)$ với $a_n=1-2n$, $b_n=n^2+n+1$, $c_n=-n^3+2n+1$, $d_n=2-3n^2$.
 Mỗi kết quả giới hạn  dưới đây đúng hay sai?
 \choiceTF 
   {\True $\lim\limits_{n \to +\infty} \dfrac{a_n}{b_n}=0$}
  {\True $\lim\limits_{n \to  +\infty} \dfrac{a_n}{c_n}=0$}
  {$\lim\limits_{n  \to +\infty} \dfrac{b_n}{d_n}=-\dfrac{1}{2}$}
  {\True $\lim\limits_{n \to  +\infty} \dfrac{d_n}{c_n}=0$}
  \loigiai{
  \begin{itemchoice}
   \itemch $\lim\limits_{n  \to +\infty} \dfrac{a_n}{b_n}
   =\lim\limits_{n \to  +\infty} \dfrac{1-2n}{n^2+n+1}
    =\lim\limits_{n \to +\infty} \dfrac{n^2\left(\dfrac{1}{n^2}-\dfrac{2}{n}\right)}{n^2\left(1+\dfrac{1}{n}+\dfrac{1}{n^2}\right)}
   =\lim\limits_{n  \to +\infty} \dfrac{\dfrac{1}{n^2}-\dfrac{2}{n}}{1+\dfrac{1}{n}+\dfrac{1}{n^2}}=0$.
    \itemch $\lim\limits_{n \to +\infty} \dfrac{a_n}{c_n}
   =\lim\limits_{n \to  +\infty} \dfrac{1-2n}{-n^3+2n+1}
   =\lim\limits_{n \to +\infty} \dfrac{n^3\left(\dfrac{1}{n^3}-\dfrac{2}{n^2}\right)}{n^3\left(-1+\dfrac{2}{n^2}+\dfrac{1}{n^3}\right)} 
   =\lim\limits_{n \to +\infty}  \dfrac{\dfrac{1}{n^3}-\dfrac{2}{n^2}}{-1+\dfrac{2}{n^2}+\dfrac{1}{n^3}}=0$.
   \itemch $\lim\limits_{n \to +\infty} \dfrac{b_n}{d_n}=\lim\limits_{n \to +\infty} \dfrac{n^2+n+1}{2-3n^2}=\lim\limits_{n \to +\infty} \dfrac{n^2\left(1+\dfrac{1}{n}+\dfrac{1}{n^2}\right)}{n^2\left(\dfrac{2}{n^2}-3\right)}=\lim\limits_{n  \to +\infty} \dfrac{1+\dfrac{1}{n}+\dfrac{1}{n^2}}{\dfrac{2}{n^2}-3}=-\dfrac{1}{3}$.
   \itemch $\lim\limits_{n \to +\infty} \dfrac{d_n}{c_n} 
   =\lim\limits_{n \to +\infty} \dfrac{2-3n^2}{-n^3+2n+1} 
   =\lim\limits_{n \to +\infty} \dfrac{n^3\left(\dfrac{2}{n^3}-\dfrac{3}{n}\right)}{n^3\left(-1+\dfrac{2}{n^2}+\dfrac{1}{n^3}\right)} 
   =\lim\limits_{n \to  +\infty} \dfrac{\dfrac{2}{n^3}-\dfrac{3}{n}}{-1+\dfrac{2}{n^2}+\dfrac{1}{n^3}}=0$.
  \end{itemchoice}
 } 
\end{ex}

% nguồn: Ok_dungsai_lop11_Lan1
\begin{ex}%[1D3H1-2]%[ma:5698beae]
  Cho $u_n=\dfrac{7^n+2^{2n-1}+3^{n+1}}{7^{n+1}+5^{n-1}}$. Biết $\lim u_n=\dfrac{a}{b}$ (với $a,b\in \mathbb{Z};\dfrac{a}{b}$ tối giản). Khi đó:
  \choiceTF
  {\True  $a+b=8$}
  {$a-b=-7$} 
  {\True Bộ ba số $a$; $b$; $13$ tạo thành một cấp số cộng có công sai $d=7$} 
  {\True Bộ ba số $a$; $b$; $49$ tạo thành  một cấp số nhân có công bội $q=7$}
  \loigiai{
  Ta có $\lim u_n=\lim \dfrac{7^n+2^{2n-1}+3^{n+1}}{7^{n+1}+5^{n-1}}=\lim \dfrac{1+\dfrac{1}{2} \left(\dfrac{4}{7}  \right)^n+3\left(\dfrac{3}{7} \right)^n}{7+\dfrac{1}{5} \left(\dfrac{5}{7} \right)^n}=\dfrac{1}{7}$.\\
  Do  đó suy ra $a=1,b=7$.
  \begin{itemchoice}
   \itemch Đúng.  $a+b=8$.
   \itemch  Sai. $a-b=-6.$
   \itemch Đúng. Bộ ba số $a$; $b$; $13$ tạo thành một cấp số cộng có công sai $d=7$  vì $a+13=14=2b$.
   \itemch Đúng. Bộ ba số  $a$; $b$; $49$ tạo thành một cấp số nhân có công bội $q=7$ vì $a\cdot 49 = b^2$.
  \end{itemchoice}
 } 
\end{ex}

% nguồn: Ok_dungsai_lop11_Lan2
\begin{ex}%[1D3H1-2]%[ma:6b0d8dd8]
	Cho dãy số $\left(u_n\right)$ với $u_n=a n^2+n-1$ với  $a \in \mathbb{R}$.
	 \choiceTF
		{Với $a=1$, giới hạn của dãy số đã cho  là $1$}
		{\True Với $a=2$, giới hạn của dãy số đã cho  là $+\infty$}
		{\True Với $a=-\dfrac{5}{2}$, giới hạn của dãy số đã cho là  $-\infty$}
		 {Với $a \leq 0$, giới hạn của dãy số đã cho là $-\infty$}
	\loigiai{ 
		\begin{itemchoice}
			\itemch Với $a=1$, $u_n=n^2+n-1$.\\ 
			Ta có $\lim \left(n^2+n-1\right)=\lim \left[n^2\left(1+\dfrac{1}{n}-\dfrac{1}{n^2}\right)\right]=+\infty$.\\ 
			Vì $\lim n^2=+\infty$ và  $\lim \left(1+\dfrac{1}{n}-\dfrac{1}{n^2}\right)=1$.
			\itemch Với $a=2$,  $u_n=2n^2+n-1$.\\
			Ta có $\lim \left(2n^2+n-1\right)=\lim \left[n^2\left(2+\dfrac{1}{n}-\dfrac{1}{n^2}\right)\right]=+\infty$.\\ 
			Vì  $\lim n^2=+\infty$ và $\lim \left(2+\dfrac{1}{n}-\dfrac{1}{n^2}\right)=2$.
			\itemch Với  $a=-\dfrac{5}{2}$, $u_n=-\dfrac{5}{2} n^2+n-1$.\\
			 Ta có $\lim \left(-\dfrac{5}{2} n^2+n-1\right)=\lim\left[n^2\left(-\dfrac{5}{2}+\dfrac{1}{n}-\dfrac{1}{n^2}\right)\right]=-\infty$.\\
			 Vì $\lim n^2=+\infty$ và $\lim \left(-\dfrac{5}{2}+\dfrac{1}{n}-\dfrac{1}{n^2}\right)=-\dfrac{5}{2}$.
			\itemch Với $a < 0$. Khi  đó $$\lim \left(a n^2+n-1\right)=\lim\left[n^2\left(a+\dfrac{1}{n}-\dfrac{1}{n^2}\right)\right]=-\infty.$$
			 Vì $\lim n^2=+\infty$ và $\lim \left(a+\dfrac{1}{n}-\dfrac{1}{n^2}\right)=a<0$.\\
			Với $a=0$,  $u_n=n-1$. Khi đó
			 $$\lim (n-1)=\lim \left[n\left(1-\dfrac{1}{n}\right)\right]=+\infty.$$
			Vì $\lim  n=+\infty$ và $\lim \left(1-\dfrac{1}{n}\right)=1$.
		\end{itemchoice}
	} 
\end{ex}

% nguồn: Ok_dungsai_lop11_Lan2
\begin{ex}%[1D3H1-2]%[ma:9b680c3b]
	Cho hai dãy số $\left(u_n\right)$ và $\left(v_n\right)$ có $u_n=4n^2-n+3$; $v_n=3n^2+7$. 
	\choiceTF 
		{\True  $\lim \dfrac{u_n}{v_n}=\dfrac{4}{3}$}
		{$\lim \dfrac{u_n}{\left(v_n\right)^2}=\dfrac{4}{9}$} 
		{$\lim  \dfrac{\left(u_n\right)^2}{v_n}=\dfrac{16}{3}$}
		{\True $\lim \dfrac{u_n+an^2+7}{v_n}=8$ khi đó  $a=20$}
	\loigiai{ 
		\begin{itemchoice}
			 \itemch Ta có\\
			$\begin{aligned}
			\lim \dfrac{u_n}{v_n}=\lim \dfrac{4n^2-n+3}{3n^2+7} 
			=&\lim  \dfrac{n^2\left(4-\dfrac{1}{n}+\dfrac{3}{n^2}\right)}{n^2\left(3+\dfrac{7}{n^2}\right)}\\
			=&\lim  \dfrac{4-\dfrac{1}{n}+\dfrac{3}{n^2}}{3+\dfrac{7}{n^2}}
			 =\dfrac{4}{3}.
		\end{aligned}$
		\itemch Ta có\\ 
		$\begin{aligned}
		\lim \dfrac{u_n}{\left(v_n\right)^2}=\lim \dfrac{4n^2-n+3}{(3n^2+7)^2} 
		=&\lim  \dfrac{n^2\left(4-\dfrac{1}{n}+\dfrac{3}{n^2}\right)}{n^4\left(3+\dfrac{7}{n^2}\right)^2}\\
		=&\lim \left[\dfrac{1}{n^2}\cdot \dfrac{4-\dfrac{1}{n}+\dfrac{3}{n^2}}{\left(3+\dfrac{7}{n^2}\right)^2}\right]=0. 
	\end{aligned}$
	\itemch  Ta có\\
	$\begin{aligned}
	\lim  \dfrac{\left(u_n\right)^2}{v_n}=\lim \dfrac{(4n^2-n+3)^2}{3n^2+7}
	 =&\lim \dfrac{n^4\left(4-\dfrac{1}{n}+\dfrac{3}{n^2}\right)^2}{n^2\left(3+\dfrac{7}{n^2}\right)} \\
	=&\lim \left[n^2\cdot  \dfrac{\left(4-\dfrac{1}{n}+\dfrac{3}{n^2}\right)^2}{3+\dfrac{7}{n^2}}\right]
	 =+\infty.
	\end{aligned}$
	\itemch  Ta có\\
	$\begin{aligned}
	 \lim \dfrac{4n^2-n+3+an^2+7}{3n^2+7}=\lim \dfrac{(a+4)n^2-n+10}{3n^2+7}
	=&\lim \dfrac{n^2\left[(a+4)-\dfrac{1}{n}+\dfrac{10}{n^2}\right]}{n^2\left(3+\dfrac{7}{n^2}\right)}  \\
	=&\lim \dfrac{(a+4)-\dfrac{1}{n}+\dfrac{10}{n^2}}{3+\dfrac{7}{n^2}}=\dfrac{a+4}{3}. 
	\end{aligned}$\\
	Khi đó $\dfrac{a+4}{3}=8\Leftrightarrow a=20$. 
	\end{itemchoice}
} 
\end{ex}

% nguồn: Ok_dungsai_lop11_Lan2
\begin{ex}%[1D3H1-2]%[ma:a29a679c]
	Cho hai dãy số $\left(u_n\right)$ và $\left(v_n\right)$ có $u_n=7^{n+1}-3$;  $v_n=2\cdot 7^n+3^n$
	\choiceTF 
		{\True $\lim  \dfrac{u_n}{v_n}=\dfrac{7}{2}$}
		 {$\lim \dfrac{u_n}{\left(v_n\right)^2}=\dfrac{7}{4}$}
		 {$\lim \dfrac{\left(u_n\right)^2}{v_n}=\dfrac{49}{2}$}
		{$\lim \dfrac{u_n+a\cdot 7^n+7}{v_n}=8$ khi đó  $a=8$}
	\loigiai{ 
		\begin{itemchoice}
			\itemch  Ta có
			$\lim \dfrac{u_n}{v_n}=\lim  \dfrac{7^{n+1}-3}{2\cdot 7^n+3^n}
			=\lim  \dfrac{7^n\left(7-\dfrac{3}{7^n}\right)}{7^n\left[2+\left(\dfrac{3}{7}\right)^n\right]}
			=\lim  \dfrac{7-\dfrac{3}{7^n}}{2+\left(\dfrac{3}{7}\right)^n}
			=\dfrac{7}{2}$. 
			\itemch Ta  có
			 $\lim \dfrac{u_n}{\left(v_n\right)^2}
			=\lim \dfrac{7^{n+1}-3}{(2\cdot 7^n+3^n)^2} 
			 =\lim \dfrac{7^n\left(7-\dfrac{3}{7^n}\right)}{7^{2n}\left[2+\left(\dfrac{3}{7}\right)^n\right]^2}
			=\lim \left[\dfrac{1}{7^n}\cdot \dfrac{7-\dfrac{3}{7^n}}{\left(2+\left(\dfrac{3}{7}\right)^n\right)^2}\right] 
			 =0$.
			\itemch Ta có 
			$\lim  \dfrac{\left(u_n\right)^2}{v_n}
			 =\lim \dfrac{(7^{n+1}-3)^2}{2\cdot 7^n+3^n}
			=\lim  \dfrac{7^{2n}\left(7-\dfrac{3}{7^n}\right)^2}{7^n\left[2+\left(\dfrac{3}{7}\right)^n\right]}
			 =\lim \left[7^n\cdot \dfrac{\left(7-\dfrac{3}{7^n}\right)^2}{2+\left(\dfrac{3}{7}\right)^n}\right]
			=+\infty$. 
			\itemch  Ta có
			$\lim \dfrac{7^{n+1}-3+a\cdot 7^n+7}{2\cdot  7^n+3^n}
			=\lim \dfrac{(a+7)\cdot 7^n+4}{7^n\left[2+\left(\dfrac{3}{7}\right)^n\right]} 
			=\lim  \dfrac{7^n\left[(a+7)+\dfrac{4}{7^n}\right]}{7^n\left[2+\left(\dfrac{3}{7}\right)^n\right]}
			=\dfrac{a+7}{2}$. \\ 
			Khi đó $\dfrac{a+7}{2}=8\Leftrightarrow a=9$. 
		\end{itemchoice}
	 }
\end{ex}

% nguồn: Ok_dungsai_lop11_Lan2
\begin{ex}%[1D3H1-2]%[ma:2b0c596a]
	Cho giới hạn  $L=\lim \dfrac{\sqrt[3]{a n^3+5n^2-7}}{\sqrt{3n^2-n+2}}$.
	 \choiceTF
		{\True Khi $a=1$  thì $L=\dfrac{1}{\sqrt{3}}$}
		{Khi $a=0$ thì  $L=\dfrac{1}{3}$}
		 {\True Khi $a > 0$ thì $L> 0$}
		{Khi $L=b \sqrt{3}+c$ với $b$, $c$ là các tham số  thì $P=\dfrac{a+c}{b^3}=\dfrac{1}{2}$}
	\loigiai{ 
		Ta  có
		\begin{eqnarray*}
			\lim  \dfrac{\sqrt[3]{a n^3+5n^2-7}}{\sqrt{3n^2-n+2}}&=&\lim \dfrac{\sqrt[3]{a+\dfrac{5}{n}-\dfrac{7}{n^3}}}{\sqrt{3-\dfrac{1}{n}+\dfrac{2}{n^2}}}\\
			&=&\dfrac{\sqrt[3]a}{\sqrt{3}}. 
		\end{eqnarray*}
		\begin{itemchoice}
			\itemch Với $a=1$  thì $L=\dfrac{1}{\sqrt{3}}$.
			\itemch Với $a=0$  thì $L=0$.
			\itemch Với $a  > 0$ thì $\sqrt[3]a>0$, suy ra $\dfrac{\sqrt[3]a}{\sqrt{3}}>0$ hay $L>0$.
			\itemch Ta  có $L=\dfrac{\sqrt[3]a}{\sqrt{3}}=\dfrac{\sqrt[3]a\cdot \sqrt{3}}{3}$.\\
			Khi $L=b \sqrt{3}+c$ với $b$,  $c$ là các tham số thì
			$$\heva{&b=\dfrac{\sqrt[3]a}{3} \\&c=0}\Leftrightarrow\heva{&b^3=\dfrac{a}{27}  \\&c=0.}$$
			Vậy  $P=\dfrac{a+c}{b^3}=\dfrac{a}{\dfrac{a}{27}}=27$.
		\end{itemchoice}
	} 
\end{ex}

% nguồn: Ok_dungsai_lop11_Lan2
\begin{ex}%[1D3H1-2]%[ma:cc9c02d4]
	Cho dãy  số $\left(u_n\right)$ có số hạng tổng quát $u_n=\dfrac{9^n+3^{n+2}}{6^n+9^{n+a}}$.
	\choiceTF 
		{\True Khi  $a=0$ thì $\lim\, u_n=1$}
		{Khi $a=1$ thì  $\lim\, u_n=9$}
		{Khi  $a=-1$ thì $\lim\, u_n=\dfrac{1}{9}$}
		{Có $2022$ giá trị nguyên của tham số  $a$ thuộc đoạn $\left[0;2024\right]$ để $\lim\, u_n\le \dfrac{1}{81}$}
	 \loigiai{
		\begin{itemchoice}
			\itemch Khi $a=0$ thì  $\lim\, u_n=\lim \dfrac{9^n+3^{n+2}}{6^n+9^n}=\lim \dfrac{1+9\cdot{\left(\dfrac{3}{9}\right)^n}}{\left(\dfrac{6}{9}\right)^n+1}=1$.
			 \itemch Khi $a=1$ thì $\lim\, u_n=\lim \dfrac{9^n+3^{n+2}}{6^n+9^{n+1}}=\lim \dfrac{1+9\cdot{\left(\dfrac{3}{9}\right)^n}}{\left(\dfrac{6}{9}\right)^n+9}=\dfrac{1}{9}$.
			\itemch Khi $a=-1$ thì $\lim\, u_n=\lim  \dfrac{9^n+3^{n+2}}{6^n+9^{n-1}}=\lim \dfrac{1+9\cdot{\left(\dfrac{3}{9}\right)^n}}{\left(\dfrac{6}{9}\right)^n+\dfrac{1}{9}}=9$.
			\itemch Ta có $\lim\, u_n=\lim \dfrac{9^n+3^{n+2}}{6^n+9^{n+a}}=\lim  \dfrac{1+9\cdot{\left(\dfrac{3}{9}\right)^n}}{\left(\dfrac{6}{9}\right)^n+9^a}=\dfrac{1}{9^a}$.\\
			Khi đó $\lim\, u_n\le \dfrac{1}{81}\Leftrightarrow  \dfrac{1}{9^a}\le \dfrac{1}{81}\Leftrightarrow a\ge 2$.\\
			 Mặt khác $a$ nguyên và thuộc đoạn $\left[0;2024\right]$ nên $a\in \left\{2;3;4;\ldots;2024\right\}$. \\
			Vậy có $2023$ giá trị nguyên của  $a$ thỏa mãn.
		\end{itemchoice}
	 }
\end{ex}

% nguồn: Ok_dungsai_lop11_Lan2
\begin{ex}%[1D3H1-2]%[ma:bc6f1274]
	Biết giới hạn $\lim \dfrac{2n+1}{2-3n}=a$. 
	 \choiceTF
		{Giá trị $a$ lớn  hơn $0$}
		{\True $x=a$ là trục đối xứng của parabol $(P)\colon y=3x^2+4x-7$} 
		{Bộ ba số $-\dfrac{5}{3}$; $a$; $\dfrac{1}{3}$ lập thành một cấp số cộng với công sai  bằng $2$}
		{Cho cấp số nhân $(u_n)$ với công bội $q=3$ và  $u_1=a$ thì $u_3=6$}
	\loigiai{ 
		\begin{itemchoice}
			 \itemch Ta có
			\begin{eqnarray*}
				\lim \dfrac{2n+1}{-3n+2}&=&\lim  \dfrac{n\left(2+\dfrac{1}{n}\right)}{n\left(-3+\dfrac{2}{n}\right)}\\
				&=&\lim  \dfrac{2+\dfrac{1}{n}}{-3+\dfrac{2}{n}}\\
				&=&\dfrac{-2}{3}. 
			\end{eqnarray*}
			 Suy ra $a=-\dfrac{2}{3}$.
			 \itemch Trục đối xứng của parabol $(P)\colon y=3x^2+4x-7$ là $x=-\dfrac{4}{2.3}=-\dfrac{2}{3}$.
			\itemch Ba số $-\dfrac{5}{3}$; $-\dfrac{2}{3}$; $\dfrac{1}{3}$ tạo thành một cấp số cộng  với công sai bằng $1$.
			\itemch Cho cấp số nhân $(u_n)$ với công bội  $q=3$ và $u_1=a=-\dfrac{2}{3}$.\\
			Suy ra $u_3=u_1  \cdot q^2=-\dfrac{2}{3} \cdot 3^2=-6$.
		\end{itemchoice}
	} 
\end{ex}

% nguồn: Ok_dungsai_lop11_Lan2
\begin{ex}%[1D3V1-2]%[ma:1270d51c]
	Cho giới  hạn $L=\lim \sqrt{3+\dfrac{a n^2-1}{3+n^2}-\dfrac{1}{2^n}}$. Xét tính đúng sai của các khẳng định sau:
	\choiceTF 
		{\True $L=2$ khi  $a=1$}
		{$L=3$ thì có $2$ giá trị nguyên  $a$ thỏa mãn}
		{\True $L> 3$ khi $a  > 6$}
		{\True Có $3$ giá trị nguyên  của $a$ thuộc $(0; 20)$ sao cho $\lim \sqrt{3+\dfrac{a n^2-1}{3+n^2}-\dfrac{1}{2^n}}$ là một số nguyên}
	 \loigiai{
		\begin{itemchoice}
			\itemch Ta  có
			\begin{eqnarray*}
				\lim  \dfrac{a n^2-1}{3+n^2}&=&\lim \dfrac{a-\dfrac{1}{n^2}}{\dfrac{3}{n^2}+1}\\
				&=&a. 
			\end{eqnarray*}
			Và  $\lim \dfrac{1}{2^n}=\lim \left(\dfrac{1}{2}\right)^n=0$.\\
			Suy ra $\lim \sqrt{3+\dfrac{a  n^2-1}{3+n^2}-\dfrac{1}{2^n}}=\sqrt{3+a}$.\\
			Khi đó $L=2$ thì $\sqrt{3+a}=2$ hay $3+a=4\Leftrightarrow a=1$. 
			 \itemch Khi $L=3$ thì $\sqrt{3+a}=3$ hay $3+a=9\Leftrightarrow a=6$.
			\itemch Khi $L>3$ thì $\sqrt{3+a}>3$ hay  $3+a>9\Leftrightarrow a>6$.
			\itemch Ta  có $\heva{&a \in(0; 20), a \in \mathbb{Z} \\&\sqrt{a+3} \in \mathbb{Z}} \Rightarrow a \in\{1; 6; 13\}$.
		\end{itemchoice}
	 }
\end{ex}

% nguồn: Ok_dungsai_lop11_Lan2
\begin{ex}%[1D3V1-2]%[ma:a68938d1]
	Cho các giới hạn $L_1=\lim  \dfrac{2^n}{3^n}=\lim f(n)$; $L_2=\lim \dfrac{3^n}{4^n}=\lim g(n)$ và \break $L_3=\lim\, \dfrac{a\cdot 3^n+b\cdot 2^n+c}{d\cdot 3^n+e\cdot 4^n}$.
	\choiceTF 
		 {$L_1=\dfrac{2}{3}$}
		{$L_1\ne  L_2$}
		{\True $L_1\cdot L_2=\lim  \left[f(n)\cdot g(n)\right]$}
		{Nếu $e=0$ thì tồn tại duy nhất giá trị $a$ sao  cho $2L_3=L_1+1$}
	 \loigiai{
		\begin{itemchoice}
			 \itemch Ta có $L_1=\lim{\left(\dfrac{2}{3}\right)^n}=0$.
			\itemch Ta có $L_2=\lim{\left(\dfrac{3}{4}\right)^n}=0=L_1$. 
			\itemch Do  các giới hạn $L_1$, $L_2$ xác định nên $L_1\cdot L_2=\lim \left[f(n)\cdot g(n)\right]$ xác định.
			\itemch  Nếu $e=0$ thì
			$L_3=\lim\, \dfrac{a\cdot  3^n+b\cdot 2^n+c}{d\cdot 3^n}
			=\lim\,  \dfrac{a+b\cdot{\left(\dfrac{2}{3}\right)^n}+c\left(\dfrac{1}{3}\right)^n}{d}
			=\dfrac{a}{d} 
			 \Rightarrow 2L_3=\dfrac{2a}{d}$.\\
			$2L_3=L_1+1\Leftrightarrow \dfrac{2a}{d}=1\Leftrightarrow a=\dfrac{d}{2}\,\, (d\ne 0)$ nên có vô  số giá trị $a$.
		\end{itemchoice}
	} 
\end{ex}

% nguồn: Ok_dungsai_lop11_Lan2
\begin{ex}%[1D3V1-2]%[ma:fa0327ff]
	Cho các giới hạn $\heva{&L=\lim  \dfrac{\sqrt{9n^2-n+1}}{4n-2}=\lim f(n)\\
	 &L'=\lim \left(\sqrt{4n^2+n}-\sqrt{4n^2+2}\right)=\lim g(n)\\
	 &L''=\lim \left(\sqrt{n^2+2n}-n+a\right)}$, với $a\in \mathbb{R}$ là tham số.
	\choiceTF 
		{\True  $L=\dfrac{3}{4}$}
		{\True $L+L'=1$} 
		{$\lim \left[f(n)\cdot  g(n)\right]=\dfrac{3}{8}$}
		{Tồn tại duy nhất giá  trị $a$ để $L''=\dfrac{2}{a}(L+L')$}
	\loigiai{ 
		\begin{itemchoice}
			 \itemch Ta có
			$L=\lim  \dfrac{\sqrt{9n^2-n+1}}{4n-2}
			=\lim  \dfrac{\sqrt{9-\dfrac{1}{n}+\dfrac{1}{n^2}}}{4-\dfrac{2}{n}}
			=\dfrac{3}{4}$. 
			\itemch Ta có  \\
			$\begin{aligned}
			L'=&\lim \left(\sqrt{4n^2+n}-\sqrt{4n^2+2}\right) 
			=\lim  \dfrac{n-2}{\sqrt{4n^2+n}+\sqrt{4n^2+2}} \\
			=&\lim  \dfrac{1-\dfrac{2}{n}}{\sqrt{4+\dfrac{1}{n}}+\sqrt{4+\dfrac{2}{n^2}}}
			 =\dfrac{1}{4}\cdot
		\end{aligned}$\\
		 Suy ra $L+L'=1$.
		\itemch Ta  có
		$\lim  \left[f(n)\cdot g(n)\right]
		=\lim  f(n)\cdot \lim g(n)
		 =\dfrac{3}{4}\cdot \dfrac{1}{4}
		=\dfrac{3}{16}\ne  \dfrac{3}{8}\cdot$
		\itemch  Ta có \\
		$\begin{aligned}
		L''=\lim  \left(\sqrt{n^2+2n}-n+a\right)
		 =&\lim \left(\dfrac{2n}{\sqrt{n^2+2n}+n}+a\right)\\
		=&\lim  \left(\dfrac{2}{\sqrt{1+\dfrac{2}{n}}+1}+a\right)=1+a.
	\end{aligned}$\\
	Ta có  $L''=\dfrac{2}{a}\left(L+L'\right)
	\Leftrightarrow 1+a=\dfrac{2}{a} 
	\Leftrightarrow  a^2+a-2=0
	\Leftrightarrow  \hoac{&a=-2 \text{ (nhận)} \\&a=1 \text{ (nhận)}}$ nên có $2$ giá trị $a$.
	\end{itemchoice}
} 
\end{ex}

% nguồn: Ok_dungsai_lop11_Lan2
\begin{ex}%[1D3V1-2]%[ma:5ea4efb5]
	Cho giới hạn $M=\lim  \left(\sqrt{n^2+18n+1}\right)$, $N=\lim \left(\sqrt{4n^2+2n}-n\right)$,\break $P=\lim \left(\sqrt{n^2-3n+1}-n\right)$ và $Q=\lim \left(\sqrt{4n^2+3n}-\sqrt[3]{8n^3+n}\right)$.
	\choiceTF 
		 {\True $M=+\infty$}
		 {$N=1$}
		{\True $P=-\dfrac{3}{2}$} 
		{$Q=\dfrac{3}{2}$} 
	\loigiai{ 
		\begin{itemchoice}
			\itemch  Ta có
			$M=\lim \left(\sqrt{n^2+18n+1}\right)=\lim  n\left(\sqrt{1+\dfrac{18}{n}+\dfrac{1}{n^2}}\right)=+\infty$
			vì $\heva{&\lim n=+\infty\\ &\lim \sqrt{1+\dfrac{18}{n}+\dfrac{1}{n^2}}=1>0.}$ 
			\itemch  Ta có
			$N=\lim \left(\sqrt{4n^2+2n}-n\right)=\lim n\left(\sqrt{4+\dfrac{2}{n}}-1\right)=+\infty$ 
			vì $\heva{&\lim  n=+\infty\\ &\lim \left(\sqrt{4+\dfrac{2}{n}}-1\right)=1>0.}$
			\itemch Ta  có\\
			$\begin{aligned}
			P=&\lim  \left(\sqrt{n^2-3n+1}-n\right)
			=\lim  \dfrac{n^2-3n+1-n^2}{\sqrt{n^2-3n+1}+n}=\lim \dfrac{-3n+1}{\sqrt{n^2-3n+1}+n}\\
			=&\lim  \dfrac{-3+\dfrac{1}{n}}{\sqrt{1-\dfrac{3}{n}+\dfrac{1}{n^2}}+1}
			=\dfrac{-3}{2}\cdot 
		\end{aligned}$
		\itemch  Ta có\\
		$\begin{aligned}
		 Q=&\lim \left(\sqrt{4n^2+3n}-\sqrt[3]{8n^3+n}\right)=\lim \left(\sqrt{4n^2+3n}-2n\right)+\lim \left(2n-\sqrt[3]{8n^3+n}\right) \\
		=&\lim  \dfrac{4n^2+3n-4n^2}{\sqrt{4n^2+3n}+2n}+\lim \dfrac{8n^3-8n^3-n}{4n^2+2n\cdot \sqrt[3]{8n^3+n}+\sqrt[3]{\left(8n^3+n\right)^2}} \\
		=&\lim \dfrac{3n}{\sqrt{4n^2+3n}+2n}-\lim \dfrac{n}{4n^2+2n\cdot \sqrt[3]{8n^3+n}+\sqrt[3]{\left(8n^3+n\right)^2}} \\ 
		=&\lim \dfrac{3}{\sqrt{4+\dfrac{3}{n}}+2}-\lim \dfrac{\dfrac{1}{n}}{4+2\cdot \sqrt[3]{8+\dfrac{1}{n^2}}+\sqrt[3]{\left(8+\dfrac{1}{n^2}\right)^2}} \\ 
		 =&\dfrac{3}{4}-0=\dfrac{3}{4}\cdot
	\end{aligned}$
	\end{itemchoice}
} 
\end{ex}

% nguồn: Ok_dungsai_lop11_Lan2
\begin{ex}%[1D3V1-2]%[ma:2f9a3432]
	 Cho giới hạn $L=\lim \dfrac{5n^2-3an^4}{(1-a)n^4+2n+1}\cdot$
	 \choiceTF
		{\True  $a=0$ thì $L=0$}
		 {\True $a>2$ thì $L>0$}
		 {$a\in \left[0;1\right]$ thì $L>0$}
		{Có $5$ giá trị nguyên của $a\in \left[-10;10\right]$ để $L=\lim \dfrac{5n^2-3an^4}{(1-a)n^4+2n+1}>0$} 
	 \loigiai{
		\begin{itemchoice}
			\itemch Với $a=0$ ta  có $L=\lim \dfrac{5n^2}{n^4+2n+1}=\lim \dfrac{5}{n^2+\dfrac{2}{n^3}+\dfrac{1}{n^2}}=0$.
			\itemch Ta có $L=\lim  \dfrac{5n^2-3an^4}{(1-a)n^4+2n+1}=\lim \dfrac{\dfrac{5}{n^2}-3a}{(1-a)+\dfrac{2}{n^3}+\dfrac{1}{n^4}}=\dfrac{-3a}{(1-a)}>0\Leftrightarrow \hoac{&a<0 \\&a>1.}$
			\itemch Ta có 
			\begin{itemize}
				\item Với $a=0$  ta có $L=\lim \dfrac{5n^2}{n^4+2n+1}=\lim \dfrac{5}{n^2+\dfrac{2}{n^3}+\dfrac{1}{n^2}}=0$. \hfill $(1)$
				\item Với $a=1$ $\Rightarrow L=\lim  \dfrac{5n^2-3n^4}{2n+1}=\lim \dfrac{\dfrac{5}{n^2}-3}{\dfrac{2}{n^3}+\dfrac{1}{n^4}}=-\infty$. \hfill $(2)$
				\item Với $a\ne 1;  a\ne 0$. \\
				 Ta có $L=\lim \dfrac{5n^2-3an^4}{(1-a)n^4+2n+1}=\lim \dfrac{\dfrac{5}{n^2}-3a}{(1-a)+\dfrac{2}{n^3}+\dfrac{1}{n^4}}=\dfrac{-3a}{(1-a)}>0\Leftrightarrow \hoac{&a<0 \\&a>1.}$\hfill $(3)$
			\end{itemize}
			 \itemch Ta có
			\begin{itemize}
				 \item Với $a=0$ ta có $L=\lim \dfrac{5n^2}{n^4+2n+1}=\lim \dfrac{5}{n^2+\dfrac{2}{n^3}+\dfrac{1}{n^2}}=0$. \hfill $(1)$
				\item Với $a=1\Rightarrow  L=\lim \dfrac{5n^2-3n^4}{2n+1}=\lim \dfrac{\dfrac{5}{n^2}-3}{\dfrac{2}{n^3}+\dfrac{1}{n^4}}=-\infty$. \hfill $(2)$
				\item Với $a\ne 1; a\ne  0$. \\
				Ta có  $L=\lim \dfrac{5n^2-3an^4}{(1-a)n^4+2n+1}=\lim \dfrac{\dfrac{5}{n^2}-3a}{(1-a)+\dfrac{2}{n^3}+\dfrac{1}{n^4}}=\dfrac{-3a}{(1-a)}>0\Leftrightarrow \hoac{&a<0 \\&a>1.}$ \hfill $(3)$
			\end{itemize}
			Kết  hợp $a$ là số nguyên và $a\in \left[-10;10\right]$, suy ra có $19$ giá trị nguyên của $a$ thoả mãn $L>0$.
		\end{itemchoice}
	 }
\end{ex}


%%%%%%%%%%%%%%%% TRẢ LỜI NGẮN %%%%%%%%%%%%%%%%

% nguồn: Ok_TraLoiNgan_Lop11_Lan1
\begin{ex}%[1D3H1-2]%[ma:6538ef8e]
	Cho các giới hạn  $\lim\limits_{n \to +\infty} \dfrac{3n^2-2n+5}{4n^2+7}=a$ và $\lim\limits_{x \to-\infty} \dfrac{\sqrt{x^2+3x+5}}{4x-1}=b$.
	Giá trị của tổng $a+b$ bằng  bao nhiêu?
	 \shortans{0,5}
	 \loigiai{
		\begin{itemize}
			\item Ta có $\lim\limits_{n  \to +\infty} \dfrac{3n^2-2n+5}{4n^2+7}=\lim\limits_{n \to +\infty} \dfrac{3-\dfrac{2}{n}+\dfrac{2}{n^2}}{4+\dfrac{7}{n^2}}=\dfrac{3}{4}$.
			\item  Ta có
			\begin{eqnarray*}
				\lim\limits_{x  \to-\infty} \dfrac{\sqrt{x^2+3x+5}}{4x-1}
				&=  & \lim\limits_{x \to-\infty} \dfrac{\sqrt{x^2\left(1+\dfrac{3}{x}+\dfrac{5}{x^2}\right)}}{x\left(4-\dfrac{1}{x}\right)}\\
				&= & \lim\limits_{x \to-\infty}  \dfrac{\left| x\right| \sqrt{1+\dfrac{3}{x}+\dfrac{5}{x^2}}}{x\left(4-\dfrac{1}{x}\right)}\\
				 &= & \lim\limits_{x \to-\infty} \dfrac{ -\sqrt{1+\dfrac{3}{x}+\dfrac{5}{x^2}}}{4-\dfrac{1}{x}}\\
				&=  &-\dfrac{1}{4}.
			\end{eqnarray*}
		\end{itemize}
		 Suy ra $a=\dfrac{3}{4}$ và $b=-\dfrac{1}{4}$.
		Vậy  $a+b=\dfrac{1}{2}=0{,}5$.
	} 
\end{ex}

% nguồn: Ok_TraLoiNgan_Lop11_Lan1
\begin{ex}%[1D3H1-2]%[ma:f465edd6]
	Tính giới hạn sau ${\lim \left(1+3n-\sqrt{9n^2-n+7}\right)}$  (làm tròn đến hàng phần trăm).
	\shortans{1,17} 
	\loigiai{ 
		$\begin{aligned}
		& \lim \left(1+3n-\sqrt{9{{n}^2}-n+7} \right) =\lim  \left(1+\dfrac{9{{n}^2}-9{{n}^2}+n-7}{3n+\sqrt{9{{n}^2}-n+7}} \right)\\
		&=\lim \left(1+\dfrac{n-7}{3n+\sqrt{9{{n}^2}-n+7}} \right)=\lim \left(1+\dfrac{1-\dfrac{7}{n}}{3+\sqrt{9-\dfrac{1}{n}+\dfrac{7}{{{n}^2}}}} \right)=\dfrac{7}{6}. 
	\end{aligned}$
} 
\end{ex}

% nguồn: Ok_TraLoiNgan_Lop11_Lan1
\begin{ex}%[1D3H1-2]%[ma:de1087df]
	 Tìm giới hạn sau: ${\lim \dfrac{1+2+\cdots+n}{n^2+3n}}$.
	 \shortans{0,5}
	 \loigiai{
		${1+2+\cdots+n}$ là tổng của ${n}$ số hạng đầu của một cấp số cộng có số hạng  đầu ${u_1=1}$ và công sai ${d=1}$. Vì vậy ${1+2+\cdots+n=\dfrac{(1+n) n}{2}}$.\\
		Ta có ${\lim \dfrac{1+2+\cdots+n}{n^2+3n}=\lim \dfrac{n(n+1)}{2\left(n^2+3n\right)}=\lim \dfrac{n^2\left(1+\dfrac{1}{n}\right)}{2n^2\left(1+\dfrac{3}{n}\right)}=\lim \dfrac{1+\dfrac{1}{n}}{2\left(1+\dfrac{3}{n}\right)}=\dfrac{1}{2}}$. 
	 }
\end{ex}

% nguồn: Ok_TraLoiNgan_Lop11_Lan1
\begin{ex}%[1D3H1-2]%[ma:8c98ba96]
	 Tìm giới hạn sau: ${\lim n \sqrt{\dfrac{9n^2}{4n^4-n^2}}}$.
	\shortans{1,5} 
	 \loigiai{
		Ta có  ${\lim n \sqrt{\dfrac{9n^2}{4n^4-n^2}}=\lim \dfrac{n^2}{n^2} \sqrt{\dfrac{9}{4-\dfrac{1}{n^2}}}=\lim \sqrt{\dfrac{9}{4-\dfrac{1}{n^2}}}=\dfrac{3}{2}}=1{,}5$.
	 }
\end{ex}

% nguồn: Ok_TraLoiNgan_Lop11_Lan1
\begin{ex}%[1D3H1-2]%[ma:910cb712]
	Tìm giới hạn sau: ${\lim  \dfrac{n \cdot \sqrt{1+3+5+\cdots+(2n-1)}}{2n^2+n+1}}$.
	\shortans{0,5} 
	 \loigiai{
		Nhận xét: ${1+3+5+\cdots+(2n-1)}$ là  tổng của ${n}$ số hạng đầu tiên của một cấp số cộng có số hạng đầu là ${u_1=1}$ và công sai là ${d=2}$. Vì vậy, ta có
		$$1+3+5+\cdots+(2 n-1)=\dfrac{\left(u_1+u_n\right)  n}{2}=\dfrac{(1+2 n-1) n}{2}=n^2.
		$$ 
		Vì  vậy
		$$ 
		\begin{aligned}
			&\lim \dfrac{n \sqrt{1+3+5+\cdots+(2n-1)}}{2 n^2+n+1}=\lim \dfrac{n  \sqrt{n^2}}{2 n^2+n+1}=\lim \dfrac{n^2}{2 n^2+n+1}\\
			&=\lim \dfrac{n^2}{n^2\left(2+\dfrac{1}{n}+\dfrac{1}{n^2}\right)}=\lim  \dfrac{1}{2+\dfrac{1}{n}+\dfrac{1}{n^2}}=\dfrac{1}{2}.
		\end{aligned}
		 $$
	 }
\end{ex}

% nguồn: Ok_TraLoiNgan_Lop11_Lan1
\begin{ex}%[1D3H1-2]%[ma:41bce925]
	Tìm giới hạn  $\lim\limits_{n \to+\infty} \dfrac{-2n^2-3}{1-n^2}$ bằng
	\shortans{2} 
	 \loigiai{
		Ta có $\lim\limits_{n \to+\infty} \dfrac{-2n^2-3}{1-n^2}=\lim\limits_{n \to+\infty}  \dfrac{-2-\dfrac{3}{n^2}}{\dfrac{1}{n^2}-1}=\dfrac{-2-0}{0-1}=2$.
	} 
\end{ex}

% nguồn: Ok_TraLoiNgan_Lop11_Lan2
\begin{ex}%[1D3H1-2]%[ma:eec472d5]
	Giới hạn  $\lim \dfrac{2n^2-3\sqrt{n}+1}{3n \sqrt{n}+2n}=\lim \dfrac{a \sqrt{n}-\dfrac{3}{n}+\dfrac{1}{n \sqrt{n}}}{b+\dfrac{2}{\sqrt{n}}}$ với $a$; $b$ là các số tự nhiên. Tính $P=a+b^2$.
	\shortans{11} 
	 \loigiai{
		 Ta có
		\begin{eqnarray*}
			\lim \dfrac{2n^2-3\sqrt{n}+1}{3n \sqrt{n}+2n}&=&\lim \dfrac{n \sqrt{n}\left(2\sqrt{n}-\dfrac{3}{n}+\dfrac{1}{n  \sqrt{n}}\right)}{n \sqrt{n}\left(3+\dfrac{2}{\sqrt{n}}\right)}\\
			&=&\lim \dfrac{2\sqrt{n}-\dfrac{3}{n}+\dfrac{1}{n  \sqrt{n}}}{3+\dfrac{2}{\sqrt{n}}}.
		\end{eqnarray*}
		Suy ra $a=2$, $b=3$.\\ 
		Vậy  $P=a+b^2=11$.
	 }
\end{ex}

% nguồn: Ok_TraLoiNgan_Lop11_Lan2
\begin{ex}%[1D3H1-2]%[ma:0bbcbd87]
	Tìm giới  hạn $\lim \left[\dfrac{1}{1\cdot3}+\dfrac{1}{3\cdot5}+\dfrac{1}{5\cdot7}+\ldots+\dfrac{1}{(2n-1)(2n+1)}\right]$.
	 \shortans{0,5}
	 \loigiai{
		Ta có 
		\begin{eqnarray*}
			 \lim \left[\dfrac{1}{1 \cdot 3}+\dfrac{1}{3 \cdot 5}+\dfrac{1}{5 \cdot 7}+\ldots+\dfrac{1}{(2 n-1)(2 n+1)}\right]&=&\lim \dfrac{1}{2}\left(1-\dfrac{1}{3}+\dfrac{1}{3}-\dfrac{1}{5}+\ldots+\dfrac{1}{2 n-1}-\dfrac{1}{2 n+1}\right)\\
			 &=&\lim \dfrac{1}{2}\left(1-\dfrac{1}{2 n+1}\right)\\
			&=&\dfrac{1}{2}. 
		\end{eqnarray*}
	} 
\end{ex}

% nguồn: Ok_TraLoiNgan_Lop11_Lan2
\begin{ex}%[1D3H1-2]%[ma:f5669174]
	Tính  giới hạn $L=\lim \left[\dfrac{\left(3^n+2\right)\left(2^n+1\right)}{3^n+2}-2^n+2\right]$.
	 \shortans{3}
	\loigiai{ 
		Ta có 
		\begin{eqnarray*}
			L=\lim \left[\dfrac{\left(3^n+2\right)\left(2^n+1\right)}{3^n+2}-2^n+2\right]&=&\lim \left[\dfrac{\left(3^n+2\right)\left(2^n+1\right)-2^n\left(3^n+2\right)}{3^n+2}\right]+2\\ 
			 &=&\lim \dfrac{6^n+3^n+2\cdot 2^n+2-6^n-2\cdot 2^n}{3^n+2}+2\\
			&=&\lim  \dfrac{3^n+2}{3^n+2}+2\\
			&=&\lim 1+2\\ 
			 &=&3.
		\end{eqnarray*}
	} 
\end{ex}

% nguồn: Ok_TraLoiNgan_Lop11_Lan2
\begin{ex}%[1D3H1-2]%[ma:96a06a16]
	 Biết giới hạn $\lim \dfrac{\sqrt[3]{2n^2-n^3}+n}{\sqrt{n^2+n}-n}=\lim \dfrac{\sqrt{1+\dfrac{1}{n}}+1}{\sqrt[3]{\left(\dfrac{a}{n}-b\right)^2}+1-\sqrt[3]{\dfrac{2}{n}-1}}$ với $a$; $b$ là các số nguyên dương. Tính $S=a+b$.
	 \shortans{3}
	\loigiai{ 
		 Ta có
		\begin{eqnarray*}
			\lim \dfrac{\sqrt[3]{2n^2-n^3}+n}{\sqrt{n^2+n}-n}&=&\lim  \left( \dfrac{2n^2-n^3+n^3}{n^2+n-n^2} \cdot \dfrac{\sqrt{n^2+n}+n}{\sqrt[3]{\left(2n^2-n^3\right)^2}+n^2-n \sqrt[3]{2n^2-n^3}}\right) \\
			 &=&\lim \dfrac{\sqrt{\left(n \sqrt{1+\dfrac{1}{n}}+n\right)}}{\sqrt[3]{n^6\cdot\left(\dfrac{2}{n}-1\right)^2}+n^2-n \cdot \sqrt[3]{n^3\left(\dfrac{2}{n}-1\right)}}\\
			&=&\lim \dfrac{\sqrt{1+\dfrac{1}{n}}+1}{\left(\dfrac{2}{n}-1\right)^{\tfrac{2}{3}}+1-\sqrt[3]{\dfrac{2}{n}-1}}\\ 
			&=&\lim  \dfrac{\sqrt{1+\dfrac{1}{n}}+1}{\sqrt[3]{\left(\dfrac{2}{n}-1\right)^2}+1-\sqrt[3]{\dfrac{2}{n}-1}}.
		\end{eqnarray*}
		Do đó, $\heva{&a=2\\&b=1} \Rightarrow  S=3$.
	} 
\end{ex}

% nguồn: Ok_TraLoiNgan_Lop11_Lan2
\begin{ex}%[1D3H1-2]%[ma:b6e24151]
	Giới hạn $\lim \sqrt[3]{\dfrac{2^{3n}-3^{n+1}}{8^{n+1}+6^{n-1}}}=\lim \sqrt[3]{\dfrac{1-3\cdot a^n}{8+\dfrac{1}{6} \cdot b^n}}$ với $a$; $b$ là các số hữu tỉ. Tính giá  trị $8(a+b)$.
	\shortans{9} 
	\loigiai{ 
		Ta  có
		\begin{eqnarray*}
			\lim \sqrt[3]{\dfrac{2^{3n}-3^{n+1}}{8^{n+1}+6^{n-1}}}&=&\lim \sqrt[3]{\dfrac{8^n-3.3^n}{8.8^n+\dfrac{1}{6}  6^n}}\\
			 &=&\lim{\sqrt[3]{\dfrac{8^n\left(1-3\cdot\left(\dfrac{3}{8}\right)^n\right)}{8^n\left(8+\dfrac{1}{6}\left(\dfrac{6}{8}\right)^n\right)}}}\\
			&=&\lim{\sqrt[3]{\dfrac{1-3\left(\dfrac{3}{8}\right)^n}{8+\dfrac{1}{6}\left(\dfrac{3}{4}\right)^n}}}. 
		\end{eqnarray*}
		Suy ra $a=\dfrac{3}{8}$ và $b=\dfrac{3}{4}$.  \\
		 Vậy $8(a+b)=9$.
	 }
\end{ex}

% nguồn: Ok_TraLoiNgan_Lop11_Lan2
\begin{ex}%[1D3H1-2]%[ma:002eeae0]
	 Tính giới hạn $L=\lim \dfrac{1+3+5+\ldots+2n-1}{5n^2-n+1}$.
	\shortans{0,2} 
	 \loigiai{
		Ta có tổng $1+3+5+\ldots+2n-1$ là tổng cấp  số cộng của $n$ số hạng với $u_1=1$, $d=2$.\\
		Suy  ra $1+3+5+\cdots+2n-1=\dfrac{n(1+2n-1)}{2}=n^2$.\\
		Do  đó
		\begin{eqnarray*}
			L&=&\lim  \dfrac{1+3+5+\ldots+2n-1}{5n^2-n+1}\\
			&=&\lim  \dfrac{n^2}{5n^2-n+1}\\
			 &=&\lim \dfrac{1}{5-\dfrac{1}{n}+\dfrac{1}{n^2}}\\
			 &=&\dfrac{1}{5}.
		\end{eqnarray*}
		Vậy  $L=0{,}2$.
	 }
\end{ex}

% nguồn: Ok_TraLoiNgan_Lop11_Lan2
\begin{ex}%[1D3H1-2]%[ma:d012416e]
	Cho dãy số $(u_n)$ với $u_n=\dfrac{3n-1}{n \cdot 2^n}+2024$. Khi đó $\lim\limits_{n \to+\infty} u_n$ bằng bao  nhiêu?
	 \shortans{2024}
	\loigiai{ 
		Ta  có
		\begin{eqnarray*}
			\lim\limits_{n \to+\infty}  u_n&=&\lim\limits_{n \to+\infty}\left(\dfrac{3n-1}{n \cdot 2^n}+2024\right)\\
			&=&\lim\limits_{n  \to+\infty}\left(\dfrac{3-\dfrac{1}{n}}{2^n}+2024\right)\\
			 &=&2024.
		\end{eqnarray*}
		Vì $\lim\limits_{n \to+\infty}\left(3-\dfrac{1}{n}\right)=3$ và $\lim\limits_{n \to+\infty} 2^n=+\infty$ nên $\lim\limits_{n \to+\infty}  \dfrac{3-\dfrac{1}{n}}{2^n}=0$.\\
		Vậy $\lim\limits_{n \to+\infty} u_n=0+2024=2024$. 
	 }
\end{ex}

% nguồn: Ok_TraLoiNgan_Lop11_Lan2
\begin{ex}%[1D3H1-2]%[ma:9fbce122]
	Cho $\lim\limits_{n \to+\infty}\left[\sqrt{4n^2-3n+1}-(a n+b)\right]=0$, ($a$,  $b \in \mathbb{R}$). Tính $S=a+8b$.
	\shortans{-4} 
	\loigiai{ 
		Nhận xét: $\lim\limits_{n \to+\infty}\left[\sqrt{4n^2-3n+1}-(a n+b)\right]=0$  nên $a=\sqrt{4}=2$.\\
		 Khi đó
		\begin{eqnarray*}
			\lim\limits_{n \to+\infty}\left[\sqrt{4n^2-3n+1}-(an+b)\right]&=&\lim\limits_{n  \to+\infty}\left[\sqrt{4n^2-3n+1}-(2n+b)\right]\\
			&=&  \lim\limits_{n \to+\infty}\left(\dfrac{4n^2-3n+1-4n^2}{\sqrt{4n^2-3n+1}+2n}-b\right)\\
			 &=& \lim\limits_{n \to+\infty}\left[\dfrac{-3n+1}{\sqrt{4n^2-3n+1}+2n}-b\right]\\
			&=&\lim\limits_{n \to+\infty}\left[\dfrac{n\left( -3+\dfrac{1}{n}\right) }{n\left(\sqrt{4-\dfrac{3}{n}+\dfrac{1}{n^2}}+2\right)  }-b\right]\\
			 &=&\lim\limits_{n \to+\infty}\left[\dfrac{-3+\dfrac{1}{n}}{\sqrt{4-\dfrac{3}{n}+\dfrac{1}{n^2}}+2}-b\right]\\
			&=&\dfrac{-3}{4}-b. 
		\end{eqnarray*}
		Theo đề $\lim\limits_{n  \to+\infty}\left[\sqrt{4n^2-3n+1}-(a n+b)\right]=0$ nên $\dfrac{-3}{4}-b$ hay $b=\dfrac{-3}{4}$.\\
		 Vậy $S=a+8b=-4$.
	} 
\end{ex}

% nguồn: Ok_TraLoiNgan_Lop11_Lan2
\begin{ex}%[1D3H1-2]%[ma:e9a95041]
	Giá trị của giới hạn $\lim \left(\sqrt{n^2+2n-2}-n\right)=\lim \dfrac{a \cdot\left(1-\dfrac{1}{n}\right)}{\sqrt{1+b  \cdot\left(\dfrac{1}{n}-\dfrac{1}{n^2}\right)}+1}$ với $a$; $b$ là các số tự nhiên. Tính $S=(a+b)^2$.
	\shortans{16} 
	\loigiai{ 
		Ta có 
		\begin{eqnarray*}
			\lim  \left(\sqrt{n^2+2n-2}-n\right)&=&\lim \dfrac{n^2+2n-2-n^2}{\sqrt{n^2+2n-2}+n}\\
			&=&\lim  \dfrac{2n-2}{n \sqrt{1+\dfrac{2}{n}-\dfrac{2}{n^2}}+n}\\
			&=&\lim  \dfrac{n\left(2-\dfrac{2}{n}\right)}{n\left(\sqrt{1+\dfrac{2}{n}-\dfrac{2}{n^2}}+1\right)}\\
			&=&\lim \dfrac{2-\dfrac{2}{n}}{\sqrt{1+\dfrac{2}{n}-\dfrac{2}{n^2}}+1}\\ 
			&=&\lim \dfrac{2\cdot\left(1-\dfrac{1}{n}\right)}{\sqrt{1+2\cdot\left(\dfrac{1}{n}-\dfrac{1}{n^2}\right)}+1}. 
		\end{eqnarray*}
		Khi đó  $a=b=2\Rightarrow S=(a+b)^2=16$.
	} 
\end{ex}

% nguồn: Ok_TraLoiNgan_Lop11_Lan2
\begin{ex}%[1D3H1-2]%[ma:2a687563]
	$\lim \dfrac{n \cdot \sqrt{n^2+1}}{-2n^2+n+1}$ bằng bao nhiêu (kết quả viết dưới dạng số thập phân)? 
	 \shortans{-0,5}
	 \loigiai{
		 Ta có
		\begin{eqnarray*}
			\lim \dfrac{n \sqrt{n^2+1}}{-2 n^2+n+1}&=&\lim  \dfrac{n^2 \sqrt{1+\dfrac{1}{n^2}}}{n^2\left(-2+\dfrac{1}{n}+\dfrac{1}{n^2}\right)}\\
			&=&\lim \dfrac{\sqrt{1+\dfrac{1}{n^2}}}{-2+\dfrac{1}{n}+\dfrac{1}{n^2}}\\ 
			 &=&-\dfrac{1}{2}.
		\end{eqnarray*}
		 Vậy $\lim \dfrac{n \cdot \sqrt{n^2+1}}{-2n^2+n+1}=-0{,}5$.
	} 
\end{ex}

% nguồn: Ok_TraLoiNgan_Lop11_Lan1
\begin{ex}%[1D3V1-2]%[ma:19e77190]
	Tìm giới hạn sau: ${\lim \dfrac{1+2+2^2+\cdots+2^n}{1+3+3^2+\cdots+3^n}}$. 
	 \shortans{0}
	 \loigiai{
		Xét ${1+2+2^2+\cdots+2^n}$ là tổng của ${n+1}$ số hạng đầu của một cấp số nhân có số hạng đầu ${u_1=1}$, công bội bằng  ${q=2}$ nên
		 $$1+2+2^2+\cdots+2^n=\dfrac{u_1\left(1-q^{n+1}\right)}{1-q}=\dfrac{1\left(1-2^{n+1}\right)}{1-2}=2^{n+1}-1.$$
		Hoàn toàn  tương tự, ta có: ${1+3+3^2+\cdots+3^n=\dfrac{1\left(1-3^{n+1}\right)}{1-3}=\dfrac{3^{n+1}-1}{2}}$.\\
		 Khi đó
		$$ 
		\begin{aligned}
			\lim \dfrac{1+2+2^2+\cdots+2^n}{1+3+3^2+\cdots+3^n}&=\lim \dfrac{2^{n+1}-1}{\dfrac{3^{n+1}-1}{2}}=\lim  2 \cdot \dfrac{2 \cdot 2^n-1}{3 \cdot 3^n-1}\\
			& =\lim 2 \cdot \dfrac{2 \cdot\left(\dfrac{2}{3}\right)^n-\dfrac{1}{3^n}}{3-\dfrac{1}{3^n}}=2  \cdot \dfrac0{3}=0.
		\end{aligned}
		$$ 
	} 
\end{ex}

% nguồn: Ok_TraLoiNgan_Lop11_Lan1
\begin{ex}%[1D3V1-2]%[ma:214eb82e]
	Cho  $a$ thỏa $\lim \dfrac{20 a n^2+a^3 n+2023}{(2 n-1)^2}=a^2-a+9$. Giá trị của $a$ là
	 \shortans{3}
	\loigiai{Ta  có \\
		 $\lim \dfrac{20 a n^2+a^3 n+2023}{(2 n-1)^2}=a^2-a+9 \Leftrightarrow \lim \dfrac{n^2\left(20 a+\dfrac{a^3}{n}+\dfrac{2023}{n^2}\right)}{n^2\left(2-\dfrac{1}{n}\right)^2}=a^2-a+9$
		 $$
		\begin{aligned}
			 & \Leftrightarrow \lim \dfrac{(20 a+0+0)}{(2-0)^2}=a^2-a+9 \Leftrightarrow \lim \dfrac{20 a}{4}=a^2-a+9 \Leftrightarrow 5 a=a^2-a+9 \\
			 & \Leftrightarrow a^2-6 a+9=0 \Leftrightarrow a=3
		\end{aligned}
		 $$
		Vậy $a=3$ thỏa mãn  yêu cầu bài toán.
	} 
\end{ex}

% nguồn: Ok_TraLoiNgan_Lop11_Lan2
\begin{ex}%[1D3V1-2]%[ma:c448007b]
	Biết 
	 \[
	\lim \dfrac{2n^3 +  n^2 - 4}{an^3 + 2} = \dfrac{1}{2}
	 \]
	 với $a$ là tham số. Khi đó $a - a^2$ bằng bao nhiêu?
	 \shortans{-12}
	\loigiai{ 
		 \[
		\lim \dfrac{2n^3 + n^2 - 4}{an^3 + 2} =\lim \dfrac{2  + \dfrac{1}{n}-\dfrac{4}{n^3}}{a+\dfrac{2}{n}} = \dfrac{2}{a}
		\] 
		Từ điều kiện đầu bài suy ra $\dfrac{2}{a} = \dfrac{1}{2} \Rightarrow a = 4$.  \\Vậy $a - a^2 = -12$.
	} 
\end{ex}

% nguồn: Ok_TraLoiNgan_Lop11_Lan2
\begin{ex}%[1D3V1-2]%[ma:c97f023a]
	Gọi  $S$ là tập hợp các tham số nguyên $a$ thỏa mãn $\lim \left(\dfrac{3n+2}{n+2}+a^2-4a\right)=0$. Tổng các phần tử của $S$ bằng
	\shortans{4} 
	\loigiai{ 
		 Ta có
		$$\begin{aligned}
		\lim  \left(\dfrac{3n+2}{n+2}+a^2-4a\right)
		&=\lim \left(\dfrac{\left(a^2-4a+3\right)n+2+2a^2-8a}{n+2}\right)\\ 
		&=\lim \left(\dfrac{a^2-4a+3+\dfrac{2+2a^2-8a}{n}}{1+\dfrac{2}{n}}\right) 
		 =a^2-4a+3.
	\end{aligned}$$
	 Theo giả thiết ta có
	$$\begin{aligned}
	 &\lim \left(\dfrac{3n+2}{n+2}+a^2-4a\right)=0
	\Leftrightarrow  a^2-4a+3=0
	\Leftrightarrow \hoac{&a=3\\ &a=1.} 
	\end{aligned}$$
	 Vậy $S=\left\{1;3\right\}\Rightarrow 1+3=4$.
} 
\end{ex}

% nguồn: Ok_TraLoiNgan_Lop11_Lan2
\begin{ex}%[1D3V1-2]%[ma:2f296139]
	Tính giới  hạn $\lim \left[\dfrac{1}{1\cdot 2}+\dfrac{1}{2\cdot 3}+\dfrac{1}{3\cdot 4}+ \cdots +\dfrac{1}{n(n+1)}\right]$.
	 \shortans{1}
	 \loigiai{
		Ta có $\dfrac{1}{1\cdot 2}+\dfrac{1}{2\cdot 3}+\dfrac{1}{3\cdot  4}+ \cdots +\dfrac{1}{n(n+1)}
		 =\dfrac{1}{1}-\dfrac{1}{2}+\dfrac{1}{2}-\dfrac{1}{3}+\cdots +\dfrac{1}{n-1}-\dfrac{1}{n}+\dfrac{1}{n}-\dfrac{1}{n+1}
		 =1-\dfrac{1}{n+1}$.\\
		Vậy  $\lim \left[\dfrac{1}{1\cdot 2}+\dfrac{1}{2\cdot 3}+\dfrac{1}{3\cdot 4}+ \cdots +\dfrac{1}{n(n+1)}\right]
		 =\lim \left(1-\dfrac{1}{n+1}\right)=1$.
	 }
\end{ex}

% nguồn: Ok_TraLoiNgan_Lop11_Lan2
\begin{ex}%[1D3V1-2]%[ma:dc8ef7d1]
	Cho giới  hạn $\lim \dfrac{(a+2)\cdot 2^{2n-2}-3\cdot 4^n-5}{2\cdot 3^{n+3}-4^{n-1}}=2$. Tìm $a$.
	 \shortans{8}
	 \loigiai{
		Ta có 
		$$\begin{aligned}
		&\lim  \dfrac{(a+2)\cdot 2^{2n-2}-3\cdot 4^n-5}{2\cdot 3^{n+3}-4^{n-1}}=2 \\
		\Leftrightarrow &\lim \dfrac{(a+2)\cdot  \dfrac{1}{4}\cdot 4^n-3\cdot 4^n-5}{2\cdot 27\cdot 3^n-\dfrac{1}{4}\cdot 4^n}=2 \\
		\Leftrightarrow &\lim  \dfrac{\dfrac{a+2}{4}-3-\dfrac{5}{4^n}}{54\left(\dfrac{3}{4}\right)^n-\dfrac{1}{4}}=2 \\
		\Leftrightarrow &-a+10=2 \Leftrightarrow  a=8.
	\end{aligned}$$
	Vậy $a=8$ là  giá trị cần tìm.
} 
\end{ex}

% nguồn: Ok_TraLoiNgan_Lop11_Lan2
\begin{ex}%[1D3V1-2]%[ma:a7852743]
	Cho dãy số $(u_n)$ với $u_n=\dfrac{1}{1\cdot 3}+\dfrac{1}{3\cdot 5}+ \cdots +\dfrac{1}{(2n-1)\cdot (2n+1)}$.  Tính $\lim{u_n}$.\\
	\begin{flushright}
	 \shortans{0,5}
	\end{flushright}
	 \loigiai{
		 Ta có
		$$\begin{aligned}
		u_n&=\dfrac{1}{1\cdot 3}+\dfrac{1}{3\cdot 5}+ \cdots  +\dfrac{1}{(2n-1)\cdot (2n+1)} \\
		&=\dfrac{1}{2}\left(\dfrac{1}{1}-\dfrac{1}{3}+\dfrac{1}{3}-\dfrac{1}{5}+ \cdots  +\dfrac{1}{2n-1}-\dfrac{1}{2n+1}\right) \\
		 &=\dfrac{1}{2}\left(\dfrac{1}{1}-\dfrac{1}{2n+1}\right)=\dfrac{n}{2n+1}.
	\end{aligned}$$
	Suy ra $\lim{u_n}=\lim \dfrac{n}{2n+1}=\dfrac{1}{2}$. 
} 
\end{ex}

% nguồn: Ok_TraLoiNgan_Lop11_Lan2
\begin{ex}%[1D3V1-2]%[ma:953cb775]
	Tìm giới hạn sau $\lim \dfrac{1+2+2^2+ \cdots +2^n}{1+3+3^2+  \cdots +3^n}$.
	\shortans{0} 
	\loigiai{ 
		Ta có \\ 
		$$\begin{aligned}
		\lim \dfrac{1+2+2^2+ \cdots +2^n}{1+3+3^2+ \cdots  +3^n}
		= &\lim  \dfrac{\dfrac{1-2^{n+1}}{-1}}{\dfrac{1-3^{n+1}}{-2}}
		= \lim  \dfrac{(1-2^{n+1})\cdot 2}{1-3^{n+1}}\\
		= &\lim \dfrac{\left[\left(\dfrac{1}{3}\right)^{n+1}-\left(\dfrac{2}{3}\right)^{n+1}\right]\cdot 2}{\left(\dfrac{1}{3}\right)^{n+1}-1}=0. 
	\end{aligned}$$
 }
\end{ex}

% nguồn: Ok_TraLoiNgan_Lop11_Lan2
\begin{ex}%[1D3V1-2]%[ma:9343b22e]
	Tìm  $\lim{u_n}$ biết $u_n=\dfrac{1}{2^2-1}+\dfrac{1}{3^2-1}+ \cdots +\dfrac{1}{n^2-1}$.
	 \shortans{0,75}
	 \loigiai{
		Ta có 
		$$\begin{aligned}
		u_n&=\dfrac{1}{2^2-1}+\dfrac{1}{3^2-1}+ \cdots  +\dfrac{1}{n^2-1} \\
		&=\dfrac{1}{1\cdot  3}+\dfrac{1}{2\cdot 4}+\dfrac{1}{3\cdot 5}+ \cdots +\dfrac{1}{(n-1)(n+1)}\\
		&=\dfrac{1}{2}\left(\dfrac{1}{1}-\dfrac{1}{3}+\dfrac{1}{2}-\dfrac{1}{4}+\dfrac{1}{3}-\dfrac{1}{5}+ \cdots  +\dfrac{1}{n-1}-\dfrac{1}{n+1}\right) \\
		 &=\dfrac{1}{2}\left(\dfrac{1}{1}+\dfrac{1}{2}-\dfrac{1}{n+1}\right)=\dfrac{3}{4}-\dfrac{1}{2(n+1)}.
	\end{aligned}$$
	Suy ra $\lim{u_n}=\lim  \left[\dfrac{3}{4}-\dfrac{1}{2(n+1)}\right]=\dfrac{3}{4}$.
} 
\end{ex}

% nguồn: Ok_TraLoiNgan_Lop11_Lan2
\begin{ex}%[1D3V1-2]%[ma:90b62996]
	Giới hạn $L=\lim \dfrac{1^2+2^2+3^2+  \cdots +n^2}{n^3}$ bằng
	 \shortans{0,33}
	 \loigiai{
		Ta  có $1^2+2^2+3^2+ \cdots +n^2=\dfrac{n(n+1)(2n+1)}{6}(*)$, thật vậy với $(*)$ đúng với $n=1$.\\
		Giả sử $(*)$ đúng  với $n=k$, ta có
		$$\begin{aligned}
		1^2+2^2+3^2+  \cdots +k^2+(k+1)^2
		&=\dfrac{k(k+1)(2k+1)}{6}+(k+1)^2 \\ 
		&=\dfrac{(k+1)(k+2)\left[2(k+1)+1\right]}{6}\cdot 
	\end{aligned}$$
	Vậy $(*)$ đúng  với $n\in \mathbb{N}*$.\\
	Do đó 
	\begin{eqnarray*}
		L&=&\lim \dfrac{1^2+2^2+3^2+ \cdots +n^2}{n^3}=\lim \dfrac{\dfrac{n(n+1)(2n+1)}{6}}{n^3}  \\
		&=&\lim \dfrac{n^3\left(1+\dfrac{1}{n}\right)\left(2+\dfrac{1}{n}\right)}{6n^3}  =\lim \dfrac{\left(1+\dfrac{1}{n}\right)\left(2+\dfrac{1}{n}\right)}{6}
		 =\dfrac{1}{3} \approx 0{,}33.
	\end{eqnarray*}
} 
\end{ex}

% nguồn: Ok_TraLoiNgan_Lop11_Lan2
\begin{ex}%[1D3V1-2]%[ma:6ad95ffe]
	Cho  dãy số $\left(u_n\right)$ với $\heva{&u_1=-2 \\&u_n=3u_{n-1}-1, \forall n \geq 2}$. Tính $\lim\dfrac{u_n}{3^n}$. (Kết quả làm tròn đến hàng phần mười)
	 \shortans{-0,8}
	\loigiai{ 
		Ta có dãy số $\left(u_n\right)$ với $\heva{&u_1=-2\\&u_n-\dfrac{1}{2}=3\left(u_{n-1}-\dfrac{1}{2}\right)}$, $\forall n \geq  2$.\\
		Đặt $v_n=u_n-\dfrac{1}{2} \Rightarrow v_{n-1}=u_{n-1}-\dfrac{1}{2}$, $\forall  n \geq 2$ và $v_1=u_1-\dfrac{1}{2}=-2-\dfrac{1}{2}=-\dfrac{5}{2}$.\\
		Khi đó dãy $\left(v_n\right)$ thoả mãn  $\heva{&v_1=-\dfrac{5}{2} \\&v_n=3v_{n-1}, \forall n \geq 2.}$\\
		Ta thấy $\left(v_n\right)$ là cấp số nhân với $v_1=-\dfrac{5}{2}$, công bội $q=3$, suy ra $v_n=-\dfrac{5}{2} \cdot 3^{n-1}$, $\forall n \geq  1$.\\
		Do đó $u_n=v_n+\dfrac{1}{2}=-\dfrac{5}{2} \cdot 3^{n-1}+\dfrac{1}{2}, \forall n \geq 1$.\\ 
		 Vậy
		\begin{eqnarray*}
			\lim \dfrac{u_n}{3^n}&=&\lim \dfrac{-\dfrac{5}{2} \cdot  3^{n-1}+\dfrac{1}{2}}{3^n}\\
			 &=&\lim \left(-\dfrac{5}{6}+\dfrac{1}{2\cdot 3^n}\right)\\
			&=&-\dfrac{5}{6}  \approx-0{,}8.
		\end{eqnarray*}
	} 
\end{ex}

% nguồn: Ok_TraLoiNgan_Lop11_Lan2
\begin{ex}%[1D3V1-2]%[ma:242154d1]
	Cho dãy số  $(u_n)$ như sau $u_n=\dfrac{n}{1+n^2+n^4}$, $\forall n\in \mathbb{N}^*$. Tính giới hạn $\lim\, \left(u_1+u_2+ \cdots +u_n\right)$.\\
	\begin{flushright}
	\shortans{0,5} 
	\end{flushright}
	 \loigiai{
		Ta có $u_n=\dfrac{n}{(1+n^2)^2-n^2}=\dfrac{n}{(n^2+n+1)(n^2-n+1)} 
		 =\dfrac{1}{2}\left(\dfrac{1}{n^2-n+1}-\dfrac{1}{n^2+n+1}\right)$. \\
		Ta  có
		$$\begin{aligned}
		u_1+u_2+ \cdots  +u_n
		&=\dfrac{1}{2}\left(1-\dfrac{1}{3}+\dfrac{1}{3}-\dfrac{1}{7}+\dfrac{1}{7}-\dfrac{1}{13}+\dfrac{1}{13}-\dfrac{1}{21}+  \cdots +\dfrac{1}{n^2-n+1}-\dfrac{1}{n^2+n+1}\right) \\
		&=\dfrac{1}{2}\left(1-\dfrac{1}{n^2+n+1}\right) 
		 =\dfrac{1}{2}\dfrac{n^2+n}{n^2+n+1}.
	\end{aligned}$$
	Suy  ra $\lim (u_1+u_2+ \cdots +u_n)=\dfrac{1}{2}\lim \dfrac{1+\dfrac{1}{n}}{1+\dfrac{1}{n}+\dfrac{1}{n^2}}=\dfrac{1}{2}\cdot$
 }
\end{ex}

% nguồn: Ok_TraLoiNgan_Lop11_Lan2
\begin{ex}%[1D3V1-2]%[ma:d43eb23f]
	Cho $L=\lim \dfrac{(an+1)^3(n^2+3n+5)}{(2n-1)^2(n^2+3)^2}$. Tìm  tất cả các giá trị của $a$ để $L=\dfrac{1}{2}$. \\
	\begin{flushright}
	\shortans{1} 
	\end{flushright}
	 \loigiai{
		Ta  có
		\begin{eqnarray*}
			L&=&\lim  \dfrac{(an+1)^3(n^2+3n+5)}{(2n-1)(n^2+3)^2}
			=\lim  \dfrac{n^3\left(a+\dfrac{1}{n}\right)^3\cdot n^2\left(1+\dfrac{3}{n}+\dfrac{5}{n^2}\right)}{n\left(2-\dfrac{1}{n}\right)\cdot n^4\left(1+\dfrac{3}{n^2}\right)^2} \\
			&=&\lim  \dfrac{\left(a+\dfrac{1}{n}\right)^3\cdot \left(1+\dfrac{3}{n}+\dfrac{5}{n^2}\right)}{\left(2-\dfrac{1}{n}\right)\cdot{\left(1+\dfrac{3}{n^2}\right)^2}}
			 =\dfrac{a^3}{2}\cdot
		\end{eqnarray*}
		Từ đề bài suy ra $\dfrac{a^3}{2}=\dfrac{1}{2}\Leftrightarrow  a=1$.\\
		Vậy $a=1$ là giá trị cần  tìm.
	} 
\end{ex}

% nguồn: Ok_TraLoiNgan_Lop11_Lan2
\begin{ex}%[1D3V1-2]%[ma:7d9593d0]
	Có bao nhiêu số  tự nhiên $a$ sao cho
	 \[
	 \lim \sqrt{\dfrac{9^n + 3^{n+1}}{6^n + 9^{n+2}}} \leq \dfrac{1}{3^a}?
	 \]
	 \shortans{3}
	 \loigiai{
		 \[
		\lim \sqrt{\dfrac{9^n + 3^{n+1}}{6^n + 9^{n+2}}} =\lim \sqrt{\dfrac{1 +  3\cdot \left(\dfrac{3}{9} \right)^n}{\left(\dfrac{6}{9} \right)^n + 9^2}} = \dfrac{1}{9}.
		\] 
		Từ điều kiện đầu bài suy ra  $\dfrac{1}{9} \leq \dfrac{1}{3^a} \Rightarrow a \leq 2$. Do $a \in \mathbb{N}$, nên $a \in \{0, 1, 2\}$.
	} 
\end{ex}

% nguồn: Ok_TraLoiNgan_Lop11_Lan2
\begin{ex}%[1D3V1-2]%[ma:7073ceef]
	 Có bao nhiêu giá trị nguyên của tham số $a$ thuộc khoảng $\left( 0; 2\,023 \right)$ để
	\[ 
	\lim \sqrt{\dfrac{9^n +  3^{n+1}}{5^n + 9^{n+a}}} \leq \dfrac{1}{2\,187}?
	 \]
	\shortans{2,02} 
	 \loigiai{
		Ta có 
		\[ 
		\lim \sqrt{\dfrac{9^n + 3^{n+1}}{5^n + 9^{n+a}}} = \lim \sqrt{\dfrac{1 + 3\cdot \left(\dfrac{1}{3} \right)^n}{\left(\dfrac{5}{9} \right)^n + 9^a}}  =\sqrt{\dfrac{1}{9^a}}.
		\] 
		 Từ điều kiện đầu bài suy ra\\
		$\sqrt{\dfrac{1}{9^a}} \leq \dfrac{1}{2\,187}  \Rightarrow 9^a \geq 4\,782\,969 \Rightarrow a \geq 7$.\\
		Do $a \in \mathbb{Z}$ và $7 \leq a  < 2\,023$, nên có $2\,016$ giá trị thỏa mãn.
	 }
\end{ex}

% nguồn: Ok_TraLoiNgan_Lop11_Lan2
\begin{ex}%[1D3V1-2]%[ma:67130b85]
	 Cho dãy số $(u_n)$ với $u_n=\dfrac{1^2+2^2+3^2+ \cdots +n^2}{n^3+2n+7}$. Tính giới hạn của dãy số $(u_n)$.
	 \shortans{0,33}
	 \loigiai{
		Ta có kết  quả quen thuộc $1^2+2^2+3^2+\cdots +n^2=\dfrac{n(n+1)(2n+1)}{6}$.\\
		 Do đó
		$$\begin{aligned}
		\lim \dfrac{1^2+2^2+3^2+4^2+  \cdots +n^2}{n^3+2n+7}
		&=\lim  \dfrac{n(n+1)(2n+1)}{6(n^3+2n+7)} \\
		&=\lim  \dfrac{\left(1+\dfrac{1}{n}\right)\left(2+\dfrac{1}{n}\right)}{6\left(1+\dfrac{2}{n^2}+\dfrac{7}{n^3}\right)}
		 =\dfrac{1\cdot 2}{6}=\dfrac{1}{3}\approx 0{,}33.
	\end{aligned}$$
 }
\end{ex}

% nguồn: Ok_TraLoiNgan_Lop11_Lan2
\begin{ex}%[1D3V1-2]%[ma:2827c575]
	Có bao nhiêu giá trị nguyên của $a$ thuộc  $\left( 0;20 \right)$ sao cho
	 \[
	 \lim \sqrt{3 + \dfrac{a n^2 - 1}{3 + n^2} - \dfrac{1}{2^n}}
	\] 
	là một số  nguyên?
	\shortans{3} 
	 \loigiai{
		\[ 
		\lim \sqrt{3 + \dfrac{a n^2 - 1}{3 + n^2} -  \dfrac{1}{2^n}} = \sqrt{3 + a}
		 \]
		Điều kiện: $a \in (0; 20)$, $a \in \mathbb{Z}$ và $\sqrt{3 + a} \in \mathbb{Z}$, nên $a \in \{1,  6, 13\}$.
	 }
\end{ex}

% nguồn: Ok_TraLoiNgan_Lop11_Lan2
\begin{ex}%[1D3V1-2]%[ma:6b0b1252]
	Cho dãy số $(u_n)$ với $u_n=\dfrac{1+2+3+ \cdots +n}{n^2+1}$. Tính giới  hạn của dãy số $(u_n)$.\\
	 \shortans{0,5}
	 \loigiai{
		Ta  có $\lim{u_n}=\lim \dfrac{1+2+3+ \cdots +n}{n^2+1}
		=\lim \dfrac{n(n+1)}{2(n^2+1)} =\dfrac{1}{2}  = 0{,}5$.
	} 
\end{ex}

% nguồn: Ok_TraLoiNgan_Lop11_Lan2
\begin{ex}%[1D3V1-2]%[ma:6418e9ba]
	Cho dãy số $(u_n)$ với $u_n=\dfrac{1+3+ \cdots +2n+1}{3n^2+4}$.  Tính giới hạn của dãy số $(u_n)$.\\
	\begin{flushright}
	 \shortans{0,33}
	\end{flushright}
	\loigiai{ 
		Ta có $1+3+5+ \cdots +(2n+1)=\dfrac{(1+2n+1)(n+1)}{2}=(n+1)^2$. \\ 
		Khi đó 
		$\lim \dfrac{1+3+5+ \cdots  +(2n+1)}{3n^2+4}
		=\lim \dfrac{(n+1)^2}{3n^2+4} 
		=\lim  \dfrac{1+\dfrac{2}{n}+\dfrac{1}{n^2}}{3+\dfrac{4}{n^2}}
		 =\dfrac{1}{3} \approx 0{,}33$.
	 }
\end{ex}


%%%%%%%%%%%%%%%% TỰ LUẬN %%%%%%%%%%%%%%%%

% nguồn: Ok_TuanLuan_Toan11_Lan2
\begin{ex}%[1D3N1-2]%[ma:e7494d69]
	 Tìm $\lim \dfrac{\sqrt{2n+2}-\sqrt{n}}{\sqrt{n}}$.
	 \loigiai{
		\begin{itemize}
			\item \textbf{Cách 1.} Ta có $\lim \dfrac{\sqrt{2n+2}-\sqrt{n}}{\sqrt{n}}=\lim  \dfrac{\sqrt{2+\dfrac{2}{n}}-1}{1}=\sqrt{2}-1$.
			\item \textbf{Cách 2.} Quan tâm đến số hạng có chứa số mũ cao nhất, ta có thể xem $u_n=\dfrac{\sqrt{2n}-\sqrt{n}}{\sqrt{n}}=\sqrt{2}-1$, sau đó rút gọn ta được $\sqrt{2}-1$. Vậy giới hạn cần  tìm là $\sqrt{2}-1$.
		\end{itemize}
	} 
\end{ex}

% nguồn: Ok_TuanLuan_Toan11_Lan2
\begin{ex}%[1D3N1-2]%[ma:338372f9]
	Tìm $\lim  \dfrac{\sqrt[3]{n^3+n}}{3n+2}$.
	\loigiai{ 
		\begin{itemize}
			\item  \textbf{Cách 1.} Ta có $\lim \dfrac{\sqrt[3]{n^3+n}}{3n+2}=\lim \dfrac{\sqrt[3]{1+\dfrac{1}{n^2}}}{3+\dfrac{2}{n}}=\dfrac{1}{3}$.
			\item  \textbf{Cách 2.} Quan tâm đến số hạng có chứa số mũ cao nhất, ta có thể xem $u_n=\dfrac{\sqrt[3]{n^3}}{3n}$, sau đó rút gọn ta được $\dfrac{1}{3}$. Vậy giới hạn cần tìm là $\dfrac{1}{3}$.
		\end{itemize}
	} 
\end{ex}

% nguồn: Ok_TuanLuan_Toan11_Lan2
\begin{ex}%[1D3H1-2]%[ma:8442427c]
	Tìm $\lim  \dfrac{4^n}{2 \cdot 3^n + 4^n}$.
	\loigiai{ 
		\begin{itemize}
			\item  \textbf{Cách 1.} Ta có
			 \[
			\lim \dfrac{4^n}{2  \cdot 3^n + 4^n} = \lim \dfrac{1}{2 \cdot \left(\dfrac{3}{4}\right)^n + 1} = 1.
			 \]
			\item \textbf{Cách 2.} Chỉ quan tâm đến biểu thức chứa hàm mũ có cơ số lớn nhất là $4$  ở tử và mẫu, ta có thể xem $u_n = \dfrac{4^n}{4^n}$, rút gọn ta được $1$, đó chính là giới hạn cần tìm.
		\end{itemize}
	 }
\end{ex}

% nguồn: Ok_TuanLuan_Toan11_Lan2
\begin{ex}%[1D3H1-2]%[ma:32eb3ecd]
	 Tìm $\lim \dfrac{4 \cdot 3^n + 7^{n+1}}{2 \cdot 5^n + 7^n}$.
	\loigiai{ 
		\begin{itemize}
			 \item \textbf{Cách 1.} Ta có
			 \[
			\lim \dfrac{4 \cdot 3^n + 7^{n+1}}{2 \cdot 5^n + 7^n} = \lim \dfrac{4 \cdot 3^n + 7 \cdot 7^n}{2 \cdot 5^n + 7^n} = \lim \dfrac{4 \cdot \left(\dfrac{3}{7}\right)^n + 7}{2 \cdot \left(\dfrac{5}{7}\right)^n  + 1} = 7.
			\] 
			\item \textbf{Cách 2.} Chỉ quan tâm đến biểu thức chứa hàm mũ có cơ số lớn nhất là $7$ ở tử và mẫu, ta có thể xem $u_n = \dfrac{7^{n+1}}{7^n}$, rút gọn ta được $7$, đó chính là giới  hạn cần tìm.
		\end{itemize}
	 }
\end{ex}

% nguồn: Ok_TuanLuan_Toan11_Lan2
\begin{ex}%[1D3H1-2]%[ma:c5dd6f73]
	Cho  dãy số $(u_n)$ xác định bởi
	\begin{enumerate}
		\item $u_1=1$; $u_{n+1}=\dfrac{u_n-4}{u_n+6}$, $\forall n \geq 1$. Tính giới  hạn $\lim \dfrac{u_n+1}{u_n+4}$.
		\item $u_1=3$, $u_{n+1}=\dfrac{u_n+1}{2}$, $\forall n \geq 1$. Tính giới hạn $\lim  u_n$.
	\end{enumerate}
	 \loigiai{
		\begin{enumerate}
			\item Đặt $v_n=\dfrac{u_n+1}{u_n+4}$. Ta có $v_{n+1}=\dfrac{u_{n+1}+1}{u_{n+1}+4}=\dfrac{2\left(u_n+1\right)}{5\left(u_n+4\right)}=\dfrac{2}{5}  v_n=\ldots=\left(\dfrac{2}{5}\right)^{n+1}$.\\
			Vậy, ta có $v_n=\left(\dfrac{2}{5}\right)^n$ do đó $\lim  \dfrac{u_n+1}{u_n+4}=\lim v_n=0$.
			\item Ta có $u_{n+1}-1=\dfrac{u_n-1}{2}=\dfrac{1}{2^2}\left(u_{n-1}-1\right)=\cdots=\dfrac{1}{2^n}\left(u_1-1\right)=\dfrac{1}{2^{n-1}}$.\\ 
			Do đó, $u_n=\dfrac{1}{2^{n-2}}+1$ nên  $\lim u_n=\lim \left(\dfrac{1}{2^{n-2}}+1\right)=1$.
		\end{enumerate}
	 }
\end{ex}

% nguồn: Ok_TuanLuan_Toan11_Lan2
\begin{ex}%[1D3H1-2]%[ma:fd18a64b]
	 Tính giới hạn $\lim \dfrac{n^4}{(n+1)(2+n)\left(n^2+1 \right)}$
	\loigiai{ 
		\textbf{Cách  1:} $\lim \dfrac{n^4}{(n+1)(2+n)\left(n^2+1 \right)}=\lim \dfrac{1}{\left(1+\dfrac{1}{n} \right)\left(\dfrac{2}{n}+1 \right)\left(1+\dfrac{1}{n^2} \right)}=1$.\\
		\textbf{Cách 2:} Ta quan tâm  đến hệ số của số hạng có số mũ cao nhất của tử, và hệ số của số hạng có bậc cao nhất trong từng thừa số của mẫu, ta có thể xem $u_n=\dfrac{n^4}{n \cdot n \cdot n^2}$, rút gọn ta được $1$. Vậy kết quả giới hạn sẽ bằng $1$.
	 }
\end{ex}

% nguồn: Ok_TuanLuan_Toan11_Lan2
\begin{ex}%[1D3H1-2]%[ma:6e134141]
	 Tìm $\lim\dfrac{(-2)^n-4\cdot5^{n+1}}{2\cdot4^n+3\cdot5^n}$.
	 \loigiai{
		\begin{itemize}
			 \item \textbf{Cách 1.} $\lim\dfrac{(-2)^n-4\cdot5^{n+1}}{2\cdot4^n+3\cdot5^n}=\lim\dfrac{\left(-\dfrac{2}{5}\right)^n-4\cdot5}{2\cdot \left(\dfrac{4}{5}\right)^n+3}=-\dfrac{20}{3}$.
			\item \textbf{Cách 2.} Chỉ quan tâm đến biểu thức chứa hàm mũ có cơ số lớn nhất là $ 5$ ở tử và mẫu, ta có thể xem $u_n=\dfrac{-4\cdot5^{n+1}}{3\cdot5^n}$ rút gọn  ta được $-\dfrac{20}{3}$, đó chính là giới hạn cần tìm.
	\end{itemize}}
\end{ex}

% nguồn: Ok_TuanLuan_Toan11_Lan2
\begin{ex}%[1D3H1-2]%[ma:ef6088b3]
	Tìm $\lim  \dfrac{2^n + 4^n}{4^n - 3^n}$.
	\loigiai{ 
		\begin{itemize}
			\item \textbf{Cách 1.} Ta  có
			 \[
			\lim \dfrac{2^n + 4^n}{4^n - 3^n} = \lim  \dfrac{\left(\dfrac{1}{2}\right)^n + 1}{1 - \left(\dfrac{3}{4}\right)^n} = 1.
			 \]
			\item \textbf{Cách 2.} Chỉ quan tâm đến biểu thức chứa hàm mũ có cơ số lớn nhất là $4$ ở tử và mẫu, ta có thể xem $u_n = \dfrac{4^n}{4^n}$, rút gọn ta được $1$,  đó chính là giới hạn cần tìm.
		\end{itemize}
	} 
\end{ex}

% nguồn: Ok_TuanLuan_Toan11_Lan2
\begin{ex}%[1D3H1-2]%[ma:94d86337]
	 Tính $\lim \dfrac{2 n^{3}-2 n+3}{1-4 n^{3}}$.
	\loigiai{Ta có  $
		\lim  \dfrac{2 n^{3}-2 n+3}{1-4 n^{3}}=\lim \dfrac{2-\dfrac{2}{n^{2}}+\dfrac{3}{n^{3}}}{\dfrac{1}{n^{3}}-4}=-\dfrac{1}{2}
		 $.
	 }
\end{ex}

% nguồn: Ok_TuanLuan_Toan11_Lan2
\begin{ex}%[1D3H1-2]%[ma:1588660d]
	 Tìm $\lim \dfrac{\sqrt{4n^2+3}-2n+1}{\sqrt{n^2+2n}+3n}$.
	 \loigiai{
		\begin{itemize}
			 \item \textbf{Cách 1.} Ta có $\lim \dfrac{\sqrt{4n^2+3}-2n+1}{\sqrt{n^2+2n}+3n}=\lim \dfrac{\sqrt{4+\dfrac{3}{n^2}}-2+\dfrac{1}{n}}{\sqrt{1+\dfrac{2}{n}}+3}=0$.
			\item \textbf{Cách 2.} Quan tâm đến số hạng có chứa số mũ cao nhất, ta có thể xem $u_n=\dfrac{\sqrt{4n^2}-2n}{\sqrt{n^2}+n}$, sau đó  rút gọn ta được $0$. Vậy giới hạn cần tìm là $0$.
		\end{itemize}
	} 
\end{ex}

% nguồn: Ok_TuanLuan_Toan11_Lan2
\begin{ex}%[1D3H1-2]%[ma:2dc0c682]
	Cho dãy số $\left(u_n \right)$ với $u_n=\dfrac{2n\sqrt{n}+1}{n^2+2\sqrt{n}-3}$. Tính  $\lim u_n$.
	\loigiai{ 
		Ta có $u_n=\dfrac{2n\sqrt{n}+1}{n^2+2\sqrt{n}-3}=\dfrac{\dfrac{2n\sqrt{n}+1}{n^2}}{\dfrac{n^2+2\sqrt{n}-3}{n^2}}=\dfrac{\dfrac{2}{\sqrt{n}}+\dfrac{1}{n^2}}{1+\dfrac{2}{n\sqrt{n}}-\dfrac{3}{n^2}}$.\\ 
	Mà $\lim \dfrac{2}{\sqrt{n}}=0$, $\lim \dfrac{1}{n^2}=0$, $\lim \dfrac{2}{n\sqrt{n}}=0$, $\lim  \dfrac{3}{n^2}=0$ nên $\lim u_n=0$. }
\end{ex}

% nguồn: Ok_TuanLuan_Toan11_Lan2
\begin{ex}%[1D3H1-2]%[ma:46cb66d5]
	Cho dãy số $\left(u_n \right)$ với $u_n=\dfrac{(-1)^n \cdot 2^{5n+1}}{3^{5n+2}}$.  Tính $\lim u_n$.
	 \loigiai{
		Ta có $u_n=\dfrac{(-1)^n \cdot 2^{5n+1}}{3^{5n+2}}=\dfrac{(-1)^n \cdot 2 \cdot 2^{5n}}{3^2 \cdot  3^{5n}}=\dfrac{(-1)^n \cdot 2}{9} \cdot \left(\dfrac{2}{3} \right)^{5n}$.\\
		Vì $\left|\dfrac{2}{3} \right|< 1$ nên $\lim \left(\dfrac{2}{3}  \right)^{5n}=0$.\\
	Do đó $\lim u_n=0$.} 
\end{ex}

% nguồn: Ok_TuanLuan_Toan11_Lan2
\begin{ex}%[1D3H1-2]%[ma:3e430ef0]
	Tính giới hạn  $\lim \dfrac{-4n^2+n+2}{2n^2+n+1}$
	 \loigiai{
		\textbf{Cách 1:} $\lim  \dfrac{-4n^2+n+2}{2n^2+n+1}=\lim \dfrac{-4+\dfrac{1}{n}+\dfrac{2}{n^2}}{2+\dfrac{1}{n}+\dfrac{1}{n^2}}=\dfrac{-4}{2}=-2$.\\
		\textbf{Cách 2:} Quan tâm đến hệ số của số hạng có số mũ cao nhất của tử và mẫu, khi đó ta có thể xem $u_n=\dfrac{-4n^2}{2n^2}$, rút gọn ta được $-2$. Vậy giới  hạn cần tìm bằng $-2$.
	} 
\end{ex}

% nguồn: Ok_TuanLuan_Toan11_Lan2
\begin{ex}%[1D3H1-2]%[ma:e65564f9]
	Cho dãy số $\left(u_n \right)$ với $u_n=\dfrac{n+\sqrt{n^2+1}}{n{{.3}^n}}$. Tính  $\lim u_n$.
	\loigiai{ 
		Ta có $u_n=\dfrac{n+\sqrt{n^2+1}}{n \cdot 3^n}=\dfrac{\dfrac{n+\sqrt{n^2+1}}{n}}{\dfrac{n \cdot 3^n}{n}}=\dfrac{1+\sqrt{1+\dfrac{1}n^2}}{3^n}=\dfrac{1}{3^n}\left(1+\sqrt{1+\dfrac{1}n^2} \right)$.\\ 
	Mà $\lim \dfrac{1}{n^2}=0$ nên $\lim \left(1+\sqrt{1+\dfrac{1}{n^n}} \right)=2$ và $\lim \dfrac{1}{3^n}=0$. Do  đó $\lim u_n=-a$.}
\end{ex}

% nguồn: Ok_TuanLuan_Toan11_Lan2
\begin{ex}%[1D3H1-2]%[ma:38366ab5]
	Chứng minh  các giới hạn sau đây:
	\begin{listEX}[3]
		\item  $\lim \dfrac{2n^2+n}{n^2+4}=2$.
		\item $\lim  \dfrac{6n+2}{n+5}=6$.
		\item $\lim  \dfrac{7^n-2.8^n}{8^n+3^n}=-2$.
		\item $\lim  \dfrac{2\cdot 3^n+5^n}{5^n+3^n}=1$.
		 \item $\lim \left(\dfrac{3.3^n-\sin 3n}{3^n}\right)=3$.
		\item $\lim  \left(\sqrt{n^2+n}-n\right)=\dfrac{1}{2}$.
	\end{listEX}
	 \loigiai{
		\begin{enumerate}
			\item Ta có $\lim  \left(\dfrac{2n^2+n}{n^2+4}-2\right)=\lim \dfrac{n-8}{n^2+4}$.\\
			Vì $0\leq\left|\dfrac{n-8}{n^2+4}\right|\leq \dfrac{n}{n^2}=\dfrac{1}{n}$ và  $\lim \dfrac{1}{n}=0$ nên $\lim \left(\dfrac{2n^2+n}{n^2+4}-2\right)=0$.\\
			Do đó  $\lim \dfrac{2n^2+n}{n^2+4}=2$.
			\item Ta có $\lim \left(\dfrac{6n+2}{n+5}-6\right)=\lim  \dfrac{-28}{n+5}$.\\
			Vì $\left|\dfrac{-28}{n+5}\right|< \dfrac{28}{n}$ và $\lim \dfrac{28}{n}=0$ nên $\lim  \left(\dfrac{6n+2}{n+5}-6\right)=0$.\\
			Do đó  $\lim \dfrac{6n+2}{n+5}=6$.
			\item Ta có $\lim \left(\dfrac{7^n-2\cdot 8^n}{8^n+3^n}+2\right)=\lim  \dfrac{7^n+2\cdot 3^n}{8^n+3^n}$.\\
			Vì $0< \left|\dfrac{7^n+2\cdot 3^n}{8^n+3^n}\right|< \dfrac{7^n+2\cdot 3^n}{8^n+3^n} < \dfrac{3\cdot  7^n}{8^n}=3\left(\dfrac{7}{8}\right)^n$ và $\lim \dfrac{3\cdot 7^n}{8^n}=0$ nên $\lim \left(\dfrac{7^n-2\cdot 8^n}{8^n+3^n}+2\right)=0$.\\
			 Do đó $\lim \dfrac{7^n-2\cdot 8^n}{8^n+3^n}=-2$.
			\item Ta có $\lim \left(\dfrac{2\cdot 3^n+5^n}{5^n+3^n}-1\right)=\lim  \dfrac{3^n}{5^n+3^n}$.\\
			Vì $0< \left|\dfrac{3^n}{5^n+3^n}\right|< \dfrac{3^n}{5^n+3^n} < \left(\dfrac{3}{5}\right)^n$ và $\lim \left(\dfrac{3}{5}\right)^n=0$ nên $\lim \left(\dfrac{2\cdot  3^n+5^n}{5^n+3^n}-1\right)=0$.\\
			Do đó $\lim \dfrac{2\cdot  3^n+5^n}{5^n+3^n}=1$.
			\item Ta có  $\lim \left(\dfrac{3\cdot 3^n-\sin 3n}{3^n}-3\right)=\lim \left(\dfrac{-\sin 3n}{3^n}\right)$.\\
			Vì $0\leq\left|\dfrac{-\sin 3n}{3^n}\right|=\dfrac{|-\sin 3n|}{3^n} \leq \dfrac{1}{3^n}=\left(\dfrac{1}{3}\right)^n$ và  $\lim \left(\dfrac{1}{3}\right)^n=0$ nên $\lim \left(\dfrac{-\sin 3n}{3^n}\right)=0$.\\
			Do đó $\lim \left(\dfrac{3\cdot 3^n-\sin 3n}{3^n}\right)=3$. 
			\item Ta có $\lim \left(\sqrt{n^2+n}-n-\dfrac{1}{2}\right)=\lim \dfrac{2\sqrt{n^2+n}-(2n+1)}{2}=\lim  \dfrac{-1}{2\left(2\sqrt{n^2+n}+(2n+1)\right)}$.\\
			Vì $0\leq\left|\dfrac{-1}{2\left(2\sqrt{n^2+n}+(2n+1)\right)}\right|\leq \dfrac{1}{2\left(2\sqrt{n^2+n}+(2n+1)\right)} \leq \dfrac{1}{2\left(2\sqrt{n^2}+2n\right)}=\dfrac{1}{8}  \cdot \dfrac{1}{n}, \forall n \in \mathbb{N}^{*}$ và $\lim \dfrac{1}{8} \cdot \dfrac{1}{n}=\dfrac{1}{8} \lim \dfrac{1}{n}=0$ nên $\lim \dfrac{-1}{2\left(2\sqrt{n^2+n}+(2n+1)\right)}=0$.\\
			Do đó  $\lim \left(\sqrt{n^2+n}-n\right)=\dfrac{1}{2}$.
		\end{enumerate}
	} 
\end{ex}

% nguồn: Ok_TuanLuan_Toan11_Lan2
\begin{ex}%[1D3H1-2]%[ma:2094c132]
	Tìm  $\lim \dfrac{1-2^n}{1+2^n}$.
	\loigiai{ 
		\begin{itemize}
			\item \textbf{Cách 1.} Ta có 
			 \[
			\lim  \dfrac{1-2^n}{1+2^n} = \lim \dfrac{\left(\dfrac{1}{2}\right)^n-1}{\left(\dfrac{1}{2}\right)^n+1} = -1.
			 \]
			\item \textbf{Cách 2.} Chỉ quan tâm đến biểu thức chứa $2^n$ ở tử và mẫu, ta có thể xem $u_n =  \dfrac{-2^n}{2^n}$, rút gọn ta được $-1$, đó chính là giới hạn cần tìm.
		\end{itemize}
	} 
\end{ex}

% nguồn: Ok_TuanLuan_Toan11_Lan2
\begin{ex}%[1D3H1-2]%[ma:99196de0]
	Tìm $\lim\dfrac{2^n+3^n-4^{n+3}}{2^n-3^{n+1}+4^{n-1}}$. 
	 \loigiai{
		\begin{itemize}
			\item \textbf{Cách 1.} $\lim\dfrac{2^n+3^n-4^{n+3}}{2^n-3^{n+1}+4^{n-1}}=\lim\dfrac{\left(\dfrac{1}{2}\right)^n+\left(\dfrac{3}{4}\right)^n-4^3}{\left(\dfrac{1}{2}\right)^n-3\left(\dfrac{3}{4}\right)^n+\dfrac{1}{4}}=-256$. 
			\item \textbf{Cách 2.} Chỉ quan tâm đến biểu thức chứa hàm mũ có cơ số lớn nhất là $ 4$ ở tử và mẫu, ta có thể xem $u_n=\dfrac{-4^{n+3}}{4^{n-1}}$ rút gọn ta được $ 256$, đó chính là giới hạn  cần tìm.
		\end{itemize}
	} 
\end{ex}

% nguồn: Ok_TuanLuan_Toan11_Lan2
\begin{ex}%[1D3H1-2]%[ma:13f13f71]
	Tính $L=\lim  \dfrac{\left(2 n^{2}+1\right)^{4}(n+2)^{9}}{n^{17}+1}$.
	 \loigiai{Ta có $
		L=\lim \dfrac{n^{8}\left(2+\dfrac{1}{n^{2}}\right)^{4} \cdot n^{9}\left(1+\dfrac{2}{n}\right)^{9}}{n^{17}\left(1+\dfrac{1}{n^{17}}\right)}=\lim \dfrac{\left(2+\dfrac{1}{n^{2}}\right)^{4}  \cdot\left(1+\dfrac{2}{n}\right)^{9}}{1+\dfrac{1}{n^{17}}}=16$.
	 }
\end{ex}

% nguồn: Ok_TuanLuan_Toan11_Lan2
\begin{ex}%[1D3H1-2]%[ma:52a03503]
	 Tìm $\lim\dfrac{\left(-1\right)^n\cdot {2^{5n+1}}}{3^{5n+2}}$.
	 \loigiai{
		\begin{itemize}
			\item \textbf{Cách 1.}  $\lim\dfrac{\left(-1\right)^n{2^{5n+1}}}{3^{5n+2}}=\lim\dfrac{\left(-1\right)^{5n}\cdot2\cdot(2)^{5n}}{3^2\cdot(3)^{5n}}=\lim\dfrac{2\cdot\left(\dfrac{-2}{3}\right)^{5n}}{3^2}=0$.
			\item \textbf{Cách 2.} Tử chứa hàm số mũ  có cơ số nhỏ hơn cơ số của hàm số mũ ở mẫu nên giới hạn là $ 0 $.
		\end{itemize}
	} 
\end{ex}

% nguồn: Ok_TuanLuan_Toan11_Lan2
\begin{ex}%[1D3H1-2]%[ma:5a1f028f]
	 $\lim u_n$ với $u_n=\dfrac{5n^2+3n-7}{n^2}$ bằng\\
	\textbf{Cách 1:} Ta có  $\lim u_n=\lim \left(\dfrac{5n^2}{n^2}+\dfrac{3n}{n^2}-\dfrac{7}{n^2} \right)=\lim \left(5+\dfrac{3}{n}-\dfrac{7}{n^2} \right)=5$.\\
	\textbf{Cách 2:} Sử dụng  máy tính bỏ túi tương tự những ví dụ trên.\\
	Đây không phải là giá trị chính xác của giới hạn cần tìm, mà chỉ là giá trị gần đúng của một số hạng với $n$ khá lớn, trong  khi $n$ dần ra vô cực. Tuy nhiên kết quả này cũng giúp ta lựa chọn xấp xỉ ra kết quả gần đúng.
\end{ex}

% nguồn: Ok_TuanLuan_Toan11_Lan2
\begin{ex}%[1D3H1-2]%[ma:f8161cce]
	Tính các  giới hạn sau đây
	\begin{listEX}[3]
		\item $\lim \dfrac{-2n+1}{n}$. 
		 \item $\lim \dfrac{\sqrt{16n^2-2}}{n}$.
		\item $\lim  \dfrac{4}{2n+1}$.
		 \item $\lim \dfrac{n^2-2n+3}{2n^2}$.
		 \item $\lim \dfrac{1}{-3n+7}$.
		 \item $\lim \dfrac{2}{(n+1)(n-2)}$.
		\item  $\lim \dfrac{3}{2^n+1}$.
		\item $\lim \dfrac{2^n}{5+7^n}$. 
		 \item $\lim \dfrac{-3n^2+4n+1}{2n^3-3n+7}$.
		 \item $\lim \dfrac{(2n+1)(n-2)+n}{n^3+n}$.
		\item $\lim  \dfrac{4\cdot 3^n+5^n}{2\cdot 5^n+7^n}$.
		 \item $\lim \dfrac{4^n+7^n}{\left(2^n+1\right)\left(5+7^n\right)}$.
	\end{listEX}
	\loigiai{ 
		\begin{enumerate}
			\item $\lim  \dfrac{-2n+1}{n}=\lim \left(-2+\dfrac{1}{n}\right)=-2$.
			\item $\lim \dfrac{\sqrt{16n^2-2}}{n}=\lim \sqrt{\dfrac{16n^2-2}{n^2}}=\sqrt{\lim  \left(16-\dfrac{2}{n^2}\right)}=\sqrt{16}=4$.
			\item  $\lim \dfrac{4}{2n+1}=\lim \dfrac{\dfrac{4}{n}}{2+\dfrac{1}{n}}=\dfrac0{2+0}=0$.
			\item $\lim  \dfrac{n^2-2n+3}{2n^2}=\lim \left(\dfrac{1}{2}-\dfrac{1}{n}+\dfrac{3}{2n^2}\right)=\dfrac{1}{2}$.
			\item $\lim \dfrac{1}{-3n+7}=\lim \dfrac{1}{n} \cdot  \dfrac{1}{\left(-3+\dfrac{7}{n}\right)}=0\cdot \dfrac{1}{-3+0}=0$.
			 \item $\lim \dfrac{2}{(n+1)(n-2)}=\lim \dfrac{1}{n^2} \cdot \dfrac{1}{\left(1+\dfrac{1}{n}\right)\left(1-\dfrac{2}{n}\right)}=0\cdot \dfrac{1}{(1+0)(1-0)}=0$.
			\item $\lim \dfrac{3}{2^n+1}=\lim \left(\dfrac{1}{2}\right)^n \cdot \dfrac{3}{1+\left(\dfrac{1}{2}\right)^n}=0\cdot \dfrac{3}{1+0}=0$. 
			\item $\lim \dfrac{2^n}{5+7^n}=\lim \left(\dfrac{2}{7}\right)^n  \cdot \dfrac{1}{5\cdot\left(\dfrac{1}{7}\right)^n+1}=0\cdot \dfrac{1}{5\cdot 0+1}=0$.
			 \item $\lim \dfrac{-3n^2+4n+1}{2n^3-3n+7}=\lim \dfrac{1}{n} \cdot\left(\dfrac{-3+\dfrac{4}{n}+\dfrac{1}{n^2}}{2-\dfrac{3}{n^2}+\dfrac{7}{n^3}}\right)=0\cdot \dfrac{-3+0+0}{2-0+0}=0$.
			\item $\lim \dfrac{(2n+1)(n-2)+n}{n^3+n}=\lim \dfrac{2n^2-2n-2}{n^3+n}=\lim \dfrac{1}{n}  \cdot\left(\dfrac{2-\dfrac{2}{n}-\dfrac{2}{n^2}}{1+\dfrac{1}{n^2}}\right)=0\cdot \dfrac{2-0-0}{1+0}=0$.
			\item $\lim \dfrac{4\cdot  3^n+5^n}{2\cdot 5^n+7^n}=\lim \left(\dfrac{5}{7}\right)^n\left(\dfrac{4\cdot\left(\dfrac{3}{5}\right)^n+1}{2\cdot\left(\dfrac{5}{7}\right)^n+1}\right)=0\cdot \dfrac{4\cdot 0+1}{2\cdot 0+1}=0$.
			\item $\lim  \dfrac{4^n+7^n}{\left(2^n+1\right)\left(5+7^n\right)}=\lim \left(\dfrac{1}{2}\right)^n\left(\dfrac{\left(\dfrac{4}{7}\right)^n+1}{\left(1+\left(\dfrac{1}{2}\right)^n\right)\left(5.\left(\dfrac{1}{7}\right)^n+1\right)}\right)=0$.
		\end{enumerate}
	 }
\end{ex}

% nguồn: Ok_TuanLuan_Toan11_Lan2
\begin{ex}%[1D3H1-2]%[ma:095ca29a]
	Tính $\lim \dfrac{\sqrt{n^{4}+2 n+2}}{n^{2}+1}$ 
	\loigiai{Ta có  $
		\lim  \dfrac{\sqrt{n^{4}+2 n+2}}{n^{2}+1}=\lim \dfrac{\sqrt{1+\dfrac{2}{n^{3}}+\dfrac{2}{n^{4}}}}{1+\dfrac{1}{n^{2}}}=1
		$. 
	} 
\end{ex}

% nguồn: Ok_TuanLuan_Toan11_Lan2
\begin{ex}%[1D3H1-2]%[ma:e148e581]
	Tìm  $\lim \dfrac{\sqrt{4n^2-n+1}-n}{\sqrt{9n^2+3n}}$.
	 \loigiai{
		\begin{itemize}
			\item  \textbf{Cách 1.} Ta có $\lim \dfrac{\sqrt{4n^2-n+1}-n}{\sqrt{9n^2+3n}}=\lim \dfrac{\sqrt{4-\dfrac{1}{n}+\dfrac{1}{n^2}}-1}{\sqrt{9+\dfrac{3}{n}}}=\dfrac{1}{3}$.
			\item \textbf{Cách 2.} Quan tâm đến  số hạng có chứa số mũ cao nhất, ta có thể xem $u_n=\dfrac{\sqrt{4n^2}-n}{\sqrt{9n^2}}$, sau đó rút gọn ta được $\dfrac{1}{3}$. Vậy giới hạn cần tìm là $\dfrac{1}{3}$.
		\end{itemize}
	 }
\end{ex}

% nguồn: Ok_TuanLuan_Toan11_Lan2
\begin{ex}%[1D3H1-2]%[ma:795b33e1]
	Cho dãy số $(u_n)$ với $u_n=\dfrac{2n+b}{5n+3}$ trong đó $b$ là tham số thực. Để dãy số $(u_n)$ có giới hạn hữu hạn, giá  trị của $b$ bằng bao nhiêu?
	 \loigiai{
		 Ta có $\lim u_n=\lim \dfrac{2n+b}{5n+3}=\lim \dfrac{2+\dfrac{b}{n}}{5+\dfrac{3}{n}}=\dfrac{2}{5}\ (\forall b \in \mathbb{R})$. \\
		Vậy với  mọi số thực $b$ thì $\lim u_n=\dfrac{2}{5}$ (hữu hạn).\\
		Giải nhanh: $\dfrac{2n+b}{5n+3} \sim \dfrac{2n}{5n}=\dfrac{2}{5}$ với mọi $b  \in \mathbb{R}$.
	 }
\end{ex}

% nguồn: Ok_TuanLuan_Toan11_Lan2
\begin{ex}%[1D3H1-2]%[ma:dec831ab]
	Tìm $\lim \dfrac{3^n - 2 \cdot 5^n}{7 + 3 \cdot  5^n}$.
	\loigiai{ 
		\begin{itemize}
			\item \textbf{Cách 1.}  Ta có
			 \[
			\lim \dfrac{3^n -  2 \cdot 5^n}{7 + 3 \cdot 5^n} = \lim \dfrac{\left(\dfrac{3}{5}\right)^n - 2}{7 \cdot \left(\dfrac{1}{5}\right)^n + 3} = -\dfrac{2}{3}.
			\] 
			\item \textbf{Cách 2.} Chỉ quan tâm đến biểu thức chứa  hàm mũ có cơ số lớn nhất là $5$ ở tử và mẫu, ta có thể xem $u_n = \dfrac{-2 \cdot 5^n}{3 \cdot 5^n}$, rút gọn ta được $-\dfrac{2}{3}$, đó chính là giới hạn cần tìm.
		\end{itemize}
	} 
\end{ex}

% nguồn: Ok_TuanLuan_Toan11_Lan2
\begin{ex}%[1D3H1-2]%[ma:8803c565]
	Tìm $\lim \dfrac{3 \cdot 2^n - 5^n}{5 \cdot 4^n + 6 \cdot 5^n}$. 
	 \loigiai{
		\begin{itemize}
			\item  \textbf{Cách 1.} Ta có
			 \[
			\lim \dfrac{3 \cdot 2^n - 5^n}{5 \cdot 4^n + 6 \cdot 5^n} = \lim \dfrac{3 \cdot \left(\dfrac{2}{5}\right)^n - 1}{5 \cdot \left(\dfrac{4}{5}\right)^n +  6} = -\dfrac{1}{6}.
			 \]
			\item \textbf{Cách 2.} Chỉ quan tâm đến biểu thức chứa hàm mũ có cơ số lớn nhất là $5$ ở tử và mẫu, ta  có thể xem $u_n = \dfrac{-5^n}{6 \cdot 5^n}$, rút gọn ta được $-\dfrac{1}{6}$, đó chính là giới hạn cần tìm.
		\end{itemize}
	} 
\end{ex}

% nguồn: Ok_TuanLuan_Toan11_Lan2
\begin{ex}%[1D3H1-2]%[ma:9d329ff7]
	Tìm  $\lim \dfrac{\sqrt{4n^2-1}+\sqrt[3]{8n^3+2n^2-3}}{\sqrt{16n^2+4n}-\sqrt[4]{n^4+1}}$.
	\loigiai{ 
		\begin{itemize}
			\item \textbf{Cách  1.} Ta có $\lim \dfrac{\sqrt{4n^2-1}+\sqrt[3]{8n^3+2n^2-3}}{\sqrt{16n^2+4n}-\sqrt[4]{n^4+1}}=\lim \dfrac{\sqrt{4-\dfrac{1}{n^2}}+\sqrt[3]{8+\dfrac{2}{n}-\dfrac{3}{n^3}}}{\sqrt{16+\dfrac{4}{n}}-\sqrt[4]{1+\dfrac{1}{n^4}}}=\dfrac{4}{3}$.
			\item \textbf{Cách 2.} Quan tâm đến số hạng có chứa số mũ cao nhất, ta có thể xem $u_n=\dfrac{\sqrt{4n^2}+\sqrt[3]{8n^3}}{\sqrt{16n^2}-\sqrt[4]{n^4}}$, sau đó rút gọn ta  được $\dfrac{4}{3}$. Vậy giới hạn cần tìm là $\dfrac{4}{3}$.
		\end{itemize}
	} 
\end{ex}

% nguồn: Ok_TuanLuan_Toan11_Lan2
\begin{ex}%[1D3H1-2]%[ma:05b908b2]
	Tính giới hạn $\lim {{(2n+1)}^2}\left(\dfrac{3}{n^2+2n}-\dfrac{1}{n^2+3n-1} \right)$. 
	\loigiai{ 
		 \[
		\lim {{(2n+1)}^2}\left(\dfrac{3}{n^2+2n}-\dfrac{1}{n^2+3n-1} \right) = \lim \dfrac{{{(2n+1)}^2} \left(2n^2 +  7n - 3\right)}{\left(n^2 + 2n\right) \left(n^2 + 3n - 1\right)}
		 \]
		\[ 
		= \lim \dfrac{{{\left(2 + \dfrac{1}{n} \right)}^2} \left(2 + \dfrac{7}{n} - \dfrac{3}{n^2} \right)}{\left(1 + \dfrac{2}{n} \right) \left(1 + \dfrac{3}{n} - \dfrac{1}{n^2} \right)} = \dfrac{2^2 \cdot 2}{1 \cdot  1} = 8.
		\] 
	 }
\end{ex}

% nguồn: Ok_TuanLuan_Toan11_Lan2
\begin{ex}%[1D3H1-2]%[ma:ad21e31b]
	[Dự án 11 HVA  2024-2025]
	 Tìm $\lim\dfrac{4^{n+2}+6^{n+1}}{5^{n-1}+2\cdot 6^{n+3}}$.
	 \loigiai{
		\begin{itemize}
			\item \textbf{Cách 1.} $\lim\dfrac{4^{n+2}+6^{n+1}}{5^{n-1}+2\cdot 6^{n+3}}=\lim\dfrac{4^2\cdot4^n+6\cdot  6^n}{\dfrac{1}{5}\cdot{5^n}+2\cdot 6^3\cdot 6^n}=\lim\dfrac{4^2\cdot \left(\dfrac{2}{3}\right)^n+6}{\dfrac{1}{5}\cdot {\left(\dfrac{5}{6}\right)^n}+2\cdot6^3}=\dfrac{1}{72}$.
			\item \textbf{Cách 2.} Chỉ quan tâm đến biểu thức chứa hàm mũ có cơ số lớn nhất là $ 6$ ở tử và mẫu, ta có thể xem $u_n=\dfrac{6^{n+1}}{2\cdot 6^{n+3}}$  rút gọn ta được $\dfrac{1}{72}$, đó chính là giới hạn cần tìm.
		\end{itemize}
	 }
\end{ex}

% nguồn: Ok_TuanLuan_Toan11_Lan2
\begin{ex}%[1D3H1-2]%[ma:271cc09c]
	Cho dãy  số $\left(u_n \right)$ với $u_n=\dfrac{-2n^3+3n^2+4}{n^4+4n^3+n}$. Tính $\lim u_n$.
	 \loigiai{
		Ta có  $u_n=\dfrac{-2n^3+3n^2+4}{n^4+4n^3+n}=\dfrac{\dfrac{-2n^3+3n^2+4}{n^4}}{\dfrac{n^4+4n^3+n}n^4}=\dfrac{-\dfrac{2}{n}+\dfrac{3}n^2+\dfrac{4}{n^4}}{1+\dfrac{4}{n}+\dfrac{1}{n^3}}$.\\
		Mà $\lim \dfrac{2}{n}=0$, $\lim  \dfrac{3}{n^2}=0$,
		$\lim \dfrac{4}{n^4}=0$,  $\lim \dfrac{4}{n}=0$ và $\lim \dfrac{1}{n^3}=0$.\\
		Do đó  $\lim u_n=\dfrac{0+0+0}{1+0+0}=0$.
	} 
\end{ex}

% nguồn: Ok_TuanLuan_Toan11_Lan2
\begin{ex}%[1D3H1-2]%[ma:6290baf9]
	Tìm ${\lim \dfrac{\sqrt{2n+1}}{\sqrt{n}+1}}$. 
	\loigiai{ 
		\textbf{Cách  1:} $\lim \dfrac{\sqrt{2n+1}}{\sqrt{n}+1}=\lim \dfrac{\sqrt{2+\dfrac{1}{n}}}{1+\sqrt{\dfrac{1}{n}}}=\sqrt{2}$.\\
		\textbf{Cách 2:} Quan tâm đến số hạng có chứa số mũ cao nhất, ta có thể xem $u_n=\dfrac{\sqrt{2n}}{\sqrt{n}}$ sau đó rút gọn ta được $\sqrt{2}$. Vậy giới hạn cần  tìm là $\sqrt{2}$.
	} 
\end{ex}

% nguồn: Ok_TuanLuan_Toan11_Lan2
\begin{ex}%[1D3H1-2]%[ma:6d46d245]
	Tính các giới hạn sau đây 
	\begin{listEX}[2]
		 \item $\lim \dfrac{n^2-4n^3}{2n^3+5n-2}$.
		\item  $\lim \dfrac{n^3-7n}{1+2n^2}$.
		\item $\lim  \dfrac{n+2}{n^2+n+1}$.
		\item $\lim  \dfrac{n^4}{(n+1)(2+n)\left(n^2+1\right)}$.
		\item $\lim  \dfrac{-4n^2+n+2}{2n^2+n+1}$.
		\item  $\lim \dfrac{5n^2+3n-7}{n^2}$.
	\end{listEX}
	\loigiai{ 
		\begin{enumerate}
			 \item Ta có $\lim \dfrac{n^2-4n^3}{2n^3+5n-2}=\lim \dfrac{n^3\left(\dfrac{1}{n}-4\right)}{n^3\left(2+\dfrac{5}{n^2}-\dfrac{2}{n^3}\right)}=\lim \dfrac{\dfrac{1}{n}-4}{2+\dfrac{5}{n^2}-\dfrac{2}{n^3}}=\dfrac{-4}{2}=-2$.
			\item Ta có $\lim \dfrac{n^3-7n}{1+2n^2}=\lim  \dfrac{n^3\left(1-\dfrac{7}{n^2}\right)}{n^2\left(\dfrac{1}{n^2}+2\right)}=\lim \left(n \cdot \dfrac{1-\dfrac{7}{n^2}}{\dfrac{1}{n^2}+2}\right)=+\infty$.\\
			Vì $\lim \left(\dfrac{1-\dfrac{7}{n^2}}{\dfrac{1}{n^2}+2}\right)=\dfrac{1}{2} > 0$;  $\lim n=+\infty$.
			\item Ta có $\lim \dfrac{n+2}{n^2+n+1}=\lim \dfrac{n\left(1+\dfrac{3}{n}\right)}{n^2\left(1+\dfrac{1}{n}+\dfrac{1}{n^2}\right)}=\lim  \dfrac{1}{n} \cdot \dfrac{1+\dfrac{3}{n}}{1+\dfrac{1}{n}+\dfrac{1}{n^2}}=0$.
			\item Ta có $\begin{aligned}
			\lim  \dfrac{n^4}{(n+1)(2+n)\left(n^2+1\right)}&=\lim \dfrac{n^4}{n\left(1+\dfrac{1}{n}\right) \cdot n\left(\dfrac{2}{n}+1\right) n^2\left(1+\dfrac{1}{n^2}\right)}\\
			&=\lim  \dfrac{1}{\left(1+\dfrac{1}{n}\right)\left(\dfrac{2}{n}+1\right)\left(1+\dfrac{1}{n^2}\right)}\\
			&=1. 
		\end{aligned}$
		\item Ta  có $\lim \dfrac{-4n^2+n+2}{2n^2+n+1}=\lim \dfrac{-4+\dfrac{1}{n}+\dfrac{2}{n^2}}{2+\dfrac{1}{n}+\dfrac{1}{n^2}}=\dfrac{-4}{2}=-2$.
		\item Ta có $\lim  \dfrac{5n^2+3n-7}{n^2}=\lim \left(\dfrac{5n^2}{n^2}+\dfrac{3n}{n^2}-\dfrac{7}{n^2}\right)=\lim \left(5+\dfrac{3}{n}-\dfrac{7}{n^2}\right)=5$.
	\end{enumerate}
 }
\end{ex}

% nguồn: Ok_TuanLuan_Toan11_Lan2
\begin{ex}%[1D3H1-2]%[ma:f17e05ff]
	Tìm $\lim \dfrac{\sqrt{2n+1}-\sqrt{n+3}}{\sqrt{4n-5}}$. 
	\loigiai{ 
		\begin{itemize}
			\item \textbf{Cách 1.} Ta có $\lim \dfrac{\sqrt{2n+1}-\sqrt{n+3}}{\sqrt{4n-5}}=\lim \dfrac{\sqrt{2+\dfrac{1}{n}}-\sqrt{1+\dfrac{3}{n}}}{\sqrt{4-\dfrac{5}{n}}}=\dfrac{\sqrt{2}-1}{2}$. 
			\item \textbf{Cách 2.} Quan tâm đến số hạng có chứa số mũ cao nhất, ta có thể xem $u_n=\dfrac{\sqrt{2n}-\sqrt{n}}{\sqrt{4n}}$, sau đó rút gọn ta  được $\dfrac{\sqrt{2}-1}{2}$. Vậy giới hạn cần tìm là $\dfrac{\sqrt{2}-1}{2}$.
		\end{itemize}
	} 
\end{ex}

% nguồn: Ok_TuanLuan_Toan11_Lan2
\begin{ex}%[1D3H1-2]%[ma:e36591b7]
	 Cho dãy số $\left(u_n \right)$ với $u_n=\dfrac{1}{2}\cdot \dfrac{3}{4}\cdot \dfrac{5}{6}\cdots \dfrac{2n-1}{2n}$. Tính $\lim u_n$.
	 \loigiai{
		Ta có  $\dfrac{2k-1}{2k}=\dfrac{2k-1}{\sqrt{4k^2}}\le \dfrac{2k-1}{\sqrt{4k^2-1}}=\sqrt{\dfrac{2k-1}{2k+1}}$, $\forall k\in \mathbb{N}^*$.
		 Do đó
		\begin{eqnarray*}
			\dfrac{1}{2} &\le& \sqrt{\dfrac{1}{3}};  \\
			\dfrac{3}{4} &\le& \sqrt{\dfrac{3}{5}};  \\
			& \cdots  &\\
			 \dfrac{2n-1}{2n} &\le& \sqrt{\dfrac{2n-1}{2n+1}}.
		\end{eqnarray*}
		Suy ra 
		\[ 
		\dfrac{1}{2} \cdot  \dfrac{3}{4}\cdots \dfrac{2n-1}{2n}\le \sqrt{\dfrac{1}{3}}\cdot\sqrt{\dfrac{3}{5}}\cdots\sqrt{\dfrac{2n-1}{2n+1}}.
		\] 
		$\Rightarrow u_n\le \dfrac{1}{\sqrt{2n+1}}$.\\ 
		Vì vậy, $\left|u_n\right|\le \dfrac{1}{\sqrt{2n+1}}$, $\forall n$.  Mà $\lim \dfrac{1}{\sqrt{2n+1}}=0$, nên $\lim u_n=0$.
	} 
\end{ex}

% nguồn: Ok_TuanLuan_Toan11_Lan2
\begin{ex}%[1D3H1-2]%[ma:4c1b8aba]
	Tính $\lim  \dfrac{3^{n+1}-4^{n}}{4^{n-1}+3}$.
	\loigiai{Ta  có $
		\lim \dfrac{3^{n+1}-4^{n}}{4^{n-1}+3}=\lim \dfrac{9 \cdot  3^{n-1}-4 \cdot 4^{n-1}}{4^{n-1}+3}=\lim \dfrac{9 \cdot\left(\dfrac{3}{4}\right)^{n-1}-4}{1+3 \cdot\left(\dfrac{1}{4}\right)^{n-1}}=-4$.
	 }
\end{ex}

% nguồn: Ok_TuanLuan_Toan11_Lan2
\begin{ex}%[1D3H1-2]%[ma:3e825b42]
	 Tìm $\lim\dfrac{(-2)^n+3^n}{(-2)^{n+1}+3^{n+1}}$.
	 \loigiai{
		\begin{itemize}
			\item \textbf{Cách 1.}  $\lim\dfrac{(-2)^n+3^n}{(-2)^{n+1}+3^{n+1}}=\lim\dfrac{\left(-\dfrac{2}{3}\right)^n+1}{-2\cdot\left(-\dfrac{2}{3}\right)^n+3}=\dfrac{1}{3}$.
			\item \textbf{Cách 2.} Chỉ quan tâm đến biểu thức chứa hàm mũ có cơ số lớn nhất là $ 3$ ở tử và mẫu, ta có thể xem $u_n=\dfrac{3^n}{3^{n+1}}$ rút gọn ta được $\dfrac{1}{3}$, đó chính là giới  hạn cần tìm.
		\end{itemize}
	} 
\end{ex}

% nguồn: Ok_TuanLuan_Toan11_Lan2
\begin{ex}%[1D3H1-2]%[ma:50ad9a2c]
	Tính các  giới hạn sau đây
	\begin{listEX}[2]
		\item $\lim \left(2^n+3^n\right)$. 
		\item $\lim  \left[-4^n+(-2)^n\right]$.
		\item $\lim  \left(\dfrac{1+3^n}{3\cdot 3^n+2^n}\right)$.
		\item $\lim \left(\dfrac{4\cdot 3^n-2^n}{2\cdot  5^n+4^n}\right)$.
		\item  $\lim \left(\dfrac{7^n+1}{-2\cdot 3^n-3\cdot 6^n}\right)$.
		 \item $\lim \left(\dfrac{5^{n+1}-4^n+1}{2\cdot 5^n-6^n}\right)$.
		\item $\lim \dfrac{1+3+3^2+\cdots+3^n}{2\cdot 3^{n+1}+2^n}$. 
		\item  $\lim \dfrac{1+\dfrac{1}{3}+\left(\dfrac{1}{3}\right)^2+\cdots+\left(\dfrac{1}{3}\right)^n}{1+\dfrac{2}{5}+\left(\dfrac{2}{5}\right)^2+\cdots+\left(\dfrac{2}{5}\right)^n}$.
		\item $\lim \dfrac{3^n-2\cdot 5^{n+1}}{2^{n+1}+5^n}$. 
		 \item $\lim \dfrac{3^n-4\cdot 2^{n+1}-3}{3\cdot 2^n+4^n}$.
		\item $\lim  \dfrac{(-1)^n 2^{5n+1}}{3^{5n+2}}$.
		\item $\lim \dfrac{3^n-4\cdot  2^{n+1}-3}{3\cdot 2^n+4^n}$.
	\end{listEX}
	 \loigiai{
		\begin{enumerate}
			\item $\lim  \left(2^n+3^n\right)=\lim \left\{3^n\left[\left(\dfrac{2}{3}\right)^n+1\right]\right\}=+\infty$ vì $\lim 3^n=+\infty$ và $\lim\left[\left(\dfrac{2}{3}\right)^n+1\right]=1$.
			\item $\lim \left[-4^n+(-2)^n\right]=\lim \left\{4^n\left[-1+\left(-\dfrac{2}{4}\right)^n\right]\right\}=-\infty$ vì $\lim 4^n=+\infty$ và  $\lim\left[-1+\left(-\dfrac{2}{4}\right)^n\right]=-1$.
			\item $\lim \left(\dfrac{1+3^n}{3\cdot  3^n+2^n}\right)=\lim \left(\dfrac{\dfrac{1+3^n}{3^n}}{\dfrac{3\cdot 3^n+2^n}{3^n}}\right)=\lim \left(\dfrac{\dfrac{1}{3^n}+1}{3+\dfrac{2^n}{3^n}}\right)=\dfrac{1}{3}$.
			\item  $\lim \left(\dfrac{4\cdot 3^n-2^n}{3\cdot 5^n+4^n}\right)=\lim \left(\dfrac{\dfrac{4\cdot 3^n-2^n}{5^n}}{\dfrac{3\cdot 5^n+4^n}{5^n}}\right)=\lim \left(\dfrac{4\cdot \dfrac{3^n}{5^n}-\dfrac{2^n}{5^n}}{2+\dfrac{4^n}{5^n}}\right)=0$.
			\item $\lim \left(\dfrac{7^n+1}{-2\cdot 3^n-3\cdot 6^n}\right)=\lim \left(\dfrac{\dfrac{7^n+1}{7^n}}{\dfrac{-2\cdot 3^n-3\cdot 6^n}{7^n}}\right)=\lim \left(\dfrac{1+\dfrac{1}{7^n}}{-2\cdot \dfrac{3^n}{7^n}-3\cdot \dfrac{6^n}{7^n}}\right)=-\infty$. 
			\item $\lim \left(\dfrac{5^{n+1}-4^n+1}{2\cdot 5^n-6^n}\right)=\lim \left(\dfrac{\dfrac{5^{n+1}-4^n+1}{6^n}}{\dfrac{2\cdot 5^n-6^n}{6^n}}\right)=\lim \dfrac{5\cdot  \left(\dfrac{5}{6}\right)^n-\left(\dfrac{2}{3}\right)^n+\left(\dfrac{1}{6}\right)^n}{2\cdot \left(\dfrac{5}{6}\right)^n-1}=0$.
			\item $\lim \dfrac{1+3+3^2+\cdots+3^n}{2\cdot 3^{n+1}+2^n}=\lim \dfrac{1-\dfrac{3}{2}\left(1-3^n\right)}{2\cdot 3^{n+1}+2^n}=\dfrac{1}{4}$. 
			\item Đặt  $u_n=1+\dfrac{1}{3}+\left(\dfrac{1}{3}\right)^2+\cdots+\left(\dfrac{1}{3}\right)^n$; $v_n=1+\dfrac{2}{5}+\left(\dfrac{2}{5}\right)^2+\cdots+\left(\dfrac{2}{5}\right)^n$.\\
			Ta có $u_n=1+\dfrac{1}{3} \cdot \dfrac{1-\left(\dfrac{1}{3}\right)^n}{1-\dfrac{1}{3}}=1+\dfrac{1}{2}\left(1-\dfrac{1}{3^n}\right)$.\\ 
			Tương tự  $v_n=1+\dfrac{2}{3}\left(1-\dfrac{2^n}{5^n}\right)$.\\
			Từ đó, $\lim u_n=\dfrac{3}{2}$,  $\lim v_n=\dfrac{5}{3}$.\\
			Vậy  $\lim \dfrac{1+\dfrac{1}{3}+\left(\dfrac{1}{3}\right)^2+\cdots+\left(\dfrac{1}{3}\right)^n}{1+\dfrac{2}{5}+\left(\dfrac{2}{5}\right)^2+\cdots+\left(\dfrac{2}{5}\right)^n}=\dfrac{9}{10}$.
			\item Giải nhanh: $\dfrac{3^n-2\cdot 5^{n+1}}{2^{n+1}+5^n} \sim \dfrac{-2\cdot  5^{n+1}}{5^n}=-10$.\\
			Hoặc cụ thể: $\lim \dfrac{3^n-2\cdot 5^{n+1}}{2^{n+1}+5^n}=\lim \dfrac{\left(\dfrac{3}{5}\right)^n-10}{2\cdot \left(\dfrac{2}{5}\right)^n+1}=-10$. 
			\item  $\lim \dfrac{3^n-4\cdot 2^{n+1}-3}{3\cdot 2^n+4^n}=\lim \dfrac{\left(\dfrac{3}{4}\right)^n-8\cdot\left(\dfrac{1}{2}\right)^n-3\cdot\left(\dfrac{1}{4}\right)^n}{3\cdot\left(\dfrac{1}{2}\right)^n+1}=\dfrac0{1}=0$.
			\item $\lim \dfrac{(-1)^n 2^{5n+1}}{3^{5n+2}}=\lim\left[(-1)^n  \cdot \dfrac{2}{9}\cdot \left(\dfrac{2}{3}\right)^n\right]=0$.
			\item Ta có  $\dfrac{3^n-4\cdot 2^{n+1}-3}{3\cdot 2^n+4^n}=\dfrac{\left(\dfrac{3}{4}\right)^n-4\cdot 2\left(\dfrac{2}{4}\right)^n-\dfrac{3}{n^4}}{3\cdot \left(\dfrac{2}{4}\right)^n+1}$ suy ra $\lim \dfrac{3^n-4\cdot 2^{n+1}-3}{3\cdot 2^n+4^n}=\dfrac0{1}=0$.
		\end{enumerate}
	 }
\end{ex}

% nguồn: Ok_TuanLuan_Toan11_Lan2
\begin{ex}%[1D3H1-2]%[ma:cd89ff0e]
	 Cho dãy số $(u_n)$ với $u_n=\dfrac{4n^2+n+2}{a n^2+5}$. Để dãy số đã cho có giới hạn bằng $2$, giá trị của $a$ bằng bao nhiêu?
	\loigiai{ 
		Ta có $2=\lim u_n=\lim \dfrac{4n^2+n+2}{a  n^2+5}=\lim \dfrac{4+\dfrac{1}{n}+\dfrac{2}{n^2}}{a+\dfrac{5}{n^2}}=\dfrac{4}{a}\ (a \neq 0) \Leftrightarrow a=2$.\\
		Giải nhanh: $2\sim \dfrac{4n^2+n+2}{a n^2+5} \sim \dfrac{4n^2}{a  n^2}=\dfrac{4}{a} \Leftrightarrow a=2$.
	} 
\end{ex}

% nguồn: Ok_TuanLuan_Toan11_Lan2
\begin{ex}%[1D3H1-2]%[ma:a5c4a05e]
	Tìm $\lim\dfrac{\pi  ^{n+1}+3^n+2^n}{5\cdot\pi ^n-4\cdot3^n+2^{n+2}}$.
	\loigiai{ 
		\begin{itemize}
			\item \textbf{Cách  1.} $\lim\dfrac{\pi ^{n+1}+3^n+2^n}{5\cdot\pi ^n-4\cdot3^n+2^{n+2}}=\lim\dfrac{\pi+\left(\dfrac{3}{\pi}\right)^n+\left(\dfrac{2}{\pi}\right)^n}{5-4\cdot\left(\dfrac{3}{\pi}\right)^n+2^2\cdot\left(\dfrac{2}{\pi}\right)^n}=\dfrac{\pi}{5}$.
			\item \textbf{Cách 2.} Chỉ quan tâm đến biểu thức chứa hàm  mũ có cơ số lớn nhất là $\pi $ ở tử và mẫu, ta có thể xem $u_n=\dfrac{\pi ^{n+1}}{5\cdot\pi ^n}$ rút gọn ta được $\dfrac{\pi}{5}$, đó chính là giới hạn cần tìm.
		\end{itemize}
	} 
\end{ex}

% nguồn: Ok_TuanLuan_Toan11_Lan2
\begin{ex}%[1D3H1-2]%[ma:ddbad7b7]
	Tìm  giới hạn của các dãy số sau
	\begin{listEX}[3]
		\item  $u_n=\dfrac{-3n^3+2n^2+1}{2n^3-n}$.
		 \item $u_n=\dfrac{2n+1}{n^2+n+3}$.
		\item  $u_n=\dfrac{2025n+2024}{n-2023}$.
	\end{listEX}
	\loigiai{ 
		\begin{enumerate}
			 \item Ta có $\lim u_n=\lim \dfrac{-3n^3+2n^2+1}{2n^3-n}=\lim \dfrac{n^3\left(-3+\dfrac{2}{n}+\dfrac{1}{n^3}\right)}{n^3\left(2-\dfrac{1}{n^2}\right)}=\lim \dfrac{-3+\dfrac{2}{n}+\dfrac{1}{n^3}}{2-\dfrac{1}{n^2}}=-\dfrac{3}{2}$.
			\item  Ta có $\lim u_n=\lim \dfrac{2n+1}{n^2+n+3}=\lim \dfrac{n\left(2+\dfrac{1}{n}\right)}{n^2\left(1+\dfrac{1}{n}+\dfrac{3}{n^2}\right)}=\lim \left(\dfrac{1}{n} \cdot \dfrac{\dfrac{2}{n}+\dfrac{1}{n^2}}{1+\dfrac{1}{n}+\dfrac{3}{n^2}}\right)=0$.
			\item Ta có $\lim u_n=\lim  \dfrac{2025n+2024}{n-2023}=\lim \dfrac{n\left(2025+\dfrac{2024}{n}\right)}{n\left(1-\dfrac{2023}{n}\right)}=\lim \dfrac{2025+\dfrac{2024}{n}}{1-\dfrac{2023}{n}}=2025$.
		\end{enumerate}
	 }
\end{ex}

% nguồn: Ok_TuanLuan_Toan11_Lan2
\begin{ex}%[1D3H1-2]%[ma:5da2d94a]
	Cho dãy số $\left(u_n \right)$ với $u_n=\dfrac{(-5)^n+4^n}{(-7)^{n+1}+4^{n+1}}$. Tính  $\lim u_n$.
	\loigiai{ 
		Ta có $u_n=\dfrac{(-5)^n \left(1+\dfrac{4^n}{(-5)^n} \right)}{(-7)^n\left(-7+\dfrac{4  \cdot 4^n}{(-7)^n} \right)}=\left(\dfrac{5}{7} \right)^n \cdot \left(\dfrac{1+\left(-\dfrac{4}{5} \right)^n}{-7+4 \cdot \left(\dfrac{-4}{7} \right)^n} \right)$.\\
		Vì $\lim \left(-\dfrac{4}{5} \right)^n =\lim \left(-\dfrac{4}{7} \right)^n=0$ nên $\lim \left(\dfrac{1+\left(-\dfrac{4}{5} \right)^n}{-7+4 \cdot \left(\dfrac{-4}{7}  \right)^n} \right)=-\dfrac{1}{7}$ và $\lim \left(\dfrac{5}{7} \right)^n=0$.\\
	Do đó $\lim u_n=0$.} 
\end{ex}

% nguồn: Ok_TuanLuan_Toan11_Lan2
\begin{ex}%[1D3H1-2]%[ma:59af4dcf]
	Tìm $\lim\dfrac{\pi ^n+3^n+2^{2n}}{3\pi ^n-3^n+2^{2n+2}}$. 
	\loigiai{ 
		\begin{itemize}
			\item \textbf{Cách 1.} $\lim\dfrac{\pi ^n+3^n+2^{2n}}{3\pi ^n-3^n+2^{2n+2}}=\lim\dfrac{\left(\dfrac{\pi}{4}\right)^n+\left(\dfrac{3}{4}\right)^n+1}{3\cdot\left(\dfrac{\pi}{4}\right)^n-\left(\dfrac{3}{4}\right)^n+2^2}=\dfrac{1}{4}$. 
			\item \textbf{Cách 2.} Chỉ quan tâm đến biểu thức chứa hàm mũ có cơ số  lớn nhất là $ 4$ ở tử và mẫu, ta có thể xem $u_n=\dfrac{2^{2n}}{2^{2n+2}}$ rút gọn ta được $\dfrac{1}{4}$, đó chính là giới hạn cần tìm.
		\end{itemize}
	} 
\end{ex}

% nguồn: Ok_TuanLuan_Toan11_Lan2
\begin{ex}%[1D3H1-2]%[ma:0472bd31]
	Tìm  $\lim\dfrac{2^n-3^n+4\cdot 5^{n+2}}{2^{n+1}+3^{n+2}+5^{n+1}}$.
	 \loigiai{
		\begin{itemize}
			 \item \textbf{Cách 1.} $\lim\dfrac{2^n-3^n+4\cdot 5^{n+2}}{2^{n+1}+3^{n+2}+5^{n+1}}=\lim\dfrac{\left(\dfrac{2}{5}\right)^n-\left(\dfrac{3}{5}\right)^n+4\cdot5^2}{2\cdot \left(\dfrac{2}{5}\right)^n+3^2\cdot \left(\dfrac{3}{5}\right)^n+5}=20$.
			\item \textbf{Cách 2.} Chỉ quan tâm đến biểu thức chứa hàm mũ có cơ số lớn nhất là $ 5$ ở tử và mẫu, ta có thể xem  $u_n=\dfrac{4\cdot 5^{n+2}}{5^{n+1}}$ rút gọn ta được $ 20$, đó chính là giới hạn cần tìm.
	\end{itemize}}
\end{ex}

% nguồn: Ok_TuanLuan_Toan11_Lan2
\begin{ex}%[1D3H1-2]%[ma:5ef04395]
	 Tìm $\lim\dfrac{\left(\sqrt 5\right)^n-2^{n+1}+1}{5\cdot2^n+\left(\sqrt 5\right)^{n+1}-3}$.
	 \loigiai{
		\begin{itemize}
			\item \textbf{Cách 1.} $\lim\dfrac{\left(\sqrt 5\right)^n-2^{n+1}+1}{5\cdot2^n+\left(\sqrt 5\right)^{n+1}-3}=\lim\dfrac{1-2\cdot\left(\dfrac{2}{\sqrt 5}\right)^n+\left(\dfrac{1}{\sqrt 5}\right)^n}{5\cdot\left(\dfrac{2}{\sqrt 5}\right)^n+\sqrt  5-3\cdot\left(\dfrac{1}{\sqrt 5}\right)^n}=\dfrac{1}{\sqrt 5}$.
			\item \textbf{Cách 2.} Chỉ quan tâm đến biểu thức chứa hàm mũ  có cơ số lớn nhất là $\sqrt 5 $ ở tử và mẫu, ta có thể xem $u_n=\dfrac{\left(\sqrt 5\right)^n}{\left(\sqrt 5\right)^{n+1}}$ rút gọn ta được $\dfrac{1}{\sqrt 5}$, đó chính là giới hạn cần tìm.
		\end{itemize}
	} 
\end{ex}

% nguồn: Ok_TuanLuan_Toan11_Lan2
\begin{ex}%[1D3H1-2]%[ma:5a4e78a0]
	Tính $\lim  \dfrac{n^{3}+1}{n^{2}-2 n+3}$.
	\loigiai{Ta  có $
		\lim \dfrac{n^{3}+1}{n^{2}-2 n+3}=\lim  \dfrac{n^{3}\left(1+\dfrac{1}{n^{3}}\right)}{n^{2}\left(1-\dfrac{2}{n}+\dfrac{3}{n^{2}}\right)}=\lim \left[n \cdot \dfrac{\left(1+\dfrac{1}{n^{3}}\right)}{\left(1-\dfrac{2}{n}+\dfrac{3}{n^{2}}\right)}\right]=+\infty$.
	 }
\end{ex}

% nguồn: Ok_TuanLuan_Toan11_Lan2
\begin{ex}%[1D3V1-2]%[ma:fc58dd57]
	Tìm tất cả các giá trị của tham số $a$  để $L=\lim \dfrac{5 n^{2}-3 a n^{4}}{(1-a) n^{4}+2 n+1}>0$.
	\loigiai{Ta  có $
		L=\lim \dfrac{5 n^{2}-3 a n^{4}}{(1-a) n^{4}+2 n+1}=\lim  \dfrac{\dfrac{5}{n^{2}}-3 a}{(1-a)+\dfrac{2}{n^{3}}+\dfrac{1}{n^{4}}}=\dfrac{-3 a}{(1-a)}$.\\
		Khi đó $L>0 \Leftrightarrow \dfrac{-3a}{(1-a)} >0 \Leftrightarrow  \hoac{&a<0\\
		&a>1.}$ 
	 }
\end{ex}

% nguồn: Ok_TuanLuan_Toan11_Lan2
\begin{ex}%[1D3V1-2]%[ma:ec322d25]
	Cho dãy số $\left(u_n\right)$  được xác định bởi $\heva{
	 &u_{1}=1\\
	&u_{n+1}=\dfrac{2\left(2 u_n+1\right)}{u_n+3},\,  n=1,2,3,\ldots }$\\
	Biết dãy số $\left(u_n\right)$ có giới hạn hữu hạn. Tính  $\lim u_n$.
	\loigiai{Bằng phương pháp quy nạp, dễ dàng chứng minh được $u_n>0$ với  mọi $n\in \mathbb{N}^*$.\\
		Đặt $\lim u_n=L \geq 0$.\\ 
		Vì $\lim u_{n+1}=\lim \dfrac{2\left(2 u_n+1\right)}{u_n+3}$ nên  $L=\dfrac{2(2 L+1)}{L+3} \Leftrightarrow L^{2}-L-2=0 \Leftrightarrow
		 \hoac{&L=2 \text{ (nhận)}\\
		&L=-1 \text{  (loại).}
		 }$\\
		Vậy $\lim  u_n=2$.
	} 
\end{ex}

% nguồn: Ok_TuanLuan_Toan11_Lan2
\begin{ex}%[1D3V1-2]%[ma:db46558a]
	Tìm $\lim u_n$ biết $\left(u_n\right)$ có giới hạn hữu hạn và 
	 $\left(u_n \right)\colon \heva{
	 &u_1=\sqrt{2}\\
	&u_{n+1}=\sqrt{2+u_n}, \, n=1,2,3,\ldots}$ 
	\loigiai{Bằng phương pháp quy nạp, dễ dàng chứng minh được $u_n>0$  với mọi $n\in \mathbb{N}^*$.\\
		Đặt $\lim u_n=a\geq 0$.\\ 
		Do $u_{n+1}=\sqrt{2+u_n}$ nên $\lim u_{n+1}=\lim  \sqrt{2+u_n}\Leftrightarrow a=\sqrt{2+a} \Leftrightarrow \hoac{
		&a=2 \text{  (nhận)}\\
		 &a=-1 \text{ (loại).}}$
	 }
\end{ex}

% nguồn: Ok_TuanLuan_Toan11_Lan2
\begin{ex}%[1D3V1-2]%[ma:20fae584]
	Có bao nhiêu giá trị nguyên của $a$ thuộc $(0; 20)$ sao cho $\lim \sqrt{3+\dfrac{a n^2-1}{3+n^2}-\dfrac{1}{2^n}}$ là  một số nguyên?
	\loigiai{ 
		Ta có  $\heva{&\lim \dfrac{a n^2-1}{3+n^2}=\lim \dfrac{a-\dfrac{1}{n^2}}{\dfrac{3}{n^2}+1}=a \\ &\lim \dfrac{1}{2^n}=\lim \left(\dfrac{1}{2}\right)^n=0} \Rightarrow \lim \sqrt{3+\dfrac{a n^2-1}{3+n^2}-\dfrac{1}{2^n}}=\sqrt{3+a}$.\\
		Ta có $\heva{&a \in(0; 20), a \in \mathbb{Z} \\ &\sqrt{a+3} \in \mathbb{Z} \Rightarrow  a \in\{1; 6; 13\}.}$\\
		Vậy có $3$ giá  trị nguyên của $a$ thỏa mãn yêu cầu bài toán.
	} 
\end{ex}

% nguồn: Ok_TuanLuan_Toan11_Lan2
\begin{ex}%[1D3V1-2]%[ma:d080531b]
	 Tính $\lim \dfrac{1+2+2^{2}+\ldots+2^{n}}{1+3+3^{2}+\ldots+3^{n}}$.
	\loigiai{Ta có 
		\begin{eqnarray*}
			 \lim \dfrac{1+2+2^{2}+\ldots+2^{n}}{1+3+3^{2}+\ldots+3^{n}}&=&\lim \dfrac{\dfrac{1-2^{n+1}}{-1}}{\dfrac{1-3^{n+1}}{-2}}\\
			&=&\lim \dfrac{\left(1-2^{n+1}\right)  \cdot 2}{1-3^{n+1}}\\
			 &=&\lim \dfrac{\left(\left(\dfrac{1}{3}\right)^{n+1}-\left(\dfrac{2}{3}\right)^{n+1}\right) \cdot 2}{\left(\dfrac{1}{3}\right)^{n+1}-1}\\
			 &=&0.
		\end{eqnarray*}
	} 
\end{ex}

% nguồn: Ok_TuanLuan_Toan11_Lan2
\begin{ex}%[1D3V1-2]%[ma:24c87ad0]
	 Tính các giới hạn sau đây
	\begin{listEX}[2]
		\item $\lim \left(\sqrt{n^2+7}-\sqrt{n^2+5}\right)$. 
		\item $\lim \left(\sqrt{n^2-n+1}-n\right)$. 
		 \item $\lim \left(\sqrt[3]{n^2-n^3}+n\right)$.
		 \item $\lim\limits[\sqrt{n}(\sqrt{n+1}-\sqrt{n})]$.
		\item $\lim \sqrt{\dfrac{8n+2}{2n-1}}$. 
		\item $\lim \sqrt{\dfrac{2n+9}{n+2}}$. 
		\item  $\lim \dfrac{\sqrt{4n^2+1}+2n-1}{\sqrt{n^2+4n+1}+n}$.
		 \item $\lim \dfrac{n^2+\sqrt[3]{1-n^6}}{\sqrt{n^4+1}+n^2}$.
	\end{listEX}
	\loigiai{ 
		\begin{enumerate}
			 \item Ta có $\lim \left(\sqrt{n^2+7}-\sqrt{n^2+5}\right)=\lim \dfrac{n^2+7-n^2-5}{\sqrt{n^2+7}+\sqrt{n^2+5}}=\lim \dfrac{2}{\sqrt{n^2+7}+\sqrt{n^2+5}}=0$.
			\item Ta có $\lim \left(\sqrt{n^2-n+1}-n\right)=\lim  \dfrac{-n+1}{\sqrt{n^2-n+1}+n}=\lim \dfrac{-1+\dfrac{1}{n}}{\sqrt{1-\dfrac{1}{n}+\dfrac{1}{n^2}}+1}=-\dfrac{1}{2}$.
			\item Ta có $\begin{aligned}
			 \lim \left(\sqrt[3]{n^2-n^3}+n\right)&=\lim \dfrac{n^2}{\sqrt[3]{\left(n^2-n^3\right)^2}-n \sqrt[3]{n^2-n^3}+n^2}\\
			&=\lim \dfrac{1}{\sqrt[3]{\left(\dfrac{1}{n}-1\right)^2}-\sqrt[3]{\dfrac{1}{n}-1}+1}\\ 
			 \\
			 &=\dfrac{1}{3}.
		\end{aligned}$
		\item Ta có $\lim \sqrt{n}\left(\sqrt{n+1}-\sqrt{n}\right)=\lim  \dfrac{\sqrt{n}}{\sqrt{n+1}+\sqrt{n}}=\lim \dfrac{1}{\sqrt{1+\dfrac{1}{n}}+1}=\dfrac{1}{2}$.
		\item Ta có $\lim \sqrt{\dfrac{8n+2}{2n-1}}=\lim  \sqrt{\dfrac{n\left(8+\dfrac{2}{n}\right)}{n\left(2-\dfrac{1}{n}\right)}}=\lim \sqrt{\dfrac{8+\dfrac{2}{n}}{2-\dfrac{1}{n}}}=\sqrt{\dfrac{8+0}{2-0}}=2$.
		\item Ta  có $\lim \sqrt{\dfrac{2n+9}{n+2}}=\lim \sqrt{\dfrac{n\left(2+\dfrac{9}{n}\right)}{n\left(1+\dfrac{2}{n}\right)}}=\lim \sqrt{\dfrac{2+\dfrac{9}{n}}{1+\dfrac{2}{n}}}=\sqrt{\dfrac{2}{1}}=\sqrt{2}$.
		\item Ta có $\begin{aligned}
		\lim \dfrac{\sqrt{4n^2+1}+2n-1}{\sqrt{n^2+4n+1}+n}&=\lim  \dfrac{n \sqrt{\left(4+\dfrac{1}{n^2}\right)}+n\left(2-\dfrac{1}{n}\right)}{n \sqrt{\left(1+\dfrac{4}{n}+\dfrac{1}{n^2}\right)}+n}\\
		&=\lim  \dfrac{\sqrt{n^2\left(4+\dfrac{1}{n^2}\right)}+n\left(2-\dfrac{1}{n}\right)}{\sqrt{n^2\left(1+\dfrac{4}{n}+\dfrac{1}{n^2}\right)}+n}\\
		 &=\lim \dfrac{\sqrt{4+\dfrac{1}{n^2}}+2-\dfrac{1}{n}}{\sqrt{1+\dfrac{4}{n}+\dfrac{1}{n^2}}+1}\\
		&=\dfrac{\sqrt{4}+2}{1+1}\\ 
		 &=2.
	\end{aligned}$
	\item Ta có $\begin{aligned}
	\lim  \dfrac{n^2+\sqrt[3]{1-n^6}}{\sqrt{n^4+1}+n^2}&=\lim \dfrac{n^2+\sqrt[3]{n^6\left(\dfrac{1}{n^6}-1\right)}}{\sqrt{n^4\left(1+\dfrac{1}{n^4}\right)}+n^2}\\
	&=\lim  \dfrac{n^2+n^2\sqrt[3]{\dfrac{1}{n^6}-1}}{n^2\sqrt{1+\dfrac{1}{n^4}}+n^2}\\
	&=\lim  \dfrac{n^2\left(1+\sqrt[3]{\dfrac{1}{n^6}-1}\right)}{n^2\left(\sqrt{1+\dfrac{1}{n^4}}+1\right)}\\
	 &=\lim \dfrac{1+\sqrt[3]{\dfrac{1}{n^6}-1}}{\sqrt{1+\dfrac{1}{n^4}}+1}\\
	 &=\dfrac{1+\sqrt[3]{-1}}{1+1}\\
	&=0. 
	\end{aligned}$
	\end{enumerate}
 }
\end{ex}

% nguồn: Ok_TuanLuan_Toan11_Lan2
\begin{ex}%[1D3V1-2]%[ma:510d5ea9]
	Tính các giới hạn  sau đây
	\begin{listEX}[2]
		\item  $A=\lim \dfrac{\sqrt{4n^2+3}-2n+1}{\sqrt{n^2+2n}+3n}$.
		\item $B=\lim \dfrac{2n+1-\sqrt{n^2+2n-4}}{3n+\sqrt{n^2+7}}$. 
		 \item $C=\lim \dfrac{\sqrt{3n^2+1}-\sqrt{n^2-1}}{n}$.
		\item $D=\lim  \dfrac{\sqrt[3]{n^6+n+1}-4\sqrt{n^4+2n-1}}{(2n+3)^2}$.
		\item $E=\lim \dfrac{\sqrt{4n^2-1}+\sqrt[3]{8n^3+2n^2-3}}{\sqrt{16n^2+4n}-\sqrt[4]{n^4+1}}$. 
		\item $F=\lim \dfrac{n^2+\sqrt[3]{1-n^6}}{\sqrt{n^4+1}-n^2}$. 
	\end{listEX}
	 \loigiai{
		\begin{enumerate}
			\item Ta có $\begin{aligned}
			 A&=\lim \dfrac{\sqrt{4n^2+3}-2n+1}{\sqrt{n^2+2n}+3n}\\
			 &=\lim \dfrac{n \sqrt{4+\dfrac{3}{n^2}}-2n+1}{n \sqrt{1+\dfrac{2}{n}}+3n}\\
			 &=\lim \dfrac{n\left(\sqrt{4+\dfrac{3}{n^2}}-2+\dfrac{1}{n}\right)}{n\left(\sqrt{1+\dfrac{2}{n}}+3\right)}\\
			 &=\lim \dfrac{\sqrt{4+\dfrac{3}{n^2}}-2+\dfrac{1}{n}}{\sqrt{1+\dfrac{2}{n}}+3}\\
			 &=0.
		\end{aligned}$
		\item Ta có $\begin{aligned}
		 B&=\lim \dfrac{2n+1-\sqrt{n^2+2n-4}}{3n+\sqrt{n^2+7}}\\
		&=\lim \dfrac{2n+1-n \sqrt{1+\dfrac{2}{n}-\dfrac{4}{n^2}}}{3n+n  \sqrt{1+\dfrac{7}{n^2}}}\\
		&=\lim \dfrac{2+\dfrac{1}{n}-\sqrt{1+\dfrac{2}{n}-\dfrac{4}{n^2}}}{3+\sqrt{1+\dfrac{7}{n}}}\\ 
		&=\dfrac{1}{4}. 
	\end{aligned}$
	\item Ta có $\begin{aligned}
	 C&=\lim \dfrac{\sqrt{3n^2+1}-\sqrt{n^2-1}}{n}\\
	&=\lim \dfrac{n  \sqrt{3+\dfrac{1}{n^2}}-n \sqrt{1-\dfrac{1}{n^2}}}{n}\\
	&=\lim  \dfrac{n\left(\sqrt{3+\dfrac{1}{n^2}}-\sqrt{1-\dfrac{1}{n^2}}\right)}{n}\\
	 &=\lim \left(\sqrt{3+\dfrac{1}{n^2}}-\sqrt{1-\dfrac{1}{n^2}}\right)\\
	 &=\sqrt{3}-1.
	\end{aligned}$
	\item Ta có $\begin{aligned}
	 D&=\lim \dfrac{\sqrt[3]{n^6+n+1}-4\sqrt{n^4+2n-1}}{(2n+3)^2}\\
	&=\lim  \dfrac{\sqrt[3]{1+\dfrac{1}{n^5}+\dfrac{1}{n^6}}-4\sqrt{1+\dfrac{2}{n^3}-\dfrac{1}{n^4}}}{\left(2+\dfrac{3}{n}\right)^2}\\
	 &=\dfrac{1-4}{4}\\
	 &=-\dfrac{3}{4}.
	\end{aligned}$
	\item Ta có $\begin{aligned}
	E&=\lim \dfrac{\sqrt{4n^2-1}+\sqrt[3]{8n^3+2n^2-3}}{\sqrt{16n^2+4n}-\sqrt[4]{n^4+1}}\\ 
	 &=\lim \dfrac{\sqrt{4-\dfrac{1}{n^2}}+\sqrt[3]{8+\dfrac{2}{n}-\dfrac{3}{n^3}}}{\sqrt{16+\dfrac{4}{n}}-\sqrt[4]{1+\dfrac{1}{n^4}}}\\
	 &=\dfrac{2+2}{4-1}\\
	 &=\dfrac{4}{3}.
	\end{aligned}$
	\item Ta có $\begin{aligned}
	F&=\lim \dfrac{n^2+\sqrt[3]{1-n^6}}{\sqrt{n^4+1}-n^2}\\ 
	&=\lim \dfrac{\left(n^6+1-n^6\right)\left(\sqrt{n^4+1}+n^2\right)}{\left(n^4+1-n^4\right)\left[n^4-n^2\cdot \sqrt[3]{1-n^6}+\left(\sqrt[3]{1-n^6}\right)^2\right]}\\ 
	&=\lim  \dfrac{\sqrt{n^4+1}+n^2}{n^4-n^2 \cdot \sqrt[3]{1-n^6}+\left(\sqrt[3]{1-n^6}\right)^2}\\
	 &=\lim \dfrac{n^4\left(\sqrt{\dfrac{1}{n^4}+\dfrac{1}{n^8}}+\dfrac{1}{n^2}\right)}{n^4 \cdot\left[1-\sqrt[3]{\dfrac{1}{n^6}-1}+\left(\sqrt[3]{\dfrac{1}{n^6}-1}\right)^2\right]}\\
	&=\lim  \dfrac{\sqrt{\dfrac{1}{n^4}+\dfrac{1}{n^8}}+\dfrac{1}{n^2}}{1-\sqrt[3]{\dfrac{1}{n^6}-1}+\left(\sqrt[3]{\dfrac{1}{n^6}-1}\right)^2}\\
	&=\lim  \dfrac{\sqrt{0+0}+0}{1-\sqrt[3]{0+1}+(\sqrt[3]{0-1})^2}\\
	 &=\dfrac0{1-1+1}\\
	 &=0.
	\end{aligned}$
	\end{enumerate}
} 
\end{ex}

% nguồn: Ok_TuanLuan_Toan11_Lan2
\begin{ex}%[1D3V1-2]%[ma:b8029af0]
	Tìm $\lim u_n$ biết  $\left(u_n\right)$ có giới hạn hữu hạn và
	 $\left(u_n\right)\colon \heva{&u_1=1\\
	&u_{n+1}=\dfrac{2u_n+3}{u_n+2},\,  n=1,2,3,\ldots
	}$ 
	\loigiai{Bằng phương pháp quy nạp, dễ dàng chứng minh được $u_n>0$ với mọi  $n\in \mathbb{N}^*$.\\
		Đặt $\lim u_n=a \geq 0$.\\ 
		Vì $u_{n+1}=\dfrac{2 u_n+3}{u_n}$ nên $\lim u_{n+1}=\lim \dfrac{2 u_n+3}{u_n+2}\Leftrightarrow a=\dfrac{2  a+3}{a+2} \Leftrightarrow \hoac{
		&a= \sqrt{3} \text{ (nhận)}\\ 
		&a=- \sqrt{3} \text{  (loại).}}$\\
		Vậy  $a=\sqrt{3}$.
	} 
\end{ex}
